% Du gène à la protéine — SVTEEHB Tle TI — Cours signé OnBuch+
% Programme officiel MINESEC. Compiler avec : tectonic N02-du-gene-a-la-proteine.tex
\newcommand{\DOCMATIERE}{SVTEEHB}
\newcommand{\DOCNIVEAU}{Tle TI}
\newcommand{\DOCMODULE}{Module I : Le Monde Vivant}
\newcommand{\DOCLECON}{N02}
\newcommand{\DOCTITRE}{Du gène à la protéine}
% OnBuch+ — préambule « cours signés » ENRICHI (compilé avec tectonic / XeLaTeX)
% Système visuel « Pop » d'OnBuch+ : fond crème (jamais blanc), bordures encre
% épaisses, ombres « dures » décalées, titres Archivo Black en capitales.
% Version MATHÉMATIQUES : graphes (pgfplots/tikz), boîtes pédagogiques
% (définition, propriété, méthode, attention, le savais-tu, exercice résolu,
% exercice, TP, bilan), page de garde et signature OnBuch+ ancrée en bas.
%
% Macros attendues dans main.tex AVANT \input{preamble} :
%   \DOCMATIERE  \DOCNIVEAU  \DOCMODULE  \DOCLECON (numéro)  \DOCTITRE
\documentclass[11pt]{article}

\usepackage{fontspec}
\usepackage[french]{babel}
\usepackage[a4paper,margin=2cm,top=3.4cm,bottom=2.8cm,headheight=1.4cm,footskip=1.3cm]{geometry}
\usepackage{amsmath,amssymb,amsfonts}
\usepackage{mathtools} % \xmapsto, \coloneqq... souvent utilisés par les modèles
\usepackage{cancel} % simplifications barrées (bilans de forces)
% \abs{z} (module, valeur absolue) et \norm{u} (norme), utilisés par les modèles
\providecommand{\abs}{}\let\abs\relax\DeclarePairedDelimiter{\abs}{\lvert}{\rvert}
\providecommand{\norm}{}\let\norm\relax\DeclarePairedDelimiter{\norm}{\lVert}{\rVert}
\usepackage[table]{xcolor} % [table] : fournit aussi \rowcolors (alternance de lignes)
\usepackage{pagecolor}
\usepackage{tikz}
\usetikzlibrary{arrows.meta,positioning,calc,shapes.geometric,shapes.misc,shapes.arrows,decorations.pathmorphing,decorations.pathreplacing,decorations.markings,patterns,shadows,mindmap,trees,fit,backgrounds,matrix,intersections,angles,quotes}
\usepackage{pgfplots}
\usepackage[version=4]{mhchem} % équations nucléaires \ce{^{235}_{92}U}
\usepackage[european,siunitx]{circuitikz} % schémas électriques
\pgfplotsset{compat=1.18}
\usepgfplotslibrary{fillbetween,groupplots} % aires entre courbes (calcul intégral)
\usepackage{enumitem}
\usepackage{fancyhdr}
% [most] charge toutes les bibliothèques tcolorbox (listings, minted, raster,
% documentation...) et gonfle inutilement la mémoire TeX sur un document long
% (~300 boîtes) : on ne charge que ce qui est réellement utilisé.
\usepackage[skins,breakable]{tcolorbox}
\usepackage{graphicx}
\usepackage{colortbl}
\usepackage{booktabs}
\usepackage{tabularx}
\usepackage{multirow}
\usepackage{diagbox} % case d'en-tête diagonale des tableaux à double entrée
% Colonnes extensibles centrée / gauche / droite (C, L, R), souvent utilisées par les modèles
\newcolumntype{C}{>{\centering\arraybackslash}X}
\newcolumntype{L}{>{\raggedright\arraybackslash}X}
\newcolumntype{R}{>{\raggedleft\arraybackslash}X}
\usepackage{array}
\usepackage{multicol}
\usepackage{siunitx}
% parse-numbers=false : accepte aussi \SI{2,5 \times 10^{-3}}{...}
\sisetup{locale=FR,output-decimal-marker={,},group-minimum-digits=4,parse-numbers=false}
\DeclareSIUnit\ohm{\ensuremath{\symup{\Omega}}} % U+2126 absent de la police
\DeclareSIUnit\division{div} % réglages d'oscilloscope (ms/div, V/div)
\DeclareSIUnit\tr{tr} % tours (vitesse de rotation, ex. \SI{1500}{\tr\per\minute})
\DeclareSIUnit\molar{M} % concentration molaire (ex. \SI{3}{\molar}, \SI{1}{\micro\molar})
\DeclareSIUnit\an{an} % année (le modèle confond parfois avec \year, un compteur TeX) — ex. \SI{5}{\kilo\meter\per\an}
\DeclareSIUnit\FCFA{FCFA} % franc CFA (ex. \SI{500}{\FCFA\per\kilo\gram})
\DeclareSIUnit\degreeBrix{\textdegree Brix} % teneur en sucre d'un sirop/jus (réfractométrie)
\DeclareSIUnit\month{mois} % le modèle utilise parfois \month, un compteur TeX, comme unité de durée
\DeclareSIUnit\microliter{\micro\liter} % le modèle écrit parfois \microliter en un seul token
\DeclareSIUnit\million{millions} % dénombrement (ex. \SI{15}{\million\per\mL} spermatozoïdes)
\usepackage{qrcode}
% \degree : commande fréquemment utilisée par les modèles pour un angle ou une
% température en dehors de \SI{}{} (non fournie par les paquets chargés).
\providecommand{\degree}{^{\circ}}
% \qed (fin de démonstration), utilisé par les modèles sans amsthm
\providecommand{\qed}{\hfill$\square$}
% \barre : double barre des valeurs interdites dans un tableau de variations (tkz-tab)
\providecommand{\barre}{\ensuremath{\|}}
% environnement proof (amsthm non chargé) utilisé par les modèles
\ifdefined\proof\else\newenvironment{proof}[1][Démonstration]{\par\noindent\textit{#1.}\ }{\qed\par}\fi
\usepackage{lastpage}
\usepackage{hyperref}
\hypersetup{hidelinks}

% ============================================================
% Palette « Pop » OnBuch+ — mêmes hex que la classe P (plus_ui.dart)
% ============================================================
\definecolor{poporange}{HTML}{F59321}
\definecolor{popdark}{HTML}{E07A0C}
\definecolor{popcream}{HTML}{FFF6E8}
\definecolor{popink}{HTML}{1C1714}
\definecolor{poppurple}{HTML}{7A5AE0}
\definecolor{popgreen}{HTML}{2E9E6A}
\definecolor{poppink}{HTML}{E0457B}
\definecolor{popblue}{HTML}{2F7DE1}
\definecolor{popgold}{HTML}{C98A12}
\definecolor{popmuted}{HTML}{8A7F76}
\definecolor{popline}{HTML}{EADFD2}
\definecolor{popgray}{HTML}{8A7F76}
\colorlet{popgrey}{popgray}
% Alias tolérants (le modèle nomme parfois les couleurs différemment)
\colorlet{popwhite}{white}
\colorlet{popblack}{popink}
\colorlet{popred}{poppink}
\colorlet{popyellow}{popgold}
% Teintes claires (fonds de tableaux/encadrés) — le modèle nomme parfois
% les couleurs avec un suffixe L (ex. popblueL) sans les avoir définies.
\colorlet{poporangeL}{poporange!20}
\colorlet{popdarkL}{popdark!20}
\colorlet{poppurpleL}{poppurple!20}
\colorlet{popgreenL}{popgreen!20}
\colorlet{poppinkL}{poppink!20}
\colorlet{popblueL}{popblue!20}
\colorlet{popgoldL}{popgold!20}
\colorlet{popredL}{popred!20}
\colorlet{popgrayL}{popgray!20}
% Teintes claires pour fonds de tableaux / zones
\colorlet{popblueL}{popblue!10!white}
\colorlet{poporangeL}{poporange!14!white}
\colorlet{popgreenL}{popgreen!12!white}
\colorlet{poppurpleL}{poppurple!12!white}
\colorlet{poppinkL}{poppink!12!white}
\colorlet{popgoldL}{popgold!14!white}

\pagecolor{popcream}

% ============================================================
% Typographie
% ============================================================
\defaultfontfeatures{Ligatures=TeX}
\setmainfont{PlusJakartaSans-Medium.ttf}[Path=../../../fonts/,BoldFont=PlusJakartaSans-Bold.ttf]
% Maths en sans-serif (Fira Math) avec chiffres et lettres droites de Plus
% Jakarta, pour que les formules se fondent dans le texte.
\usepackage[math-style=ISO,bold-style=ISO]{unicode-math}
\setmathfont{FiraMath-Regular.otf}[Path=../../../fonts/]
\setmathfont{PlusJakartaSans-Medium.ttf}[Path=../../../fonts/,range={up/num,up/{Latin,latin},\mathup}]
\usepackage{icomma} % virgule décimale sans espace en mode math
% Caractères mathématiques Unicode absents des polices de texte (Plus Jakarta,
% Archivo Black) : rendus par leur équivalent mathématique, en texte comme en
% formule — sinon ils s'affichent en carré vide (ex. « Arithmétique dans ℤ »).
\usepackage{newunicodechar}
\newunicodechar{ℝ}{\ensuremath{\mathbb{R}}}
\newunicodechar{ℤ}{\ensuremath{\mathbb{Z}}}
\newunicodechar{ℂ}{\ensuremath{\mathbb{C}}}
\newunicodechar{ℕ}{\ensuremath{\mathbb{N}}}
\newunicodechar{ℚ}{\ensuremath{\mathbb{Q}}}
\newunicodechar{𝔻}{\ensuremath{\mathbb{D}}}
\newunicodechar{α}{\ensuremath{\alpha}}
\newunicodechar{β}{\ensuremath{\beta}}
\newunicodechar{γ}{\ensuremath{\gamma}}
\newunicodechar{δ}{\ensuremath{\delta}}
\newunicodechar{ε}{\ensuremath{\varepsilon}}
\newunicodechar{θ}{\ensuremath{\theta}}
\newunicodechar{λ}{\ensuremath{\lambda}}
\newunicodechar{μ}{\ensuremath{\mu}}
\newunicodechar{σ}{\ensuremath{\sigma}}
\newunicodechar{τ}{\ensuremath{\tau}}
\newunicodechar{φ}{\ensuremath{\varphi}}
\newunicodechar{ψ}{\ensuremath{\psi}}
\newunicodechar{ω}{\ensuremath{\omega}}
\newunicodechar{Σ}{\ensuremath{\Sigma}}
% majuscules produites par \MakeUppercase dans les titres (α-aminés -> Α)
\newunicodechar{Α}{\ensuremath{\alpha}}
\newunicodechar{Β}{\ensuremath{\beta}}
\newunicodechar{Γ}{\ensuremath{\Gamma}}
\newunicodechar{Δ}{\ensuremath{\Delta}}
\newunicodechar{Ε}{\ensuremath{\varepsilon}}
\newunicodechar{Θ}{\ensuremath{\Theta}}
\newunicodechar{Λ}{\ensuremath{\Lambda}}
\newunicodechar{Μ}{\ensuremath{\mu}}
\newunicodechar{Τ}{\ensuremath{\tau}}
\newunicodechar{Φ}{\ensuremath{\Phi}}
\newunicodechar{Ψ}{\ensuremath{\Psi}}
\newunicodechar{Ω}{\ensuremath{\Omega}}
\newunicodechar{↦}{\ensuremath{\mapsto}}
\newunicodechar{∈}{\ensuremath{\in}}
\newunicodechar{∉}{\ensuremath{\notin}}
\newunicodechar{∧}{\ensuremath{\wedge}}
\newunicodechar{⇔}{\ensuremath{\Leftrightarrow}}
\newunicodechar{⟺}{\ensuremath{\Longleftrightarrow}}
\newunicodechar{⇒}{\ensuremath{\Rightarrow}}
\newunicodechar{⟹}{\ensuremath{\Longrightarrow}}
\newunicodechar{∘}{\ensuremath{\circ}}
\newunicodechar{∪}{\ensuremath{\cup}}
\newunicodechar{∩}{\ensuremath{\cap}}
\newunicodechar{≡}{\ensuremath{\equiv}}
\newunicodechar{⊂}{\ensuremath{\subset}}
\newunicodechar{⊥}{\ensuremath{\perp}}
\newunicodechar{∃}{\ensuremath{\exists}}
\newunicodechar{∀}{\ensuremath{\forall}}
\newunicodechar{∛}{\ensuremath{\sqrt[3]{\,}}}
\newunicodechar{∜}{\ensuremath{\sqrt[4]{\,}}}
\newunicodechar{⃗}{\ensuremath{{}^{\to}}}
\newunicodechar{ⁿ}{\ifmmode^{n}\else\textsuperscript{\ensuremath{n}}\fi}
\newunicodechar{ˣ}{\ifmmode^{x}\else\textsuperscript{\ensuremath{x}}\fi}
\newunicodechar{ᵇ}{\ifmmode^{b}\else\textsuperscript{\ensuremath{b}}\fi}
\newunicodechar{ʳ}{\ifmmode^{r}\else\textsuperscript{\ensuremath{r}}\fi}
\newunicodechar{ᵘ}{\ifmmode^{u}\else\textsuperscript{\ensuremath{u}}\fi}
\newunicodechar{ʸ}{\ifmmode^{y}\else\textsuperscript{\ensuremath{y}}\fi}
\newunicodechar{ᵃ}{\ifmmode^{a}\else\textsuperscript{\ensuremath{a}}\fi}
\newunicodechar{ᵉ}{\ifmmode^{e}\else\textsuperscript{\ensuremath{e}}\fi}
\newunicodechar{ᵏ}{\ifmmode^{k}\else\textsuperscript{\ensuremath{k}}\fi}
\newunicodechar{ᵖ}{\ifmmode^{p}\else\textsuperscript{\ensuremath{p}}\fi}
\newunicodechar{ᴺ}{\ifmmode^{N}\else\textsuperscript{\ensuremath{N}}\fi}
\newunicodechar{ᶜ}{\ifmmode^{c}\else\textsuperscript{\ensuremath{c}}\fi}
\newunicodechar{ʰ}{\ifmmode^{h}\else\textsuperscript{\ensuremath{h}}\fi}
\newunicodechar{ᵠ}{\ifmmode^{q}\else\textsuperscript{\ensuremath{q}}\fi}
\newunicodechar{ₙ}{\ifmmode_{n}\else\textsubscript{\ensuremath{n}}\fi}
\newunicodechar{ₐ}{\ifmmode_{a}\else\textsubscript{\ensuremath{a}}\fi}
\newunicodechar{ᵢ}{\ifmmode_{i}\else\textsubscript{\ensuremath{i}}\fi}
\newunicodechar{ᵣ}{\ifmmode_{r}\else\textsubscript{\ensuremath{r}}\fi}
\newunicodechar{ₖ}{\ifmmode_{k}\else\textsubscript{\ensuremath{k}}\fi}
\newunicodechar{ᵦ}{\ifmmode_{\beta}\else\textsubscript{\ensuremath{\beta}}\fi}
\newunicodechar{ₓ}{\ifmmode_{x}\else\textsubscript{\ensuremath{x}}\fi}
\newfontfamily\popdisplay{ArchivoBlack-Regular.ttf}[Path=../../../fonts/]
\newfontfamily\poplogo{SpaceGrotesk-Bold.ttf}[Path=../../../fonts/]
\newfontfamily\popsemi{PlusJakartaSans-SemiBold.ttf}[Path=../../../fonts/]
\newfontfamily\popxbold{PlusJakartaSans-ExtraBold.ttf}[Path=../../../fonts/]
\newfontfamily\popxboldls{PlusJakartaSans-ExtraBold.ttf}[Path=../../../fonts/,LetterSpace=3.0]

\setlist[enumerate]{leftmargin=1.7em,itemsep=5pt,topsep=4pt}
\setlist[itemize]{leftmargin=1.7em,itemsep=4pt,topsep=4pt,label={\color{poporange}\textbullet}}
\setlength{\parindent}{0pt}
\setlength{\parskip}{5pt}
\renewcommand{\arraystretch}{1.25}

% pgfplots : style Pop par défaut
\pgfplotsset{
  every axis/.append style={axis line style={popink,line width=1pt},tick style={popink},
    grid style={popline,line width=0.5pt},label style={font=\small},tick label style={font=\footnotesize}},
  popcurve/.style={poporange,line width=1.8pt,no markers,smooth},
}
% Même style disponible dans un \draw TikZ simple (hors pgfplots)
\tikzset{popcurve/.style={poporange,line width=1.8pt}}
\tikzset{popgrid/.style={popline,very thin}}
\colorlet{popcurve}{poporange} % le nom du style est parfois utilisé comme couleur

% ============================================================
% Pill / squiggle
% ============================================================
\newcommand{\poppill}[3]{%
  \tikz[baseline=-0.55ex]\node[draw=popink,line width=1.2pt,rounded corners=2.6pt,%
    fill=#2,inner xsep=6.5pt,inner ysep=3pt]{%
    \popxboldls\fontsize{7}{7}\selectfont\color{#3}\MakeUppercase{#1}};%
}
\newcommand{\popsquiggle}[1]{%
  \tikz[baseline=-0.4ex]{
    \draw[#1,line width=2.6pt,line cap=round]
      plot [smooth] coordinates {(0,0.09) (0.34,0.01) (0.68,0.21) (1.02,0.01) (1.36,0.10)};
  }%
}
% Mot-clé surligné (marqueur orange sous le texte)
\newcommand{\cle}[1]{{\popxbold\color{popdark}#1}}

% ============================================================
% En-tête / pied
% ============================================================
\pagestyle{fancy}
\fancyhf{}
\renewcommand{\headrulewidth}{0pt}
\renewcommand{\footrulewidth}{0pt}
\lhead{\raisebox{-0.15ex}{\poplogo\fontsize{13}{13}\selectfont\color{popink}OnBuch\color{poporange}+}}
\rhead{\poppill{\DOCMATIERE\ \textbullet\ \DOCNIVEAU\ \textbullet\ Leçon \DOCLECON}{white}{popink}}
\lfoot{\popxbold\fontsize{7.6}{7.6}\selectfont\color{popmuted}\MakeUppercase{Cours signé OnBuch+}\ \textbullet\ {\popsemi\DOCTITRE}}
\rfoot{\tikz[baseline=-0.6ex]\node[rounded corners=8.5pt,fill=popink,minimum height=17pt,minimum width=17pt,inner xsep=5pt,inner sep=0]{\popxbold\fontsize{8}{8}\selectfont\color{popcream}\thepage\,/\,\pageref*{LastPage}};}

% ============================================================
% Titres
% ============================================================
\newcommand{\chaptitre}[2]{%
  \begin{tcolorbox}[enhanced,colback=poporange,colframe=popink,boxrule=2.6pt,arc=11pt,
      drop shadow=popink,left=15pt,right=15pt,top=12pt,bottom=11pt,before skip=3pt,after skip=18pt]
    \poppill{#2}{popink}{popcream}\par\vspace{9pt}
    {\raggedright\hyphenpenalty=10000\exhyphenpenalty=10000\popdisplay\fontsize{22}{24}\selectfont\color{popcream}\MakeUppercase{#1}\par}
  \end{tcolorbox}
}
% Section numérotée : pastille orange avec le numéro + titre Archivo
\newcounter{popsec}
\newcommand{\coursec}[1]{%
  \stepcounter{popsec}\setcounter{popsub}{0}%
  \par\vspace{16pt}\noindent
  \begin{minipage}{\linewidth}
  \tikz[baseline=-0.7ex]\node[draw=popink,line width=1.6pt,fill=poporange,rounded corners=4pt,minimum size=22pt,inner sep=0]{\popdisplay\fontsize{12}{12}\selectfont\color{popcream}\thepopsec};%
  \hspace{8pt}{\popdisplay\fontsize{15}{17}\selectfont\color{popink}\MakeUppercase{#1}}\par
  \vspace{4pt}\noindent{\color{poporange}\rule{36pt}{3.6pt}}
  \end{minipage}\par\nopagebreak\vspace{8pt}
}
\newcounter{popsub}
\newcommand{\courssub}[1]{%
  \stepcounter{popsub}%
  \par\vspace{9pt}\noindent{\popxbold\fontsize{11.5}{13}\selectfont\color{popink}{\color{poporange}\thepopsec.\thepopsub}\hspace{6pt}#1}\par\nopagebreak\vspace{4pt}
}

% ============================================================
% Boîtes pédagogiques — même recette : fond blanc, bordure épaisse colorée,
% ombre dure encre, badge plein. Titre optionnel [#1].
% ============================================================
\tcbset{popbase/.style={breakable,enhanced,colback=white,boxrule=2.3pt,arc=9pt,
  shadow={3pt}{-3pt}{0pt}{popink},left=14pt,right=14pt,top=11pt,bottom=11pt,before skip=11pt,after skip=15pt,
  fontupper=\popsemi\fontsize{10.6}{14.5}\selectfont\color{popink},
  fontlower=\popsemi\fontsize{10.6}{14.5}\selectfont\color{popink}}}
% Coche dessinée (le glyphe ✓ n'existe pas dans les polices du document)
\newcommand{\popcheck}[1]{\tikz[baseline=-0.5ex]\draw[#1,line width=1.6pt,line cap=round,line join=round](-0.13,0.0)--(-0.04,-0.09)--(0.13,0.1);}
\providecommand{\coche}{\popcheck{popgreen}}
\providecommand{\resultat}[1]{\par\smallskip\noindent\fcolorbox{popgreen}{popgreenL}{\parbox{\dimexpr\linewidth-2\fboxsep-2\fboxrule\relax}{\textbf{Résultat.} #1}}\par\smallskip}
\newcommand{\popboxhead}[3]{\poppill{#1}{#2}{white}\ifx\relax#3\relax\else\hspace{6pt}{\popxbold\fontsize{10.6}{12}\selectfont\color{popink}#3}\fi\par\vspace{7pt}}

\newtcolorbox{aretenir}[1][]{popbase,colframe=popgreen,before upper={\popboxhead{À retenir}{popgreen}{#1}}}
\newtcolorbox{exemplebox}[1][]{popbase,colframe=poporange,before upper={\popboxhead{Exemple}{poporange}{#1}}}
\newtcolorbox{definition}[1][]{popbase,colframe=popblue,before upper={\popboxhead{Définition}{popblue}{#1}}}
\newtcolorbox{propriete}[1][]{popbase,colframe=poppurple,before upper={\popboxhead{Propriété}{poppurple}{#1}}}
\newtcolorbox{methode}[1][]{popbase,colframe=popgold,before upper={\popboxhead{Méthode}{popgold}{#1}}}
\newtcolorbox{attention}[1][]{popbase,colframe=poppink,before upper={\popboxhead{Attention}{poppink}{#1}}}
\newtcolorbox{experience}[1][]{popbase,colframe=popblue,colback=popblueL,before upper={\popboxhead{Expérience}{popblue}{#1}}}
% « Le savais-tu ? » : carte violette pleine (contexte, culture, Cameroun)
\newtcolorbox{savaistu}[1][]{popbase,colback=poppurple,colframe=popink,
  fontupper=\popsemi\fontsize{10.4}{14.2}\selectfont\color{white},
  before upper={\poppill{Le savais-tu\,?}{white}{poppurple}\ifx\relax#1\relax\else\hspace{6pt}{\popxbold\color{white}#1}\fi\par\vspace{7pt}}}
% Exercice résolu : énoncé en haut, \tcblower puis la solution sur fond crème
\newcounter{popexo}\newcounter{popexr}
\newtcolorbox{exoresolu}[1][]{popbase,colframe=poporange,colbacklower=popcream,
  segmentation style={popink,line width=1.2pt,dashed},
  before upper={\stepcounter{popexr}\popboxhead{Exercice résolu \thepopexr}{poporange}{#1}},
  before lower={\poppill{Solution}{popink}{popcream}\par\vspace{6pt}}}
% Exercice (non résolu en place) — la correction est dans la partie Corrigés
\newtcolorbox{exercice}[1][]{popbase,colframe=popink,
  before upper={\stepcounter{popexo}\popboxhead{Exercice \thepopexo}{popink}{#1}}}
% Corrigé
\newtcolorbox{corrige}[1][]{popbase,colframe=popgreen,colback=popgreenL,
  before upper={\popboxhead{Corrigé}{popgreen}{#1}}}
% Objectifs / prérequis / bilan
\newtcolorbox{objectifs}{popbase,colframe=popink,colback=poporangeL,
  before upper={\popboxhead{Objectifs de la leçon}{popink}{}}}
\newtcolorbox{prerequis}{popbase,colframe=popink,colback=popblueL,
  before upper={\popboxhead{Prérequis}{popblue}{}}}
\newtcolorbox{bilan}{popbase,colframe=popink,colback=white,boxrule=2.8pt,
  before upper={\popboxhead{Fiche bilan}{poporange}{L'essentiel à connaître pour l'examen}}}
% Figure encadrée avec légende
\newtcolorbox{popfigure}[1][]{enhanced,breakable=false,colback=white,colframe=popline,boxrule=1.4pt,arc=8pt,
  left=10pt,right=10pt,top=10pt,bottom=8pt,before skip=10pt,after skip=12pt,halign=center,
  lower separated=false,halign lower=center,
  fontlower=\popsemi\fontsize{9}{11.5}\selectfont\color{popmuted},
  #1}
\newcommand{\legende}[1]{\par\vspace{6pt}{\popsemi\fontsize{9}{11.5}\selectfont\color{popmuted}#1}}

% ============================================================
% Signature OnBuch+
% ============================================================
\newcommand{\poplogoinline}[1]{{\poplogo\fontsize{#1}{#1}\selectfont\color{popink}OnBuch\color{poporange}+}}
\newcommand{\signatureonbuch}{%
  \begin{tcolorbox}[enhanced,colback=popink,colframe=popink,boxrule=2.2pt,arc=10pt,
      drop shadow=poporange,left=14pt,right=14pt,top=10pt,bottom=10pt,before skip=0pt,after skip=0pt]
    \tikz[baseline=-0.7ex]\node[circle,fill=popgreen,draw=popcream,line width=1.3pt,minimum size=22pt,inner sep=0]{\popcheck{white}};%
    \hspace{9pt}\begin{minipage}[c]{0.62\linewidth}
      {\popxbold\fontsize{12}{13}\selectfont\color{popcream}Signé\ \textbullet\ {\poplogo OnBuch\color{poporange}+}}\par\vspace{2pt}
      {\raggedright\popsemi\fontsize{8}{10.5}\selectfont\color{popline}\MakeUppercase{Équipe pédagogique OnBuch+}\\\MakeUppercase{Conforme au programme officiel MINESEC}\par}
    \end{minipage}\hfill
    {\poplogo\fontsize{22}{22}\selectfont\color{popcream}OnBuch\color{poporange}+}
  \end{tcolorbox}%
}

% ============================================================
% Page de garde
% #1 titre · #2 matière · #3 niveau · #4 module · #5 n° leçon · #6 durée ·
% #7 description / objectifs (texte ou itemize)
% ============================================================
\newcommand{\pagegarde}[7]{%
  \thispagestyle{empty}
  \newgeometry{a4paper,margin=1.7cm,top=1.6cm,bottom=1.4cm}%
  \begin{tikzpicture}[remember picture,overlay]
    \fill[poporange!18] ([xshift=-3.2cm,yshift=-3.6cm]current page.north east) circle (4.6cm);
    \fill[poppurple!14] ([xshift=2.4cm,yshift=7.6cm]current page.south west) circle (3.8cm);
  \end{tikzpicture}%
  {\poplogo\fontsize{30}{30}\selectfont\color{popink}OnBuch\color{poporange}+}\hfill
  \raisebox{0.6ex}{\poppill{Cours signé \textbullet\ Premium}{popink}{popcream}}\par
  \vspace{1pt}\popsquiggle{poppurple}\par
  \vspace{8pt}
  \begin{tcolorbox}[enhanced,colback=poporange,colframe=popink,boxrule=2.8pt,arc=13pt,
      drop shadow=popink,left=18pt,right=18pt,top=12pt,bottom=13pt,after skip=10pt]
    \poppill{#2 \textbullet\ #3}{popink}{popcream}\hspace{5pt}\poppill{#4}{white}{popink}\par\vspace{8pt}
    {\popdisplay\fontsize{14}{14}\selectfont\color{popink}LEÇON #5}\par\vspace{4pt}
    {\raggedright\hyphenpenalty=10000\exhyphenpenalty=10000\popdisplay\fontsize{25}{27}\selectfont\color{popcream}\MakeUppercase{#1}\par}
    \vspace{8pt}
    {\popxbold\fontsize{9.5}{9.5}\selectfont\color{popink}\MakeUppercase{Durée conseillée : #6}}
  \end{tcolorbox}
  \begin{tcolorbox}[enhanced,colback=white,colframe=popink,boxrule=2.2pt,arc=10pt,
      drop shadow=popink,left=16pt,right=16pt,top=10pt,bottom=10pt,after skip=8pt]
    \poppill{Dans cette leçon}{poppurple}{white}\par\vspace{5pt}
    \popsemi\fontsize{9.3}{12.6}\selectfont\color{popink}#7
  \end{tcolorbox}
  \noindent\begin{minipage}[t]{0.487\linewidth}
    \begin{tcolorbox}[enhanced,colback=popgreen,colframe=popink,boxrule=2.2pt,arc=10pt,
        drop shadow=popink,left=11pt,right=11pt,top=10pt,bottom=10pt,height=3.1cm,valign=center]
      \tcbox[colback=white,colframe=popink,boxrule=1.3pt,arc=5pt,boxsep=0pt,nobeforeafter,
        left=3pt,right=3pt,top=3pt,bottom=3pt,valign=center]{\qrcode[height=1.9cm]{https://play.google.com/store/apps/details?id=com.luvvix.onbuch&hl=fr}}%
      \hspace{8pt}\begin{minipage}[c]{0.5\linewidth}
        {\popdisplay\fontsize{11}{12.5}\selectfont\color{white}\MakeUppercase{Télécharge OnBuch}}\par\vspace{4pt}
        {\raggedright\popsemi\fontsize{8.4}{11}\selectfont\color{white}Gratuit \textbullet\ Google Play\\Tuteur IA, annales, concours, résultats\par}
      \end{minipage}
    \end{tcolorbox}
  \end{minipage}\hfill
  \begin{minipage}[t]{0.487\linewidth}
    \begin{tcolorbox}[enhanced,colback=poppurple,colframe=popink,boxrule=2.2pt,arc=10pt,
        drop shadow=popink,left=11pt,right=11pt,top=10pt,bottom=10pt,height=3.1cm,valign=center]
      \tikz[baseline=-0.7ex]\node[circle,fill=white,draw=popink,line width=1.3pt,minimum size=1.6cm,inner sep=0]{\popdisplay\fontsize{24}{24}\selectfont\color{poppurple}+};%
      \hspace{8pt}\begin{minipage}[c]{0.6\linewidth}
        {\popdisplay\fontsize{11}{12.5}\selectfont\color{white}\MakeUppercase{Passe à OnBuch+}}\par\vspace{4pt}
        {\raggedright\popsemi\fontsize{8.4}{11}\selectfont\color{white}Tous les cours signés, Tuteur IA illimité, fiches PDF — onglet OnBuch+ de l'app\par}
      \end{minipage}
    \end{tcolorbox}
  \end{minipage}
  \vspace{8pt}
  \popcontenu
  \vfill
  \signatureonbuch
  \restoregeometry
  \newpage
}

% Rangée « Ce cours contient » (page de garde)
\newcommand{\poptile}[3]{% #1 couleur #2 gros texte #3 légende
  \begin{tcolorbox}[enhanced,colback=#1,colframe=popink,boxrule=2pt,arc=9pt,shadow={2.5pt}{-2.5pt}{0pt}{popink},
      left=6pt,right=6pt,top=9pt,bottom=9pt,halign=center,valign=center,height=1.75cm,width=\linewidth,nobeforeafter]
    {\popdisplay\fontsize{12.5}{13}\selectfont\color{white}\MakeUppercase{#2}}\par\vspace{3pt}
    {\centering\popsemi\fontsize{7.8}{9.5}\selectfont\color{white}#3\par}
  \end{tcolorbox}}
\newcommand{\popcontenu}{%
  \par\noindent{\popxboldls\fontsize{7.5}{7.5}\selectfont\color{popmuted}CE COURS SIGNÉ CONTIENT}\par\vspace{7pt}
  \noindent\begin{minipage}[t]{0.235\linewidth}\poptile{popblue}{Cours}{complet, illustré, pas à pas}\end{minipage}\hfill
  \begin{minipage}[t]{0.235\linewidth}\poptile{poporange}{Méthodes}{et exercices résolus}\end{minipage}\hfill
  \begin{minipage}[t]{0.235\linewidth}\poptile{poppink}{Exercices}{type examen + corrigés}\end{minipage}\hfill
  \begin{minipage}[t]{0.235\linewidth}\poptile{popgreen}{Fiche bilan}{l'essentiel à retenir}\end{minipage}\par}

% Sommaire
\newtcolorbox{sommaire}{enhanced,colback=white,colframe=popink,boxrule=2.2pt,arc=10pt,
  drop shadow=popink,left=18pt,right=18pt,top=13pt,bottom=13pt,before skip=6pt,after skip=20pt,
  before upper={\poppill{Sommaire}{poppurple}{white}\par\vspace{9pt}\popsemi\fontsize{10.6}{14.5}\selectfont\color{popink}}}

% Signature de fin de cours — poussée tout en bas de la dernière page
\newcommand{\findecours}{%
  \par\vfill\nopagebreak
  \begin{center}\popsquiggle{poporange}\end{center}\vspace{4pt}
  \signatureonbuch
}
\AtBeginDocument{\def\llbracket{\mathopen{[\mkern-3.5mu[}}\def\rrbracket{\mathclose{]\mkern-3.5mu]}}} % intervalles d'entiers (glyphe ⟦ absent de la police)
% Glyphes absents de Fira Math : redéfinis à partir de glyphes disponibles
\AtBeginDocument{%
  \def\setminus{\mathbin{\backslash}}\let\smallsetminus\setminus
  \def\vdots{\mathord{\vbox{\baselineskip3.2pt\lineskiplimit0pt\kern4pt\hbox{.}\hbox{.}\hbox{.}}}}%
  \def\ddots{\mathinner{\mkern1mu\raise6.5pt\hbox{.}\mkern2mu\raise3.6pt\hbox{.}\mkern2mu\raise0.7pt\hbox{.}\mkern1mu}}%
  \def\triangle{\mathord{\scalebox{0.9}{$\symup{\Delta}$}}}%
  \def\nabla{\mathord{\raisebox{\depth}{\scalebox{1}[-1]{$\symup{\Delta}$}}}}%
  \def\checkmark{\popcheck{popdark}}%
  \def\star{\mathbin{\ast}}\def\bigstar{\mathord{\ast}}%
  \def\diamond{\mathbin{\scalebox{0.6}{$\Diamond$}}}%
  \def\bigcup{\mathop{\mathchoice{\raisebox{-0.3em}{\scalebox{1.7}{$\cup$}}}{\scalebox{1.25}{$\cup$}}{\cup}{\cup}}\displaylimits}%
  \def\bigcap{\mathop{\mathchoice{\raisebox{-0.3em}{\scalebox{1.7}{$\cap$}}}{\scalebox{1.25}{$\cap$}}{\cap}{\cap}}\displaylimits}%
  \def\lesseqqgtr{\mathrel{\lessgtr}}\def\bigoplus{\mathop{\oplus}\displaylimits}%
}
\providecommand{\ssi}{si et seulement si} % abréviation fréquente
\AtBeginDocument{\providecommand{\wideparen}{\overparen}} % arc de cercle

%% FIGFIX-BEGIN v5 — correctifs globaux des figures (docs/onbuch-generation/tools/figfix)
\usepackage{adjustbox}\usepackage{etoolbox}\tcbuselibrary{raster}
% toute figure TikZ/pgfplots plus large que la ligne est réduite à la largeur disponible
\BeforeBeginEnvironment{tikzpicture}{\begin{adjustbox}{max width=\linewidth}}
\AfterEndEnvironment{tikzpicture}{\end{adjustbox}}
% légendes pgfplots lisibles par-dessus les courbes
\pgfplotsset{every axis/.append style={legend style={fill=white,fill opacity=0.92,text opacity=1,draw=black!15,inner sep=3pt}}}
% étiquettes d'arêtes (syntaxe edge["..."]) avec fond blanc
\tikzset{every edge quotes/.append style={fill=white,inner sep=1.2pt}}
%% FIGFIX-END
\begin{document}

\pagegarde{Du gène à la protéine}{SVTEEHB}{Tle TI}{Module I : Le Monde Vivant}{N02}{4 h}{Cette leçon explore les mécanismes moléculaires qui assurent la transmission et l'expression de l'information génétique : structure des acides nucléiques, réplication semi-conservative de l'ADN, puis biosynthèse des protéines (transcription et traduction). Elle établit enfin le lien entre gène, protéine et caractère à travers le code génétique, en s'appuyant sur les expériences historiques qui ont fondé la génétique moléculaire.
\vspace{4pt}
\begin{multicols}{2}\small
\begin{itemize}
\item À la fin de cette leçon, je sais décrire l'organisation architecturale des molécules d'ADN et d'ARN (nucléotides, liaisons, double hélice)
\item À la fin de cette leçon, je sais comparer l'ADN et les différents types d'ARN (structure, localisation, rôle)
\item À la fin de cette leçon, je sais réaliser une maquette de molécule d'ADN respectant la complémentarité des bases
\item À la fin de cette leçon, je sais expliquer le mécanisme de la réplication semi-conservative à partir des expériences de Meselson et Stahl
\item À la fin de cette leçon, je sais identifier et localiser les étapes et les acteurs de la transcription et de la traduction
\item À la fin de cette leçon, je sais utiliser le tableau du code génétique pour traduire une séquence d'ARNm en chaîne polypeptidique
\end{itemize}\end{multicols}}

\begin{sommaire}
\begin{multicols}{2}
\begin{enumerate}
\item I. Structure architecturale des acides nucléiques
\item II. La duplication de l'ADN : la réplication semi-conservative
\item III. La transcription : de l'ADN à l'ARN messager
\item IV. La traduction : de l'ARNm à la protéine
\item V. Le code génétique et la relation gène – protéine – caractère
\item VI. Variation de la quantité d'ADN au cours de la méiose
\item Activité d'intégration
\item Méthodes et exercices résolus
\item Exercices
\item Corrigés des exercices
\item Fiche bilan
\end{enumerate}
\end{multicols}
\end{sommaire}

\begin{objectifs}
\begin{itemize}
\item À la fin de cette leçon, je sais décrire l'organisation architecturale des molécules d'ADN et d'ARN (nucléotides, liaisons, double hélice)
\item À la fin de cette leçon, je sais comparer l'ADN et les différents types d'ARN (structure, localisation, rôle)
\item À la fin de cette leçon, je sais réaliser une maquette de molécule d'ADN respectant la complémentarité des bases
\item À la fin de cette leçon, je sais expliquer le mécanisme de la réplication semi-conservative à partir des expériences de Meselson et Stahl
\item À la fin de cette leçon, je sais identifier et localiser les étapes et les acteurs de la transcription et de la traduction
\item À la fin de cette leçon, je sais utiliser le tableau du code génétique pour traduire une séquence d'ARNm en chaîne polypeptidique
\item À la fin de cette leçon, je sais établir la relation entre une mutation ponctuelle, la protéine modifiée et le caractère altéré (exemple de la drépanocytose)
\item À la fin de cette leçon, je sais tracer et interpréter la courbe de variation de la quantité d'ADN au cours de la méiose
\end{itemize}
\end{objectifs}

\begin{prerequis}
\begin{itemize}
\item La cellule eucaryote : noyau, cytoplasme, ribosomes, réticulum (classe de Première)
\item Les chromosomes et les divisions cellulaires : mitose et méiose (Première / début de Terminale)
\item Les protéines : chaînes polypeptidiques d'acides aminés, rôles enzymatique et structural (Seconde/Première)
\item Notion de gène et de caractère héréditaire (Première, brassage génétique)
\item Grandeurs et mesures : lecture et exploitation de courbes expérimentales
\end{itemize}
\end{prerequis}

\newpage

\coursec{I. Structure architecturale des acides nucléiques}

Au marché Mokolo de Yaoundé, deux jumeaux se ressemblent trait pour trait, tandis que leurs cousins présentent des traits très différents. À l'hôpital central, on diagnostique parfois la drépanocytose, une maladie où un seul « mot » de l'héritage génétique modifié transforme l'hémoglobine et déforme les globules rouges. Comment une information inscrite dans le noyau de nos cellules peut-elle commander la fabrication d'une protéine précise et déterminer ainsi nos caractères ? Comment cette information est-elle fidèlement recopiée avant chaque division cellulaire ? C'est le voyage du gène à la protéine que nous allons entreprendre.

\courssub{Mise en évidence expérimentale : l'ADN, support de l'information génétique}

Avant d'étudier la structure de l'ADN, il est fondamental de comprendre comment les biologistes ont identifié cette molécule comme le support de l'hérédité. Deux expériences historiques, menées avec rigueur, ont levé le doute qui persistait sur le rôle respectif des protéines et des acides nucléiques.

\begin{experience}[L'expérience de transformation de Griffith (1928) — Le « principe transformant »]
\textbf{Observation :} Frederick Griffith travaille sur \emph{Streptococcus pneumoniae} (pneumocoque), responsable de pneumonies. Il existe deux souches :
\begin{itemize}
    \item La souche \textbf{S} (Smooth) : capsule polysaccharidique lisse, \textbf{virulente} (tue la souris).
    \item La souche \textbf{R} (Rough) : sans capsule, colonie rugueuse, \textbf{non virulente}.
\end{itemize}
\textbf{Expérience :} Griffith injecte à des souris :
\begin{enumerate}
    \item Souche S vivante $\rightarrow$ souris meurt.
    \item Souche R vivante $\rightarrow$ souris survit.
    \item Souche S tuée par la chaleur $\rightarrow$ souris survit.
    \item \textbf{Souche S tuée par la chaleur + Souche R vivante} $\rightarrow$ \textbf{souris meurt}. On retrouve des bactéries S vivantes dans son sang.
\end{enumerate}
\textbf{Interprétation :} Un « principe transformant » chimique, résistant à la chaleur, passe des bactéries S mortes aux bactéries R vivantes et leur confère la virulence (et la capsule). Ce principe est l'information génétique.
\textbf{Conclusion :} L'information héréditaire peut être transférée entre cellules par une molécule chimique stable.
\end{experience}

\begin{experience}[L'identification du principe transformant : Avery, MacLeod et McCarty (1944)]
\textbf{Problématique :} Quelle est la nature chimique de ce principe ? Protéine, ADN ou ARN ?
\textbf{Protocole :} Ils préparent un extrait purifié de souche S (tuée par la chaleur). Ils traitent cet extrait par des enzymes spécifiques :
\begin{itemize}
    \item Protéases (détruisent les protéines) $\rightarrow$ transformation \textbf{possible}.
    \item RNases (détruisent l'ARN) $\rightarrow$ transformation \textbf{possible}.
    \item \textbf{DNases (détruisent l'ADN)} $\rightarrow$ transformation \textbf{impossible}.
\end{itemize}
\textbf{Conclusion :} Seule la destruction de l'ADN abolit le pouvoir transformant. \textbf{L'ADN est le support de l'information génétique.}
\end{experience}

\begin{experience}[L'expérience de Hershey et Chase (1952) — Confirmation sur bactériophage]
\textbf{Contexte :} Le bactériophage T2 (virus infectant \emph{E. coli}) est constitué d'une capside protéique et d'ADN interne. Seule l'ADN pénètre dans la bactérie pour diriger la synthèse de nouveaux phages.
\textbf{Stratégie :} Marquage radioactif différentiel des deux composants.
\begin{itemize}
    \item \textbf{Marquage au \ce{^{32}P} (Phosphore 32) :} Incorporé dans l'\textbf{ADN} (phosphate du squelette), absent des protéines.
    \item \textbf{Marquage au \ce{^{35}S} (Soufre 35) :} Incorporé dans les \textbf{protéines} (acides aminés cystéine, méthionine), absent de l'ADN.
\end{itemize}
\textbf{Protocole :}
\begin{enumerate}
    \item Infection de \emph{E. coli} par phages \ce{^{32}P} (ADN marqué) $\rightarrow$ agitation violente (mélangeur Waring Blendor) pour détacher les capsides vides $\rightarrow$ centrifugation.
    \item Infection de \emph{E. coli} par phages \ce{^{35}S} (protéines marquées) $\rightarrow$ même traitement.
\end{enumerate}
\textbf{Résultats :}
\begin{center}
\begin{tabularx}{\linewidth}{|l|X|X|}\hline
\textbf{Marquage} & \textbf{Radioactivité dans le culot (bactéries)} & \textbf{Radioactivité dans le surnageant (capsides)} \\ \hline
\ce{^{32}P} (ADN) & \textbf{Forte} (pénètre dans la bactérie) & Faible \\ \hline
\ce{^{35}S} (Protéines) & Faible & \textbf{Forte} (reste à l'extérieur) \\ \hline
\end{tabularx}
\end{center}
\textbf{Interprétation :} Seul l'ADN marqué pénètre dans la cellule hôte et dirige la production de nouveaux phages.
\textbf{Conclusion définitive :} \textbf{L'ADN est la molécule support de l'information génétique} (chez les virus et, par extension, chez tous les êtres vivants).
\end{experience}

\begin{popfigure}
\begin{tikzpicture}[scale=0.9, every node/.style={font=\small}]
    % Phage 32P
    \begin{scope}[xshift=-7cm]
        \draw[popblue, thick] (0,0) ellipse (0.8 and 0.4);
        \draw[popblue, thick] (0,0.4) -- (0,1.5);
        \draw[popblue, thick] (-0.8,0.4) -- (-1.2,1.5);
        \draw[popblue, thick] (0.8,0.4) -- (1.2,1.5);
        \node[popblue] at (0, -0.8) {\textbf{Phage \ce{^{32}P}}};
        \node[popblue] at (0, 0) {ADN*};
        \node[popblue] at (0, 1.8) {Capside};
        \draw[->, thick, popdark] (1.5, 0.7) -- (2.5, 0.7) node[midway, above] {Infection};
    \end{scope}
    
    % Bacterie
    \begin{scope}[xshift=0cm]
        \draw[popgreen, thick, rounded corners=5pt] (-1.5,-1) rectangle (1.5, 1);
        \node[popgreen] at (0, -1.5) {\textbf{Bactérie \emph{E. coli}}};
        \draw[popblue, thick, fill=popblue!20] (-0.5,0) rectangle (0.5,0.5) node[midway] {ADN*};
        \draw[->, thick, popdark] (2.5, 0) -- (3.5, 0) node[midway, above] {Agitation};
    \end{scope}
    
    % Centrifugation 32P
    \begin{scope}[xshift=7cm]
        \draw[thick] (0,-2) -- (0,2) -- (3,2) -- (3,-2) -- cycle; % Tube
        \draw[fill=popgreen!30] (0,-2) rectangle (3,-0.5); % Culot
        \node at (1.5, -1.2) {\textbf{Culot : Bactéries + ADN*}};
        \draw[fill=popblue!10] (0,0.5) rectangle (3,2); % Surnageant
        \node at (1.5, 1.2) {Surnageant : Capsides};
        \node at (1.5, 2.5) {\textbf{\ce{^{32}P} : Radioactivité dans le CULOT}};
    \end{scope}
    
    % Phage 35S (bas)
    \begin{scope}[xshift=-7cm, yshift=-5cm]
        \draw[poporange, thick] (0,0) ellipse (0.8 and 0.4);
        \draw[poporange, thick] (0,0.4) -- (0,1.5);
        \draw[poporange, thick] (-0.8,0.4) -- (-1.2,1.5);
        \draw[poporange, thick] (0.8,0.4) -- (1.2,1.5);
        \node[poporange] at (0, -0.8) {\textbf{Phage \ce{^{35}S}}};
        \node[poporange] at (0, 1.8) {Capside*};
        \node[poporange] at (0, 0) {ADN};
        \draw[->, thick, popdark] (1.5, 0.7) -- (2.5, 0.7) node[midway, above] {Infection};
    \end{scope}
    
    \begin{scope}[xshift=0cm, yshift=-5cm]
        \draw[popgreen, thick, rounded corners=5pt] (-1.5,-1) rectangle (1.5, 1);
        \node[popgreen] at (0, -1.5) {\textbf{Bactérie \emph{E. coli}}};
        \draw[poporange, thick, fill=poporange!20] (-0.5,0) rectangle (0.5,0.5) node[midway] {Capside*};
        \draw[->, thick, popdark] (2.5, 0) -- (3.5, 0) node[midway, above] {Agitation};
    \end{scope}
    
    \begin{scope}[xshift=7cm, yshift=-5cm]
        \draw[thick] (0,-2) -- (0,2) -- (3,2) -- (3,-2) -- cycle;
        \draw[fill=popgreen!30] (0,-2) rectangle (3,-0.5);
        \node at (1.5, -1.2) {Culot : Bactéries};
        \draw[fill=poporange!10] (0,0.5) rectangle (3,2);
        \node at (1.5, 1.2) {\textbf{Surnageant : Capsides*}};
        \node at (1.5, 2.5) {\textbf{\ce{^{35}S} : Radioactivité dans le SURNAGEANT}};
    \end{scope}
\end{tikzpicture}
\legende{Schéma simplifié de l'expérience de Hershey et Chase (1952). Haut : marquage de l'ADN au \ce{^{32}P} $\rightarrow$ radioactivité retrouvée dans le culot (bactéries). Bas : marquage des protéines au \ce{^{35}S} $\rightarrow$ radioactivité retrouvée dans le surnageant (capsides vides). Seule l'ADN pénètre dans la cellule.}
\end{popfigure}

\courssub{Le nucléotide : unité structurale commune}

Les acides nucléiques (ADN et ARN) sont des polymères linéaires. Leur monomère est le \textbf{nucléotide}.

\begin{definition}[Nucléotide]
Un \cle{nucléotide} est l'unité monomère des acides nucléiques. Il est constitué de trois éléments covalemment liés :
\begin{enumerate}
    \item Un \textbf{groupe phosphate} (\ce{PO4^{3-}}), acide, porteur des charges négatives.
    \item Un \textbf{sucre (pentose)} à 5 atomes de carbone.
    \item Une \textbf{base azotée} (ou base nucléique), cycle hétérocyclique basique.
\end{enumerate}
La liaison entre le phosphate (carbone 5' du sucre) et le carbone 3' du sucre suivant forme le \textbf{squelette phosphodiester} de la chaîne.
\end{definition}

\begin{popfigure}
\begin{tikzpicture}[scale=1.1, every node/.style={font=\small}]
    % Nucleotide unique
    \begin{scope}[xshift=-5cm]
        % Phosphate
        \draw[popblue, thick, fill=popblue!10] (-1.2,0.5) -- (-0.2,1) -- (0.8,0.5) -- (0.8,-0.5) -- (-0.2,-1) -- (-1.2,-0.5) -- cycle;
        \node[popblue] at (-0.2, 0) {\textbf{P}};
        \node[popblue] at (-0.2, -1.5) {Phosphate};
        
        % Sucre (Pentose) - Représentation simplifiée cycle
        \draw[popgreen, thick, fill=popgreen!10] (1.2,1) -- (2.2,1.5) -- (3.2,1) -- (3.2,-1) -- (2.2,-1.5) -- (1.2,-1) -- cycle;
        \node[popgreen] at (2.2, 0) {\textbf{S}};
        \node[popgreen] at (2.2, -2) {Sucre (Pentose)};
        % Carbons numbers
        \node[popgreen, font=\tiny] at (1.2, 1.2) {1'};
        \node[popgreen, font=\tiny] at (3.2, 1.2) {2'};
        \node[popgreen, font=\tiny] at (3.5, 0) {3'};
        \node[popgreen, font=\tiny] at (3.2, -1.2) {4'};
        \node[popgreen, font=\tiny] at (1.2, -1.2) {5'};
        
        % Base
        \draw[poporange, thick, fill=poporange!10] (4.2,1) -- (5.2,1.5) -- (6.2,1) -- (6.7,0) -- (6.2,-1) -- (5.2,-1.5) -- (4.2,-1) -- cycle;
        \node[poporange] at (5.2, 0) {\textbf{B}};
        \node[poporange] at (5.2, -2) {Base azotée};
        
        % Liaisons
        \draw[thick, popdark] (0.8,0) -- (1.2,0) node[midway, above, font=\tiny] {C5'-O-P};
        \draw[thick, popdark] (3.2,0) -- (4.2,0) node[midway, above, font=\tiny] {C1'-N};
        
        \node at (2.2, 2.5) {\textbf{Nucléotide}};
    \end{scope}
    
    % Difference Ribose / Desoxyribose
    \begin{scope}[xshift=6cm]
        \node at (0, 2.5) {\textbf{Différence des sucres}};
        % Ribose (ARN)
        \draw[popgreen, thick, fill=popgreen!10] (-2,1) -- (-1,1.5) -- (0,1) -- (0,-1) -- (-1,-1.5) -- (-2,-1) -- cycle;
        \node[popgreen] at (-1, 0) {Ribose};
        \draw[poppink, thick, fill=poppink!20] (-1, 1.3) circle (0.15) node[right=0.1cm, font=\tiny] {OH en 2'};
        \draw[poppink, thick, fill=poppink!20] (-1, -1.3) circle (0.15) node[right=0.1cm, font=\tiny] {OH en 3'};
        
        % Desoxyribose (ADN)
        \draw[popblue, thick, fill=popblue!10] (2,1) -- (3,1.5) -- (4,1) -- (4,-1) -- (3,-1.5) -- (2,-1) -- cycle;
        \node[popblue] at (3, 0) {Désoxyribose};
        \draw[poppink, thick, fill=white] (3, 1.3) circle (0.15) node[right=0.1cm, font=\tiny] {H en 2'};
        \draw[poppink, thick, fill=poppink!20] (3, -1.3) circle (0.15) node[right=0.1cm, font=\tiny] {OH en 3'};
        
        \node[align=center, font=\small] at (0, -1.8) {L'absence d'OH en 2' sur le désoxyribose\\rend l'ADN plus stable chimiquement\\que l'ARN (moins sensible à l'hydrolyse basique).};
    \end{scope}
\end{tikzpicture}
\legende{Structure d'un nucléotide (à gauche) et comparaison des sucres pentoses : ribose (ARN, OH en 2') et désoxyribose (ADN, H en 2'). La numérotation des carbones (1' à 5') est essentielle pour définir le sens 5' $\rightarrow$ 3' de la chaîne.}
\end{popfigure}

\courssub{Les bases azotées : alphabet de l'information génétique}

Les bases azotées sont des molécules hétérocycliques aromatiques, dérivées de la purine (double cycle) ou de la pyrimidine (cycle simple). Leur séquence le long de la chaîne porte l'information.

\begin{definition}[Bases azotées]
On distingue deux familles chimiques :
\begin{itemize}
    \item \textbf{Bases \cle{puriques}} (double cycle fusionné) : \textbf{Adénine (A)} et \textbf{Guanine (G)}. Présentes dans l'ADN \textbf{et} l'ARN.
    \item \textbf{Bases \cle{pyrimidiques}} (cycle simple) : 
    \begin{itemize}
        \item \textbf{Cytosine (C)} : présente dans l'ADN \textbf{et} l'ARN.
        \item \textbf{Thymine (T)} : \textbf{propre à l'ADN} (groupe méthyle \ce{-CH3} en position 5).
        \item \textbf{Uracile (U)} : \textbf{propre à l'ARN} (remplace la thymine, pas de groupe méthyle).
    \end{itemize}
\end{itemize}
La liaison base-sucre se fait par une liaison \textbf{N-glycosidique} : sur l'atome N9 pour les purines, N1 pour les pyrimidiques (carbone 1' du sucre).
\end{definition}

\begin{attention}[Piège fréquent : Confusion Thymine / Uracile]
\textbf{Jamais d'Uracile dans l'ADN (sauf rares ARN viraux ou erreurs de réparation), jamais de Thymine dans l'ARN (sauf ARNt matures modifiés).}
\begin{itemize}
    \item ADN : A, T, C, G.
    \item ARN : A, U, C, G.
\end{itemize}
Au Bac, écrire « L'ADN contient de l'uracile » ou « L'ARN contient de la thymine » est une \textbf{erreur éliminatoire} pour la notion correspondante.
\end{attention}

\courssub{Différences structurales majeures entre ADN et ARN}

\begin{center}
\begin{tabularx}{\linewidth}{|l|X|X|}\hline
\rowcolor{popblueL}
\textbf{Caractère} & \textbf{ADN (Acide Désoxyribonucléique)} & \textbf{ARN (Acide Ribonucléique)} \\ \hline
\textbf{Sucre} & \textbf{Désoxyribose} (H en 2') & \textbf{Ribose} (OH en 2') \\ \hline
\textbf{Bases pyrimidiques} & Cytosine (C), \textbf{Thymine (T)} & Cytosine (C), \textbf{Uracile (U)} \\ \hline
\textbf{Nombre de brins} & \textbf{Double brin} (hélice double) & \textbf{Simple brin} (le plus souvent) \\ \hline
\textbf{Structure 3D} & Double hélice régulière (B-ADN) & Structure secondaire complexe (brins-boucles, feuillets, pseudonœuds) \\ \hline
\textbf{Localisation principale} & Noyau (chromosomes), Mitochondries, Chloroplastes & Noyau (transcription), Cytoplasme (traduction) \\ \hline
\textbf{Stabilité chimique} & Élevée (absence OH 2') & Faible (OH 2' sensible à l'hydrolyse alcaline) \\ \hline
\textbf{Rôle principal} & Stockage stable de l'information héréditaire & Expression, transfert, traitement de l'information (ARNm, ARNt, ARNr) \\ \hline
\end{tabularx}
\end{center}

\courssub{La double hélice d'ADN : modèle de Watson et Crick (1953)}

Grâce aux données de diffraction aux rayons X de \textbf{Rosalind Franklin} (cliché « Photo 51 ») et de Maurice Wilkins, Watson et Crick proposent en 1953 le modèle de la double hélice.

\begin{propriete}[Modèle de la double hélice d'ADN (Watson-Crick, 1953)]
L'ADN est constitué de \textbf{deux brins polynucléotidiques} enroulés en \textbf{double hélice droite} (forme B la plus courante \textit{in vivo}).
\begin{enumerate}
    \item \textbf{Sens antiparallèle :} Les deux brins ont des sens chimiques opposés. L'un court dans le sens \textbf{5' $\rightarrow$ 3'} (phosphate libre en 5', OH libre en 3'), l'autre dans le sens \textbf{3' $\rightarrow$ 5'}.
    \item \textbf{Complémentarité des bases (Appariement de Watson-Crick) :} Les bases des deux brins s'associent par \textbf{liaisons hydrogène} selon des règles strictes :
    \begin{itemize}
        \item \textbf{Adénine (A) = Thymine (T)} : \textbf{2 liaisons hydrogène}.
        \item \textbf{Cytosine (C) $\equiv$ Guanine (G)} : \textbf{3 liaisons hydrogène}.
    \end{itemize}
    \item \textbf{Extérieur / Intérieur :} Les squelettes sucre-phosphate sont à l'extérieur (hydrophiles, chargées négativement), les bases sont à l'intérieur (hydrophobes, empilées par interactions $\pi$-$\pi$).
    \item \textbf{Géométrie :} Diamètre constant $\approx$ \SI{2}{nm}. Pas hélical $\approx$ \SI{3,4}{nm} (10 paires de bases par tour). Rainure majeure et rainure mineure (sites de fixation des protéines).
\end{enumerate}
\end{propriete}

\begin{popfigure}
\begin{tikzpicture}[scale=0.7, every node/.style={font=\small}]
    % Double helix schematic - 2 strands with base pairs
    \def\basepair#1#2#3#4#5{ % x, y, base1, base2, color1, color2
        \draw[thick, #5] (#1, #2) -- (#1+1.5, #2); % Backbone 1
        \draw[thick, #5] (#1, #2-0.6) -- (#1+1.5, #2-0.6); % Backbone 2
        % Base 1 (purine large)
        \draw[fill=#5!20, thick] (#1+0.2, #2-0.1) rectangle (#1+0.7, #2-0.5) node[midway] {#3};
        % Base 2 (pyrimidine narrow)
        \draw[fill=#5!20, thick] (#1+0.8, #2-0.1) rectangle (#1+1.3, #2-0.5) node[midway] {#4};
        % H-bonds
        \draw[popdark, dashed, thick] (#1+0.45, #2-0.1) -- (#1+0.45, #2-0.5);
        \draw[popdark, dashed, thick] (#1+1.05, #2-0.1) -- (#1+1.05, #2-0.5);
    }
    
    % Strand 1 (5' -> 3' top to bottom)
    % Strand 2 (3' -> 5' top to bottom)
    
    % Base pairs stack
    \basepair{0}{0}{A}{T}{popblue} % 1
    \basepair{0}{-0.8}{T}{A}{popblue} % 2
    \basepair{0}{-1.6}{C}{G}{poporange} % 3
    \basepair{0}{-2.4}{G}{C}{poporange} % 4
    \basepair{0}{-3.2}{A}{T}{popblue} % 5
    \basepair{0}{-4.0}{T}{A}{popblue} % 6
    
    % Backbone lines continuous
    \draw[popblue, very thick] (0.2, 0.1) -- (0.2, -4.7); % Left backbone
    \draw[popblue, very thick] (1.3, 0.1) -- (1.3, -4.7); % Right backbone
    
    % Direction labels
    \node[popblue, left] at (0.2, 0.3) {\textbf{5'}};
    \node[popblue, left] at (0.2, -4.9) {\textbf{3'}};
    \node[popblue, right] at (1.3, 0.3) {\textbf{3'}};
    \node[popblue, right] at (1.3, -4.9) {\textbf{5'}};
    
    % Antiparallel arrows
    \draw[->, popdark, thick] (1.8, -0.5) -- (1.8, -2.5) node[midway, right] {Sens 5'$\rightarrow$3'};
    \draw[->, popdark, thick] (-0.5, -2.5) -- (-0.5, -0.5) node[midway, left] {Sens 5'$\rightarrow$3'};
    
    % Dimensions
    \draw[<->, popmuted] (2.2, 0.1) -- (2.2, -0.7) node[midway, right] {\SI{0,34}{nm}};
    \draw[<->, popmuted] (-0.5, 0.1) -- (-0.5, -3.3) node[midway, left] {\SI{3,4}{nm} (10 paires)};
    \draw[<->, popmuted] (0.2, 0.3) -- (1.3, 0.3) node[midway, above] {\SI{2}{nm}};
    
    % Legend
    \node[align=left, font=\small] at (4, -2) {
        \textbf{Légende :}\\
        \textcolor{popblue}{■} Paire A=T (2 liaisons H)\\
        \textcolor{poporange}{■} Paire C$\equiv$G (3 liaisons H)\\
        \textcolor{popblue}{—} Squelette sucre-phosphate\\
        \textcolor{popdark}{- - -} Liaisons hydrogène\\
        Brins \textbf{antiparallèles}
    };
\end{tikzpicture}
\legende{Modèle simplifié de la double hélice d'ADN (forme B). Les deux brins sont antiparallèles (5'$\rightarrow$3' opposés). Les bases sont appariées par complémentarité (A=T, C$\equiv$G) via des liaisons hydrogène (pointillés). Le diamètre est constant (\SI{2}{nm}).}
\end{popfigure}

\begin{savaistu}[Rosalind Franklin : l'ombre de la « Photo 51 »]
Rosalind Franklin (1920--1958), cristallographe au King's College de Londres, a obtenu en 1952 le cliché de diffraction X n°51 de l'ADN (forme B). Ce cliché révélait sans ambiguïté une structure hélicale, un diamètre de \SI{2}{nm} et une périodicité de \SI{3,4}{nm}. Ses données, communiquées sans son accord explicite à Watson et Crick (via Wilkins et un rapport interne), furent décisives pour la construction du modèle. Watson, Crick et Wilkins reçurent le prix Nobel de physiologie ou médecine en 1962. Franklin, décédée en 1958 d'un cancer de l'ovaire (peut-être lié à son travail aux rayons X), ne put être récompensée (le Nobel n'est pas décerné à titre posthume). Son rôle est aujourd'hui pleinement reconnu comme fondateur de la biologie moléculaire.
\end{savaistu}

\courssub{Conséquence fondamentale : la complémentarité, base de la fidélité}

La règle de complémentarité (A=T, C$\equiv$G) a une conséquence majeure : \textbf{chaque brin porte l'information nécessaire pour reconstruire l'autre}.

\begin{propriete}[Règles de Chargaff (1950) — Conséquence de la complémentarité]
Dans l'ADN double brin de n'importe quel organisme :
\begin{itemize}
    \item \textbf{\%A = \%T} et \textbf{\%C = \%G} (égalité des bases complémentaires).
    \item \textbf{\%A + \%T + \%C + \%G = 100\%}.
    \item \textbf{(\%A + \%G) = (\%T + \%C)} (égalité purines = pyrimidiques).
    \item Le rapport \textbf{(\%A + \%T) / (\%C + \%G)} varie selon les espèces (signature taxonomique), mais est constant au sein d'une espèce.
\end{itemize}
Ces règles, établies par analyse chimique (hydrolyse + chromatographie) \textbf{avant} la découverte de la structure, ont guidé Watson et Crick.
\end{propriete}

\begin{methode}[Calculer la composition en bases d'un ADN double brin]
\textbf{Donnée :} Un pourcentage de base (ex: \%A = 30\%).
\textbf{Étapes :}
\begin{enumerate}
    \item Appliquer \textbf{\%A = \%T} $\rightarrow$ \%T = 30\%.
    \item Calculer la somme restante : 100\% - (\%A + \%T) = 100\% - 60\% = 40\%.
    \item Appliquer \textbf{\%C = \%G} $\rightarrow$ \%C = \%G = 40\% / 2 = 20\%.
    \item Vérifier : 30+30+20+20 = 100\%.
\end{enumerate}
\textbf{Attention :} Cette méthode ne s'applique qu'à l'\textbf{ADN double brin}. Pour un brin simple (ARN, ADN viral simple brin), les règles de Chargaff \textbf{ne s'appliquent pas}.
\end{methode}

\begin{exoresolu}[Calcul de composition basique]
\textbf{Énoncé :} Un fragment d'ADN double brin contient 22\% de cytosine (C). Déterminer les pourcentages des trois autres bases.
\textbf{Solution :}
\begin{align*}
    \%C &= 22\% \\
    \%G &= \%C = 22\% \quad \text{(règle de Chargaff)} \\
    \%C + \%G &= 44\% \\
    \%A + \%T &= 100\% - 44\% = 56\% \\
    \%A &= \%T = 56\% / 2 = 28\% \\
    \textbf{Bilan :} &\quad \%A = 28\%,\; \%T = 28\%,\; \%C = 22\%,\; \%G = 22\%.
\end{align*}
\end{exoresolu}

\begin{exercice}[Application des règles de Chargaff]
\textbf{Énoncé :} Le génome d'une bactérie thermophile contient 38\% de bases puriques (A+G).
\begin{enumerate}
    \item Calculer le pourcentage de bases pyrimidiques.
    \item Si le pourcentage d'adénine est de 18\%, déterminer la composition complète en bases.
    \item Cette bactérie subit une mutation qui double la quantité d'adénine sans changer la quantité totale de bases. Quel est le nouveau \%G ? (Indication : la complémentarité est conservée).
\end{enumerate}
\end{exercice}

\begin{corrige}[Exercice n°1]
\begin{enumerate}
    \item Purines (A+G) = 38\% $\rightarrow$ Pyrimidiques (T+C) = 100\% - 38\% = \textbf{62\%}.
    \item \%A = 18\% $\rightarrow$ \%T = 18\%. A+T = 36\%. C+G = 64\%. \%G = (A+G) - \%A = 38\% - 18\% = 20\%. \%C = \%G = 20\%.
    \textbf{Composition :} A=18\%, T=18\%, G=20\%, C=20\%.
    \item Nouveau \%A = 36\% $\rightarrow$ Nouveau \%T = 36\%. Nouveau A+T = 72\%. Nouveau C+G = 28\%. Nouveau \%G = \%C = 14\%.
\end{enumerate}
\end{corrige}

\courssub{Les trois types d'ARN : structure et rôles}

Contrairement à l'ADN (stockage), l'ARN assure des fonctions variées : messager, adaptateur, catalytique/structurel. Tous sont des polymères de ribonucléotides (Ribose, A, U, C, G), généralement simple brin, mais capables de s'apparier localement (structure secondaire).

\begin{definition}[Les trois grands types d'ARN cellulaires]
\begin{enumerate}
    \item \textbf{ARN messager (ARNm) :} Copie fidèle d'un gène (transcrit). Porte le code génétique sous forme de \textbf{codons} (triplets de bases). Existe en versions \textbf{précurseur (pré-ARNm / hnRNA)} avec introns/exons (eucaryotes) et \textbf{mature} (épissé, coiffé \ce{7mG}, polyadénylé \ce{poly(A)}). Durée de vie courte (minutes à heures). \textbf{Rôle :} Intermédiaire entre gène et protéine.
    \item \textbf{ARN de transfert (ARNt) :} Petit ARN (\SI{\approx 70-90}{nucléotides}), structure en \textbf{trèfle} (secondaire) / \textbf{L} (tertiaire). Deux sites fonctionnels clés :
    \begin{itemize}
        \item \textbf{Anticodon} (boucle 1) : triplet complémentaire au codon de l'ARNm.
        \item \textbf{Extrémité 3' CCA} (bras accepteur) : site de fixation covalent de l'acide aminé correspondant (esterification par aminoacyl-ARNt synthétase).
    \end{itemize}
    Contient de nombreuses \textbf{bases modifiées} (ex: pseudouridine $\Psi$, dihydrouridine D, inosine I) stabilisant la structure. \textbf{Rôle :} Adaptateur traduisant le langage nucléique (codons) en langage protéique (acides aminés).
    \item \textbf{ARN ribosomique (ARNr) :} Composant majeur des \textbf{ribosomes} (jusqu'à 60\% de la masse). Existe en plusieurs tailles (ex: chez eucaryotes \textbf{28S, 18S, 5.8S, 5S} ; chez procaryotes \textbf{23S, 16S, 5S}). Structure secondaire très complexe (domaines hélicoïdaux). \textbf{Rôle :} Structurel (échafaudage) \textbf{et catalytique} (activité \textbf{ribozymique} : peptidyl-transférase pour la liaison peptidique, assurée par l'ARNr 23S/28S).
\end{enumerate}
\end{definition}

\begin{center}
\begin{tabularx}{\linewidth}{|l|X|X|X|X|}\hline
\rowcolor{popblueL}
\textbf{Caractère} & \textbf{ARNm} & \textbf{ARNt} & \textbf{ARNr} & \textbf{Autres ARNs (ex: snARN, miARN)} \\ \hline
\textbf{Taille} & Variable (centaines à milliers nt) & Petite (\SI{\approx 70-90}{nt}) & Grande (\SI{\approx 120-5000}{nt}) & Variable (souvent petits \SI{< 300}{nt}) \\ \hline
\textbf{Structure} & Linéaire, simple brin, peu de structure secondaire stable & Trèfle (2D) / L (3D), nombreuses bases modifiées & Complexe, domaines hélicoïdaux imbriqués & Spécifique à la fonction \\ \hline
\textbf{Localisation synthèse} & Noyau (Pol II) & Noyau (Pol III) & Nucléole (Pol I \& III) & Noyau (Pol II/III) \\ \hline
\textbf{Localisation fonction} & Cytoplasme (ribosomes) & Cytoplasme (ribosomes) & Cytoplasme (ribosomes) & Noyau / Cytoplasme \\ \hline
\textbf{Fonction principale} & \textbf{Codage} : Véhicule l'info génétique (codons) & \textbf{Adaptation} : Apporte l'acide aminé (anticodon) & \textbf{Catalyse/Structure} : Usine à protéines (ribozyme) & \textbf{Régulation} : Épissage, silencing, etc. \\ \hline
\textbf{Stabilité / Durée de vie} & Faible (min à h), dégradation contrôlée & Élevée (jours), recyclée & Très élevée (stable dans ribosome) & Variable \\ \hline
\textbf{Modifications post-transcriptionnelles} & Coiffe 5' (\ce{7mG}), Queue \ce{poly(A)}, Épissage & Bases modifiées ($\Psi$, D, I, m\ce{G}), ajout CCA 3' & Méthylation, pseudouridylation, clivage & Spécifiques \\ \hline
\end{tabularx}
\end{center}

\begin{popfigure}
\begin{tikzpicture}[scale=0.85, every node/.style={font=\small}]
    % ARNm
    \begin{scope}[xshift=-7cm]
        \node at (0, 2.5) {\textbf{ARNm (eucaryote mature)}};
        \draw[popblue, thick, ->] (-3,0) -- (3,0) node[midway, above] {5' $\rightarrow$ 3'};
        \draw[popblue, thick, fill=popblue!30] (-3, -0.2) rectangle (-2.2, 0.2) node[midway, font=\tiny] {Coiffe \ce{7mG}};
        \draw[popblue, thick, fill=popblue!10] (-2.2, -0.2) rectangle (-0.5, 0.2) node[midway, font=\tiny] {5' UTR};
        \foreach \x in {-0.5, 0.2, 0.9, 1.6} {
            \draw[poporange, thick, fill=poporange!20] (\x, -0.2) rectangle (\x+0.6, 0.2) node[midway, font=\tiny] {Exon};
            \ifdim\x pt<1.6pt \draw[popmuted, thick, fill=popmuted!20] (\x+0.6, -0.2) rectangle (\x+0.9, 0.2) node[midway, font=\tiny] {Intron (retiré)}; \fi
        }
        \draw[popblue, thick, fill=popblue!10] (2.5, -0.2) rectangle (3, 0.2) node[midway, font=\tiny] {3' UTR};
        \draw[popblue, thick, fill=popblue!30] (3, -0.3) rectangle (3.5, 0.3) node[midway, rotate=90, font=\tiny] {Poly(A)};
    \end{scope}
    
    % ARNt (Cloverleaf)
    \begin{scope}[xshift=0cm]
        \node at (0, 2.5) {\textbf{ARNt (Trèfle)}};
        % Acceptor stem
        \draw[popgreen, thick] (-0.5, -1.5) -- (-0.5, -0.5) node[midway, left] {3' CCA};
        \draw[popgreen, thick] (0.5, -1.5) -- (0.5, -0.5);
        \draw[popgreen, thick] (-0.5, -0.5) -- (0.5, -0.5);
        % D arm
        \draw[popgreen, thick] (-0.5, -0.5) -- (-1.2, 0.2) -- (-0.5, 0.9) node[midway, left, font=\tiny] {Bras D};
        % Anticodon arm
        \draw[popgreen, thick] (0.5, -0.5) -- (1.2, 0.2) -- (0.5, 0.9) node[midway, right, font=\tiny] {Anticodon};
        % T$\Psi$C arm
        \draw[popgreen, thick] (-0.5, 0.9) -- (-0.5, 1.5);
        \draw[popgreen, thick] (0.5, 0.9) -- (0.5, 1.5);
        \draw[popgreen, thick] (-0.5, 1.5) -- (0.5, 1.5) node[midway, above, font=\tiny] {Bras T$\Psi$C};
        % Variable loop
        \draw[popgreen, thick] (-0.5, 0.9) .. controls (-0.8, 1.2) and (0.8, 1.2) .. (0.5, 0.9);
        \node[popgreen, font=\tiny] at (0, 1.1) {Boucle variable};
    \end{scope}
    
    % ARNr (Ribosome schematic)
    \begin{scope}[xshift=7cm]
        \node at (0, 2.5) {\textbf{Ribosome (ex: Eucaryote 80S)}};
        % Grosse sous-unité 60S
        \draw[poppurple, thick, fill=poppurple!10, rounded corners=5pt] (-1.5, -0.5) rectangle (1.5, 1.5);
        \node[poppurple] at (0, 1.8) {\textbf{Grande SU (60S)}};
        \node[poppurple, align=center] at (0, 0.5) {ARNr \textbf{28S} \\ ARNr \textbf{5.8S} \\ ARNr \textbf{5S} \\ $\approx$ 45 protéines};
        \node[poppurple, font=\tiny, align=center] at (0, -0.8){Site P / Site A / Site E \\ Activité Peptidyl-Transférase (Ribozyme)};
        
        % Petite sous-unité 40S
        \draw[poppink, thick, fill=poppink!10, rounded corners=5pt] (-1.5, -2.5) rectangle (1.5, -0.5);
        \node[poppink] at (0, -2.8) {\textbf{Petite SU (40S)}};
        \node[poppink, align=center] at (0, -1.5) {ARNr \textbf{18S} \\ $\approx$ 33 protéines};
        \node[poppink, font=\tiny, align=center] at (0, -3.2){Décodage (codon-anticodon) \\ Fixation ARNm};
    \end{scope}
\end{tikzpicture}
\legende{Représentations schématiques des trois ARNs majeurs. Haut gauche : ARNm eucaryote mature (coiffe 5', UTRs, exons codants, queue poly(A) 3'). Haut centre : ARNt en structure secondaire « trèfle » (bras accepteur 3' CCA, anticodon, bras D, bras T$\Psi$C). Droite : Assemblage des ARNr et protéines ribosomiques dans les deux sous-unités du ribosome 80S eucaryote (60S + 40S = 80S ; S = coefficient de sédimentation de Svedberg, non additive).}
\end{popfigure}

\begin{aretenir}[Bilan : L'architecture de l'information]
\begin{itemize}
    \item L'\textbf{ADN} (Désoxyribose, A,T,C,G, double brin antiparallèle, hélice double) est le support \textbf{stable} de l'information héréditaire (Griffith $\rightarrow$ Avery $\rightarrow$ Hershey-Chase).
    \item L'\textbf{ARN} (Ribose, A,U,C,G, simple brin) en assure l'\textbf{expression} sous trois formes principales : \textbf{ARNm} (message/codons), \textbf{ARNt} (adaptateur/anticodon+acide aminé), \textbf{ARNr} (structure/catalyse ribosome).
    \item La \textbf{complémentarité des bases} (A=T/U, C$\equiv$G, liaisons H) est le principe universel qui permet : la réplication (fidélité), la transcription (fidélité), la traduction (spécificité codon-anticodon).
    \item Les \textbf{règles de Chargaff} (\%A=\%T, \%C=\%G) sont la signature chimique de la double hélice et un outil de calcul incontournable au Bac.
\end{itemize}
\end{aretenir}

\coursec{II. La duplication de l'ADN : la réplication semi-conservative}

Nous venons de découvrir l'architecture de la molécule d'ADN : une double hélice formée de deux brins complémentaires et antiparallèles, unis par des liaisons hydrogène entre bases appariées (A avec T, G avec C). Cette structure, élucidée par Watson et Crick en 1953, contient en germe une idée géniale : si chaque brin est le « négatif » de l'autre, alors chacun peut servir de \cle{matrice} pour fabriquer l'autre. C'est exactement ce que la cellule réalise avant chaque division. Cette section explique comment, et comment les scientifiques l'ont démontré.

\courssub{Le problème : doubler le stock d'ADN avant chaque division}

\textbf{Observation.} Une cellule de ton foie, une cellule de la moelle osseuse, une cellule de méristème de racine d'igname : toutes se divisent, et chaque cellule fille reçoit un programme génétique \emph{complet et identique} à celui de la cellule mère. Or la division partage le cytoplasme en deux, mais on ne peut pas couper l'ADN en deux moitiés : chaque cellule fille a besoin de la \emph{totalité} de l'information génétique.

\textbf{Problème posé.} Comment la cellule fait-elle pour passer d'un stock d'ADN à un stock double, \emph{sans erreur}, avant de se diviser ?

\textbf{Solution.} La cellule réalise la \cle{réplication} (ou duplication) de son ADN au cours de la \cle{phase S} (S pour \emph{synthèse}) de l'\cle{interphase}, c'est-à-dire la période qui précède la division proprement dite. La quantité d'ADN passe ainsi de $Q$ à $2Q$ avant la mitose (ou la méiose), puis revient à $Q$ dans chaque cellule fille après la division.

\begin{attention}[Le piège classique du Bac]
La réplication n'a \textbf{pas} lieu « pendant » la division cellulaire. Elle a lieu \textbf{en interphase}, pendant la phase S, c'est-à-dire \emph{avant} la mitose ou la méiose. Quand la division commence, la quantité d'ADN est déjà doublée ($2Q$) et chaque chromosome est déjà formé de deux chromatides. Écrire « l'ADN se réplique pendant la mitose » est une erreur éliminatoire dans une dissertation.
\end{attention}

\courssub{Les trois hypothèses historiques}

Une fois la structure en double hélice connue, une question s'impose aux biologistes des années 1950 : comment les deux brins sont-ils distribués dans les molécules filles ? Trois modèles sont en compétition.

\begin{itemize}
\item \textbf{Modèle conservatif} : la molécule parentale reste intacte et sert simplement de gabarit ; une molécule entièrement nouvelle est synthétisée à côté. Après une génération, on obtient une molécule « 100 \% parentale » et une molécule « 100 \% néoformée ».
\item \textbf{Modèle semi-conservatif} : les deux brins parentaux se séparent, et chacun sert de matrice à la synthèse d'un brin complémentaire. Chaque molécule fille est donc formée d'\emph{un brin parental} et d'\emph{un brin néoformé}.
\item \textbf{Modèle dispersif} : les molécules parentales sont fragmentées ; chaque molécule fille est une mosaïque de segments parentaux et de segments néoformés dispersés le long des deux brins.
\end{itemize}

\begin{popfigure}
\begin{center}
\begin{tikzpicture}[scale=0.9, every node/.style={font=\small}]
% ===== Hypothèse conservative =====
\begin{scope}
\node[font=\bfseries, popdark] at (0,4.6) {Conservative};
% mère
\draw[line width=2.2pt, popblue] (-1.2,3.6) -- (1.2,3.6);
\draw[line width=2.2pt, popblue] (-1.2,3.2) -- (1.2,3.2);
\node[popdark] at (0,2.75) {\footnotesize Molécule mère};
\draw[->, thick, popmuted] (0,2.4) -- (0,1.9);
% filles
\draw[line width=2.2pt, popblue] (-2.6,1.4) -- (-0.6,1.4);
\draw[line width=2.2pt, popblue] (-2.6,1.0) -- (-0.6,1.0);
\draw[line width=2.2pt, poporange] (0.6,1.4) -- (2.6,1.4);
\draw[line width=2.2pt, poporange] (0.6,1.0) -- (2.6,1.0);
\node[popdark] at (-1.6,0.55) {\footnotesize 100 \% parentale};
\node[popdark] at (1.6,0.55) {\footnotesize 100 \% néoformée};
\end{scope}
% ===== Hypothèse semi-conservative =====
\begin{scope}[xshift=6.2cm]
\node[font=\bfseries, popdark] at (0,4.6) {Semi-conservative};
\draw[line width=2.2pt, popblue] (-1.2,3.6) -- (1.2,3.6);
\draw[line width=2.2pt, popblue] (-1.2,3.2) -- (1.2,3.2);
\node[popdark] at (0,2.75) {\footnotesize Molécule mère};
\draw[->, thick, popmuted] (0,2.4) -- (0,1.9);
\draw[line width=2.2pt, popblue] (-2.6,1.4) -- (-0.6,1.4);
\draw[line width=2.2pt, poporange] (-2.6,1.0) -- (-0.6,1.0);
\draw[line width=2.2pt, popblue] (0.6,1.4) -- (2.6,1.4);
\draw[line width=2.2pt, poporange] (0.6,1.0) -- (2.6,1.0);
\node[popdark] at (-1.6,0.55) {\footnotesize 1 brin parental};
\node[popdark] at (-1.6,0.2) {\footnotesize + 1 brin néoformé};
\node[popdark] at (1.6,0.55) {\footnotesize 1 brin parental};
\node[popdark] at (1.6,0.2) {\footnotesize + 1 brin néoformé};
\end{scope}
% ===== Hypothèse dispersive =====
\begin{scope}[xshift=12.4cm]
\node[font=\bfseries, popdark] at (0,4.6) {Dispersive};
\draw[line width=2.2pt, popblue] (-1.2,3.6) -- (1.2,3.6);
\draw[line width=2.2pt, popblue] (-1.2,3.2) -- (1.2,3.2);
\node[popdark] at (0,2.75) {\footnotesize Molécule mère};
\draw[->, thick, popmuted] (0,2.4) -- (0,1.9);
% fille 1 : segments alternés
\draw[line width=2.2pt, popblue] (-2.6,1.4) -- (-2.1,1.4);
\draw[line width=2.2pt, poporange] (-2.1,1.4) -- (-1.6,1.4);
\draw[line width=2.2pt, popblue] (-1.6,1.4) -- (-1.1,1.4);
\draw[line width=2.2pt, poporange] (-1.1,1.4) -- (-0.6,1.4);
\draw[line width=2.2pt, poporange] (-2.6,1.0) -- (-2.1,1.0);
\draw[line width=2.2pt, popblue] (-2.1,1.0) -- (-1.6,1.0);
\draw[line width=2.2pt, poporange] (-1.6,1.0) -- (-1.1,1.0);
\draw[line width=2.2pt, popblue] (-1.1,1.0) -- (-0.6,1.0);
% fille 2
\draw[line width=2.2pt, poporange] (0.6,1.4) -- (1.1,1.4);
\draw[line width=2.2pt, popblue] (1.1,1.4) -- (1.6,1.4);
\draw[line width=2.2pt, poporange] (1.6,1.4) -- (2.1,1.4);
\draw[line width=2.2pt, popblue] (2.1,1.4) -- (2.6,1.4);
\draw[line width=2.2pt, popblue] (0.6,1.0) -- (1.1,1.0);
\draw[line width=2.2pt, poporange] (1.1,1.0) -- (1.6,1.0);
\draw[line width=2.2pt, popblue] (1.6,1.0) -- (2.1,1.0);
\draw[line width=2.2pt, poporange] (2.1,1.0) -- (2.6,1.0);
\node[popdark] at (-1.6,0.55) {\footnotesize mosaïque};
\node[popdark] at (1.6,0.55) {\footnotesize mosaïque};
\end{scope}
% Légende couleur
\draw[line width=2.2pt, popblue] (3.4,-0.6) -- (4.2,-0.6); \node[right, popdark] at (4.3,-0.6) {\footnotesize brin parental};
\draw[line width=2.2pt, poporange] (7.4,-0.6) -- (8.2,-0.6); \node[right, popdark] at (8.3,-0.6) {\footnotesize brin néoformé};
\end{tikzpicture}
\end{center}
\legende{Les trois hypothèses historiques de réplication de l'ADN. Chaque trait représente un brin polynucléotidique ; en bleu, le brin parental (préexistant) ; en orange, le brin nouvellement synthétisé. Seule l'expérience permettra de départager ces trois modèles.}
\end{popfigure}

Comment trancher ? Il faut pouvoir \emph{distinguer} l'ADN parental de l'ADN néoformé, puis suivre leur répartition au fil des générations. C'est l'exploit réalisé par Matthew Meselson et Franklin Stahl en 1958, dans ce qui est souvent appelé « la plus belle expérience de biologie moléculaire ».

\courssub{L'expérience de Meselson et Stahl (1958)}

\begin{experience}[Meselson et Stahl : le marquage isotopique de l'ADN]
\textbf{Principe.} L'azote est un constituant des bases azotées de l'ADN. Il existe deux isotopes : l'azote « léger » $^{14}$N (le plus abondant dans la nature) et l'azote « lourd » $^{15}$N, non radioactif mais plus dense. Un ADN contenant du $^{15}$N est \emph{plus dense} qu'un ADN contenant du $^{14}$N : on peut donc les séparer par \cle{centrifugation en gradient de chlorure de césium} (CsCl). Après centrifugation prolongée, chaque ADN migre jusqu'à la position du gradient dont la densité égale la sienne et forme une \emph{bande} visible : plus l'ADN est lourd, plus la bande est basse dans le tube.

\textbf{Protocole.}
\begin{enumerate}
\item Des bactéries \emph{Escherichia coli} sont cultivées pendant de nombreuses générations sur un milieu où la seule source d'azote est le $^{15}$N : tout leur ADN devient « lourd » ($^{15}$N/$^{15}$N).
\item Les bactéries sont ensuite transférées sur un milieu contenant uniquement du $^{14}$N : tout ADN \emph{nouvellement synthétisé} sera « léger ».
\item On prélève des échantillons à la génération 0 (avant le transfert), après une génération (un cycle de réplication) et après deux générations, puis on centrifuge l'ADN extrait dans un gradient de CsCl.
\end{enumerate}

\textbf{Résultats observés.}
\begin{itemize}
\item \textbf{Génération 0} : une seule bande, en position basse (ADN lourd $^{15}$N/$^{15}$N).
\item \textbf{Génération 1} : une seule bande, en position \emph{intermédiaire} (ADN hybride $^{15}$N/$^{14}$N). Aucune bande lourde, aucune bande légère.
\item \textbf{Génération 2} : deux bandes d'intensité égale, l'une intermédiaire ($^{15}$N/$^{14}$N), l'autre haute, c'est-à-dire légère ($^{14}$N/$^{14}$N).
\end{itemize}
\end{experience}

\begin{popfigure}
\begin{center}
\begin{tikzpicture}[scale=1.0, every node/.style={font=\small}]
% gradients : trois tubes
\foreach \x/\gen in {0/Génération 0, 4.6/Génération 1, 9.2/Génération 2}{
  % tube
  \draw[thick, popdark, rounded corners=2pt] (\x-1.1,0) rectangle (\x+1.1,4.2);
  % gradient de densité (fond)
  \fill[popblue!8] (\x-1.08,0.02) rectangle (\x+1.08,4.18);
  \node[popdark, font=\bfseries] at (\x,4.7) {\gen};
  % repères de densité
  \node[left, popmuted, font=\footnotesize] at (\x-1.15,3.4) {légère};
  \node[left, popmuted, font=\footnotesize] at (\x-1.15,2.2) {intermédiaire};
  \node[left, popmuted, font=\footnotesize] at (\x-1.15,1.0) {lourde};
}
% bandes génération 0 : lourde
\fill[popblue] (-1.08,0.85) rectangle (1.08,1.15);
\node[popdark, font=\footnotesize] at (0,0.35) {$^{15}$N/$^{15}$N};
% bandes génération 1 : intermédiaire
\fill[poppurple] (3.52,2.05) rectangle (5.68,2.35);
\node[popdark, font=\footnotesize] at (4.6,1.55) {$^{15}$N/$^{14}$N};
% bandes génération 2 : intermédiaire + légère
\fill[poppurple] (8.12,2.05) rectangle (10.28,2.35);
\fill[poporange] (8.12,3.25) rectangle (10.28,3.55);
\node[popdark, font=\footnotesize] at (9.2,1.55) {$^{15}$N/$^{14}$N};
\node[popdark, font=\footnotesize] at (9.2,3.85) {$^{14}$N/$^{14}$N};
% flèche temps
\draw[->, very thick, popmuted] (-1.3,-0.7) -- (10.5,-0.7) node[right, popdark] {temps (générations sur milieu $^{14}$N)};
\end{tikzpicture}
\end{center}
\legende{Résultats de l'expérience de Meselson et Stahl. Position des bandes d'ADN après centrifugation en gradient de CsCl. La densité du milieu augmente vers le bas du tube : une bande basse correspond à un ADN dense (lourd), une bande haute à un ADN léger.}
\end{popfigure}

\textbf{Interprétation (démarche d'investigation).}
\begin{itemize}
\item Le modèle \textbf{conservatif} prédit, à la génération 1, \emph{deux} bandes : une lourde (molécule parentale intacte) et une légère (molécule entièrement néoformée). Or on observe une \emph{seule} bande intermédiaire : le modèle conservatif est \textbf{réfuté}.
\item Le modèle \textbf{dispersif} prédit des molécules dont la densité diminue \emph{progressivement} à chaque génération : à la génération 2, on attendrait une seule bande, plus haute que la bande intermédiaire. Or on observe \emph{deux} bandes distinctes, dont une reste exactement intermédiaire : le modèle dispersif est \textbf{réfuté}.
\item Le modèle \textbf{semi-conservatif} prédit exactement les résultats observés : à la génération 1, chaque molécule fille garde un brin $^{15}$N et acquiert un brin $^{14}$N (bande hybride unique) ; à la génération 2, les brins $^{15}$N servent de matrice à des brins $^{14}$N (moitié des molécules hybrides) et les brins $^{14}$N de la génération 1 donnent des molécules $^{14}$N/$^{14}$N (moitié des molécules légères).
\end{itemize}

\begin{methode}[Départager les hypothèses à partir des bandes : la méthode type Bac]
Face à un document de ce type, procède toujours en quatre étapes :
\begin{enumerate}
\item \textbf{Prédire} : pour chaque hypothèse, écris la composition isotopique des molécules attendues à chaque génération et la position des bandes correspondantes.
\item \textbf{Comparer} : confronte chaque prédiction aux bandes réellement observées, génération par génération.
\item \textbf{Éliminer} : une seule observation contradictoire suffit à réfuter une hypothèse (ici, l'absence de bande lourde à la génération 1 élimine le modèle conservatif).
\item \textbf{Conclure} : l'hypothèse compatible avec \emph{tous} les résultats est retenue ; formule alors la conclusion générale (réplication semi-conservative).
\end{enumerate}
\end{methode}

\begin{definition}[Réplication semi-conservative]
La réplication de l'ADN est dite \cle{semi-conservative} car chaque molécule d'ADN fille est constituée d'\textbf{un brin parental} (conservé, servant de matrice) et d'\textbf{un brin néoformé} (nouvellement synthétisé par appariement complémentaire des bases). La moitié de chaque molécule parentale est ainsi conservée dans chaque molécule fille.
\end{definition}

\begin{exemplebox}[Combien de molécules contiennent encore un brin parental ?]
Partons d'\textbf{une} molécule d'ADN. À chaque cycle de réplication, le nombre de molécules \emph{double} : après $n$ cycles, on obtient $2^{n}$ molécules. Mais les deux brins parentaux, eux, ne se multiplient pas : ils sont simplement distribués, chacun dans une molécule fille, et le restent à jamais. Après 3 cycles : $2^{3} = 8$ molécules, dont seulement \textbf{2} contiennent encore un brin parental (celles qui hébergent les deux brins d'origine). Après 10 cycles : $2^{10} = 1024$ molécules, dont toujours seulement 2 marquées par un brin parental. C'est exactement ce qu'illustre la dilution progressive de la bande hybride $^{15}$N/$^{14}$N chez Meselson et Stahl au-delà de la génération 2.
\end{exemplebox}

\courssub{Les acteurs et les étapes de la réplication}

\textbf{Les acteurs moléculaires.} La réplication est un chantier collectif faisant intervenir plusieurs enzymes :
\begin{itemize}
\item l'\cle{ADN hélicase} : enzyme de déroulement qui rompt les liaisons hydrogène entre bases complémentaires et sépare les deux brins parentaux, ouvrant ainsi la double hélice ;
\item l'\cle{ADN polymérase} : enzyme clé qui assemble les nucléotides libres du milieu intranucléique en les appariant à la matrice (A face à T, G face à C) et en catalysant la formation des liaisons phosphodiester du nouveau brin ;
\item les \cle{amorces} : courtes séquences qui fournissent à l'ADN polymérase un point de départ, car cette enzyme ne peut allonger qu'un brin déjà amorcé.
\end{itemize}

\begin{attention}[Le sens de synthèse : 5' vers 3']
L'ADN polymérase ne synthétise le nouveau brin que dans le sens \cle{5' $\rightarrow$ 3'} : chaque nucléotide ajouté est fixé par son carbone 5' sur l'extrémité 3' libre du brin en construction. Comme les deux brins parentaux sont \emph{antiparallèles}, la lecture des deux matrices se fait dans des directions opposées au niveau de la fourche. Retiens : on lit la matrice en 3' $\rightarrow$ 5', on synthétise le brin néoformé en 5' $\rightarrow$ 3'.
\end{attention}

\textbf{Les étapes.} On distingue schématiquement trois phases :
\begin{enumerate}
\item \textbf{Ouverture de la double hélice} : l'hélicase désapparie les bases et sépare les deux brins parentaux, créant une zone en « Y » appelée \cle{fourche de réplication} (ou œil de réplication). Chez les eucaryotes, la réplication démarre en de nombreux points d'origine le long de chaque chromosome.
\item \textbf{Appariement des nucléotides} : sur chaque brin matrice exposé, les nucléotides libres du nucléoplasme viennent se fixer par complémentarité (A--T, G--C) grâce à des liaisons hydrogène.
\item \textbf{Formation des liaisons covalentes} : l'ADN polymérase relie les nucléotides entre eux (liaisons phosphodiester), édifiant le squelette sucre--phosphate du brin néoformé. La fourche progresse le long de la molécule jusqu'à réplication complète.
\end{enumerate}

\begin{popfigure}
\begin{center}
\begin{tikzpicture}[scale=1.0, every node/.style={font=\small}]
% Brins parentaux non ouverts (à gauche)
\draw[line width=2.2pt, popblue] (-4.5,0.5) -- (-1.6,0.5);
\draw[line width=2.2pt, popblue] (-4.5,-0.5) -- (-1.6,-0.5);
% petits ponts (bases appariées) dans la zone non ouverte
\foreach \x in {-4.2,-3.8,-3.4,-3.0,-2.6,-2.2,-1.8}{
  \draw[popblue!50, thin] (\x,-0.5) -- (\x,0.5);
}
% Ouverture en Y
\draw[line width=2.2pt, popblue] (-1.6,0.5) .. controls (-0.8,0.9) and (-0.2,1.3) .. (1.2,1.5);
\draw[line width=2.2pt, popblue] (-1.6,-0.5) .. controls (-0.8,-0.9) and (-0.2,-1.3) .. (1.2,-1.5);
% prolongement des matrices
\draw[line width=2.2pt, popblue] (1.2,1.5) -- (4.5,1.5);
\draw[line width=2.2pt, popblue] (1.2,-1.5) -- (4.5,-1.5);
% brins néoformés
\draw[line width=2.2pt, poporange] (0.2,1.05) -- (4.5,1.05);
\draw[line width=2.2pt, poporange] (0.2,-1.05) -- (4.5,-1.05);
% ponts hydrogène néoformés
\foreach \x in {0.6,1.0,1.4,1.8,2.2,2.6,3.0,3.4,3.8,4.2}{
  \draw[poporange!45, thin] (\x,1.05) -- (\x,1.5);
  \draw[poporange!45, thin] (\x,-1.05) -- (\x,-1.5);
}
% hélicase à la fourche
\node[draw, circle, fill=popgreenL, inner sep=2.5pt, font=\footnotesize\bfseries, popdark] at (-1.35,0) {H};
% sens de progression de la fourche
\draw[->, very thick, popgreen] (-0.9,0) -- (0.5,0);
\node[popgreen, font=\footnotesize] at (-0.2,-0.35) {sens de progression};
% sens de synthèse
\draw[->, thick, poporange] (3.2,0.72) -- (4.4,0.72) node[right, font=\footnotesize, popdark] {5' $\rightarrow$ 3'};
\draw[->, thick, poporange] (3.2,-0.72) -- (4.4,-0.72) node[right, font=\footnotesize, popdark] {5' $\rightarrow$ 3'};
% légendes
\node[popblue, font=\footnotesize] at (-3.0,1.0) {double hélice parentale};
\node[popblue, font=\footnotesize, align=center] at (5.6,1.5) {brin matrice\\ parental};
\node[poporange, font=\footnotesize, align=center] at (5.9,1.05) {brin néoformé};
\node[popdark, font=\footnotesize, align=center] at (-1.35,-1.15) {hélicase : ouverture\\ (fourche de réplication)};
\end{tikzpicture}
\end{center}
\legende{Schéma de la fourche de réplication. L'hélicase (H) sépare les deux brins parentaux (bleu) ; sur chaque brin matrice, l'ADN polymérase assemble un brin néoformé (orange) par appariement complémentaire des bases, dans le sens 5' $\rightarrow$ 3'. La fourche progresse le long de la molécule.}
\end{popfigure}

\begin{exoresolu}[Composition isotopique de l'ADN à la génération 3]
\textbf{Énoncé.} Des bactéries \emph{E. coli} cultivées sur milieu $^{15}$N sont transférées sur milieu $^{14}$N. On laisse se dérouler trois générations. Indique le nombre de molécules d'ADN par cellule initiale, la proportion de molécules hybrides ($^{15}$N/$^{14}$N) et la proportion de molécules légères ($^{14}$N/$^{14}$N) à la troisième génération. Prévois l'aspect du tube après centrifugation en gradient de CsCl.
\tcblower
\textbf{Solution détaillée.} Le nombre de molécules double à chaque génération : après 3 générations, $2^{3} = 8$ molécules par cellule initiale.

Les deux brins parentaux $^{15}$N ne se dupliquent pas : ils sont conservés chacun dans une molécule fille. Il y a donc toujours exactement \textbf{2 molécules hybrides} $^{15}$N/$^{14}$N, quelle que soit la génération. Les $8 - 2 = 6$ autres molécules sont entièrement légères $^{14}$N/$^{14}$N.

Proportions : hybrides $= \dfrac{2}{8} = \dfrac{1}{4} = 25\,\%$ ; légères $= \dfrac{6}{8} = \dfrac{3}{4} = 75\,\%$.

Aspect du tube : \textbf{deux bandes}, l'une en position intermédiaire (hybrides, 25 \% de l'ADN, bande moins intense) et l'autre en position haute (légères, 75 \% de l'ADN, bande trois fois plus intense). Aucune bande lourde, car il n'existe plus de molécule $^{15}$N/$^{15}$N dès la première génération sur milieu $^{14}$N.
\end{exoresolu}

\courssub{Résultat, fidélité et lien avec la division cellulaire}

\textbf{Résultat.} À l'issue de la réplication, on obtient \textbf{deux molécules d'ADN identiques} entre elles et identiques à la molécule mère, chacune formée d'un brin parental et d'un brin néoformé. Dans un chromosome, ces deux molécules constituent les deux \cle{chromatides sœurs}, réunies au niveau du centromère jusqu'à la division.

\textbf{Fidélité.} La complémentarité des bases est le premier garant de l'exactitude : A ne s'apparie correctement qu'avec T, G qu'avec C. S'y ajoutent des \cle{mécanismes de correction} : l'ADN polymérase « relit » le brin qu'elle vient de synthétiser et élimine la plupart des nucléotides mal appariés (notion admise à ton niveau). Le taux d'erreur final est d'environ une erreur pour $10^{9}$ nucléotides incorporés. Les rares erreurs échappant à la correction sont des \emph{mutations}, source de diversité génétique mais aussi de certaines maladies.

\begin{savaistu}[Des machines à copier ultra-rapides]
L'ADN polymérase de la bactérie \emph{E. coli} incorpore jusqu'à \num{1000} nucléotides par seconde : son génome de $4{,}6$ millions de paires de bases est répliqué en moins de 40 minutes. Chez l'Homme, l'ADN polymérase est plus lente (environ 50 nucléotides par seconde), mais la réplication démarre simultanément en des milliers de points d'origine : les $3 \times 10^{9}$ paires de bases du génome humain sont ainsi copiées en quelques heures, le temps de la phase S. Un record d'organisation industrielle... dans chacune de tes cellules !
\end{savaistu}

\textbf{Lien avec la division cellulaire.} La réplication est le prélude indispensable de toute division : elle a lieu en phase S de l'interphase, \emph{avant} la mitose comme \emph{avant} la méiose. Grâce à elle, la quantité d'ADN passe de $Q$ à $2Q$, ce qui permettra à chaque cellule fille de recevoir un stock complet $Q$ (mitose) ou de préparer la réduction à $Q/2$ puis la fécondation (méiose). Sans réplication préalable, la division priverait les cellules filles d'une partie du patrimoine génétique : la réplication semi-conservative est donc le mécanisme qui garantit la \emph{continuité de l'information génétique} de génération en génération de cellules.

\begin{aretenir}[L'essentiel de la section]
\begin{itemize}
\item La réplication a lieu en \textbf{phase S de l'interphase}, avant toute division (mitose ou méiose) ; elle fait passer la quantité d'ADN de $Q$ à $2Q$.
\item L'expérience de \textbf{Meselson et Stahl} (marquage $^{15}$N/$^{14}$N, gradient de CsCl) réfute les modèles conservatif et dispersif et démontre que la réplication est \textbf{semi-conservative} : chaque molécule fille = 1 brin parental + 1 brin néoformé.
\item Acteurs : \textbf{hélicase} (ouverture de la fourche), \textbf{amorces}, \textbf{ADN polymérase} (synthèse dans le sens 5' $\rightarrow$ 3' par complémentarité A--T, G--C).
\item Après $n$ cycles, une molécule donne $2^{n}$ molécules, dont seulement 2 conservent un brin parental.
\item La fidélité repose sur la complémentarité des bases et sur des mécanismes de correction ; les erreurs résiduelles sont des mutations.
\end{itemize}
\end{aretenir}

\coursec{III. La transcription : de l'ADN à l'ARN messager}

\courssub{Pourquoi un intermédiaire ? Observation et problème}

À l'issue de la réplication, nous avons vu que l'ADN se duplique fidèlement dans le noyau pour transmettre l'information génétique aux cellules filles. Pourtant, l'ADN ne quitte jamais le noyau chez les eucaryotes. Or, les ribosomes, usines à protéines, sont situés dans le cytoplasme (libres ou fixés sur le réticulum endoplasmique). Comment l'information codée dans les gènes parvient-elle jusqu'aux ribosomes ?

\begin{experience}[Preuve de l'existence d'un intermédiaire mobile : l'expérience de Brenner, Jacob et Meselson (1961)]
\textbf{Observation :} On fait croître des bactéries \emph{E. coli} dans un milieu riche en azote lourd (\ce{^{15}N}) pour marquer tout leur ADN. On les transfère ensuite dans un milieu à azote léger (\ce{^{14}N}) en présence d'un virus (bactériophage) qui détourne la machinerie cellulaire pour synthétiser ses propres protéines.
\textbf{Protocole :} À différents instants après l'infection, on isole les ribosomes et on mesure leur densité par centrifugation sur gradient de saccharose.
\textbf{Résultat :} Dès les premières minutes, de nouveaux ribosomes « légers » (\ce{^{14}N}) apparaissent et synthétisent des protéines virales. L'ADN viral, lui, reste « lourd » (\ce{^{15}N}) dans le nucléoïde bactérien.
\textbf{Interprétation :} L'information génétique a été copiée sur une molécule instable, de courte durée de vie, synthétisée à partir de l'ADN viral et migrante vers les ribosomes : l'\textbf{ARN messager (ARNm)}.
\textbf{Conclusion :} L'ARNm est l'intermédiaire mobile qui transporte l'information du noyau (ADN) vers le cytoplasme (ribosomes).
\end{experience}

\begin{definition}[Transcription]
La \textbf{transcription} est la synthèse d'une molécule d'ARN messager (ARNm) à partir d'un brin d'ADN (le \textbf{brin transcrit} ou brin matrice) servant de modèle, selon les règles de complémentarité des bases. Elle a lieu dans le \textbf{noyau} chez les eucaryotes (dans le cytoplasme chez les procaryotes).
\end{definition}

\courssub{Acteurs et localisation}

\begin{center}
\begin{tabularx}{\linewidth}{|l|X|X|}\hline
\rowcolor{popblueL}
\textbf{Acteur} & \textbf{Rôle} & \textbf{Localisation (Eucaryote)} \\ \hline
\textbf{Gène (ADN)} & Support de l'information ; seul le \textbf{brin transcrit} (3' $\to$ 5') est lu. & Noyau (chromatine décondensée) \\ \hline
\textbf{ARN polymérase} & Enzyme qui catalyse l'assemblage des ribonucléotides. Elle reconnaît le \textbf{promoteur} (séquence de départ). & Noyau \\ \hline
\textbf{Ribonucléotides libres} (ATP, UTP, CTP, GTP) & Substrats ; apportent les bases A, U, C, G. L'énergie est fournie par l'hydrolyse des deux groupements phosphates terminaux (pyrophosphate libéré). & Nucléoplasme \\ \hline
\textbf{Facteurs de transcription} & Protéines aidant l'ARN polymérase à se fixer au promoteur et à ouvrir la double hélice. & Noyau \\ \hline
\end{tabularx}
\end{center}

\courssub{Étapes de la transcription (Démarche d'investigation)}

L'observation au microscope électronique de gènes en cours de transcription (gènes des ARN ribosomiques d'amphibiens, dits « arbres de Noël ») montre des brins d'ADN décorés d'ARN naissants de longueur croissante. Cela suggère un processus séquentiel.

\begin{methode}[Étapes de la transcription]
\begin{enumerate}
\item \textbf{Initiation (Fixation) :} L'ARN polymérase (aidée de facteurs de transcription) se fixe spécifiquement sur le \textbf{promoteur}, une séquence régulatrice située en amont (5') du gène. La double hélice s'ouvre localement sur une dizaine de paires de bases : c'est la \textbf{bulle de transcription}.
\item \textbf{Élongation :} L'ARN polymérase se déplace le long du brin transcrit dans le sens \textbf{3' $\to$ 5'}. Elle assemble les ribonucléotides complémentaires (A$\leftrightarrow$U, T$\leftrightarrow$A, C$\leftrightarrow$G, G$\leftrightarrow$C) en catalysant la formation de liaisons phosphodiester \textbf{5' $\to$ 3'} sur l'ARN naissant. L'ADN se referme derrière l'enzyme.
\item \textbf{Termination :} L'ARN polymérase atteint une séquence \textbf{terminateur}. L'ARNm complet est libéré, l'ARN polymérase se détache, la double hélice d'ADN se reforme intégralement.
\end{enumerate}
\end{methode}

\begin{popfigure}
\begin{tikzpicture}[scale=0.9, every node/.style={font=\small}]
% Couleurs
\colorlet{dnaColor}{popdark}
\colorlet{rnaColor}{poporange}
\colorlet{enzymeColor}{popblue}
\colorlet{baseColor}{popgreen}

% Brin transcrit (matrice) 3' -> 5' (dessiné de gauche à droite pour la lecture)
\draw[thick, dnaColor] (0,0) -- (12,0) node[right] {Brins transcrit (matrice) \textbf{3' $\to$ 5'}};
\foreach \x/\base in {1/T, 2/A, 3/C, 4/G, 5/G, 6/A, 7/T, 8/T, 9/C} {
    \draw[thick, dnaColor] (\x,0) -- (\x,0.5);
    \node[dnaColor] at (\x,0.7) {\base};
}

% Brin non transcrit (codant) 5' -> 3'
\draw[thick, dnaColor] (0,-2) -- (12,-2) node[right] {Brins non transcrit (codant) \textbf{5' $\to$ 3'}};
\foreach \x/\base in {1/A, 2/T, 3/G, 4/C, 5/C, 6/T, 7/A, 8/A, 9/G} {
    \draw[thick, dnaColor] (\x,-2) -- (\x,-1.5);
    \node[dnaColor] at (\x,-1.3) {\base};
}

% Liaisons H (partiellement ouvertes)
\foreach \x in {1,2,3,9} {
    \draw[dashed, popmuted] (\x,0) -- (\x,-2);
}

% Bulle de transcription (zone ouverte)
\draw[popline, thick, decorate, decoration={brace, amplitude=5pt}] (3.5,-2.3) -- (8.5,-2.3) node[midway, below=6pt] {Bulle de transcription};

% ARN Polymérase
\draw[enzymeColor, thick, rounded corners=5pt, fill=popblueL] (4.5,1.5) rectangle (8,3) node[midway, align=center, enzymeColor] {\textbf{ARN Polymérase}\\(avec facteurs\\de transcription)};
\draw[->, enzymeColor, thick] (6.25,1.5) -- (6.25,0.5);

% ARNm naissant (5' -> 3') synthétisé sous l'enzyme
% Séquence complémentaire du brin transcrit : AUG CCU AAG (pour TAC GGA TTC)
\draw[thick, rnaColor] (4.5,3.5) -- (8.5,3.5) node[right] {ARNm naissant \textbf{5' $\to$ 3'}};
\foreach \x/\base in {5/A, 6/U, 7/G, 7.5/C, 8/C, 8.5/U, 9/A, 9.5/A, 10/G} {
    \draw[thick, rnaColor] (\x,3.5) -- (\x,3);
    \node[rnaColor] at (\x,2.8) {\base};
}
% Liaisons H temporaires ARNm/ADN (hybride)
\foreach \x in {5,6,7,7.5,8,8.5,9,9.5} {
    \draw[dotted, popmuted] (\x,0) -- (\x,3);
}

% Flèche de déplacement
\draw[->, thick, poporange] (9,3.8) -- (10.5,3.8) node[above] {Déplacement 3' $\to$ 5' sur matrice};

% Légende bases
\node[align=left, anchor=north west] at (0.5,4) {
\textbf{Complémentarité :}\\
A (ADN) $\leftrightarrow$ U (ARN)\\
T (ADN) $\leftrightarrow$ A (ARN)\\
C (ADN) $\leftrightarrow$ G (ARN)\\
G (ADN) $\leftrightarrow$ C (ARN)
};
\end{tikzpicture}
\legende{Schéma de la transcription : bulle de transcription, ARN polymérase lisant le brin transcrit (3' $\to$ 5') et synthétisant l'ARNm (5' $\to$ 3') complémentaire. Le brin non transcrit a la même séquence que l'ARNm (sauf U/T).}
\end{popfigure}

\courssub{Produit : l'ARN messager (ARNm)}

L'ARNm est une molécule \textbf{monobrin}, généralement plus courte que l'ADN chromosomique (quelques centaines à milliers de nucléotides). Sa séquence est la \textbf{copie conforme du brin non transcrit} (brin codant), à la différence près que la \textbf{thymine (T) est remplacée par l'uracile (U)}.

\begin{propriete}[Relation entre les trois brins]
Soit un gène à deux brins :
\begin{itemize}
\item Brin transcrit (matrice) : lu en \textbf{3' $\to$ 5'}.
\item Brin non transcrit (codant) : \textbf{5' $\to$ 3'}.
\item ARNm produit : \textbf{5' $\to$ 3'}, complémentaire du brin transcrit, \textbf{identique au brin non transcrit (T $\leftrightarrow$ U)}.
\end{itemize}
\end{propriete}

\begin{methode}[Écrire la séquence de l'ARNm à partir du brin transcrit (matrice)]
\begin{enumerate}
\item Identifier le sens de lecture du brin transcrit : \textbf{3' $\to$ 5'}.
\item Appliquer la complémentarité \textbf{ARN/ADN} :
\begin{center}
\begin{tabular}{c|c}
\hline
Base sur brin transcrit (ADN) & Base sur ARNm \\ \hline
\textbf{A} & \textbf{U} \\
\textbf{T} & \textbf{A} \\
\textbf{C} & \textbf{G} \\
\textbf{G} & \textbf{C} \\
\hline
\end{tabular}
\end{center}
\item Écrire l'ARNm dans le sens \textbf{5' $\to$ 3'} (sens de synthèse).
\end{enumerate}
\end{methode}

\begin{exemplebox}[Application : Détermination de la séquence d'ARNm]
\textbf{Donné :} Brin transcrit (matrice) lu en 3' $\to$ 5' : \texttt{3'-TAC GGA TTC-5'}
\textbf{Cherché :} Séquence de l'ARNm (5' $\to$ 3').

\textbf{Raisonnement :}
\begin{align*}
\text{Matrice (3' $\to$ 5')} &: \quad T \quad A \quad C \quad G \quad G \quad A \quad T \quad T \quad C \\
\text{Complémentarité} &: \quad \downarrow \quad \downarrow \quad \downarrow \quad \downarrow \quad \downarrow \quad \downarrow \quad \downarrow \quad \downarrow \quad \downarrow \\
\text{ARNm (5' $\to$ 3')} &: \quad A \quad U \quad G \quad C \quad C \quad U \quad A \quad A \quad G
\end{align*}
\textbf{Réponse :} \texttt{5'-AUG CCU AAG-3'}

\textit{Remarque :} Le brin non transcrit (codant) aurait la séquence \texttt{5'-ATG CCT AAG-3'} (T au lieu de U).
\end{exemplebox}

\begin{attention}[Pièges fréquents au Baccalauréat]
\begin{enumerate}
\item \textbf{Utiliser T au lieu de U dans l'ARNm.} L'ARN ne contient \textbf{jamais} de thymine. Toute séquence d'ARN écrite avec des T est \textbf{fausse} (0 point).
\item \textbf{Confondre brin transcrit et brin non transcrit.} Le brin transcrit est la \emph{matrice} (lu 3'$\to$5'). Le brin non transcrit a la \emph{même séquence} que l'ARNm (sauf U/T). Si l'énoncé donne le brin \textbf{codant} (non transcrit), recopiez-le simplement en remplaçant T par U. S'il donne le brin \textbf{transcrit}, appliquez la complémentarité.
\item \textbf{Inverser le sens 5'/3'.} L'ARNm est \textbf{toujours} écrit 5' $\to$ 3'. La polymérase lit la matrice 3' $\to$ 5'.
\end{enumerate}
\end{attention}

\courssub{Complément indispensable pour le Bac : La maturation de l'ARNm chez les eucaryotes}

\begin{savaistu}[Au-delà du programme officiel : L'épiçage (splicing)]
Chez les eucaryotes (donc chez l'Homme, le maïs, le cacaoyer, \emph{Trypanosoma}...), l'ARNm fraîchement transcrit (pré-ARNm ou ARNm précurseur) contient des séquences non codantes : les \textbf{introns}, intercalées entre les séquences codantes : les \textbf{exons}. Avant de sortir du noyau par les pores nucléaires, le pré-ARNm subit une \textbf{maturation} :
\begin{enumerate}
\item \textbf{Coiffage (5' cap) :} Ajout d'une guanosine triphosphate modifiée (7-méthyl-G) en 5', protectrice et nécessaire à la fixation sur le ribosome.
\item \textbf{Polyadénylation (queue poly-A) :} Ajout d'une queue d'une centaine d'adénines (AAAA...) en 3', stabilisatrice.
\item \textbf{Épiçage (splicing) :} Excision précise des introns et jonction des exons par le \textbf{spliceosome} (complexe ARN-protéines). \emph{Un même gène peut donner plusieurs ARNm matures différents (épissage alternatif) $\to$ diversité protéique.}
\end{enumerate}
Seul l'ARNm mature (exons uniquement, coiffé, polyadénylé) est exporté vers le cytoplasme.
\end{savaistu}

\begin{popfigure}
\begin{tikzpicture}[scale=0.85, every node/.style={font=\small}]
% Noyau
\draw[popblue, thick, fill=popblueL!20] (0,0) ellipse (4.5cm and 3.5cm);
\node[popblue, align=center] at (0,3) {\textbf{NOYAU}\\(Transcription + Maturation)};

% ADN (chromatine)
\draw[popdark, thick, decorate, decoration={snake, amplitude=1pt, segment length=4pt}] (-3,1.5) -- (-1,1.5) node[midway, above] {Gène (ADN)};
\draw[popdark, thick, decorate, decoration={snake, amplitude=1pt, segment length=4pt}] (-3,0.5) -- (-1,0.5);

% Transcription
\draw[poporange, thick, ->] (-1,1) -- (0,1.5) node[midway, above, sloped] {Transcription};
\node[poporange, align=center] at (0,2.2) {Pré-ARNm\\(Exons + Introns)};

% Maturation (Splicing)
\draw[poppurple, thick, ->] (0,1.5) -- (1.5,1.5);
\node[poppurple, align=center] at (1.5,2.2) {Maturation :\textbullet\ Coiffage 5'\textbullet\ Épiçage (retrait introns)\textbullet\ Poly-A 3'};

% ARNm mature
\node[popgreen, align=center, fill=popgreenL!30, rounded corners, draw=popgreen] at (3,1.5) {ARNm mature\\(Exons seulement)};

% Export
\draw[popgreen, thick, ->] (3.5,1.5) -- (4.5,1.5);
\draw[popblue, thick] (4.5,0.5) -- (4.5,2.5) -- (5.5,2.5) -- (5.5,0.5) -- cycle; % Pore nucléaire
\node[popblue, rotate=90] at (4.5,1.5) {Pore nucléaire};

% Cytoplasme
\draw[poporange, thick, fill=poporangeL!20] (6,0) ellipse (4cm and 3.5cm);
\node[poporange, align=center] at (6,3) {\textbf{CYTOPLASME}\\(Traduction)};

% Ribosome
\draw[popdark, thick, fill=popcream, rounded corners=3pt] (5.5,1) rectangle (7.5,2) node[midway, align=center] {Ribosome\\(Traduction)};
\draw[popgreen, thick, ->] (4.5,1.5) -- (5.5,1.5);

% Protéine naissante
\draw[poppink, thick, decorate, decoration={coil, aspect=0.3, segment length=3pt}] (7.5,1.5) -- (9,1.5) node[right] {Protéine};
\end{tikzpicture}
\legende{Localisation cellulaire : Transcription et maturation de l'ARNm dans le noyau (bleu), export par les pores nucléaires, traduction sur les ribosomes dans le cytoplasme (orange).}
\end{popfigure}

\courssub{Exercices résolus et entraînement}

\begin{exoresolu}[Détermination de brins et sens de lecture]
\textbf{Énoncé :} Un fragment d'ADN double brin a la séquence suivante (brin 1 donné en 5' $\to$ 3') :
\[ \text{Brin 1 : } \texttt{5'-ATG GCT TAA-3'} \]
\begin{enumerate}
\item Déduire la séquence du brin 2 (complémentaire).
\item On sait que le brin 1 est le \textbf{brin codant} (non transcrit). Écrire la séquence de l'ARNm produit (sens 5' $\to$ 3').
\item Si le brin 2 était le brin transcrit, quelle serait la séquence de l'ARNm ?
\end{enumerate}

\textbf{Solution détaillée :}
\begin{enumerate}
\item \textbf{Brins 2 (complémentaire, antiparallèle) :} Lecture du brin 1 en 5'$\to$3' impose brin 2 en 3'$\to$5'.
\begin{align*}
\text{Brin 1 (5' $\to$ 3')} &: \quad A \quad T \quad G \quad G \quad C \quad T \quad T \quad A \quad A \\
\text{Complémentarité ADN} &: \quad \downarrow \quad \downarrow \quad \downarrow \quad \downarrow \quad \downarrow \quad \downarrow \quad \downarrow \quad \downarrow \quad \downarrow \\
\text{Brin 2 (3' $\to$ 5')} &: \quad T \quad A \quad C \quad C \quad G \quad A \quad A \quad T \quad T \\
\text{Écriture conventionnelle (5' $\to$ 3')} &: \quad \texttt{5'-TTA AGC CAT-3'} \quad (\text{inverse de } TAC CGA ATT)
\end{align*}
\textit{Attention :} On écrit toujours les séquences en 5' $\to$ 3'. Brin 2 = \texttt{5'-TTA AGC CAT-3'}.

\item \textbf{Brins 1 = Brin codant (non transcrit).} L'ARNm a la même séquence (U au lieu de T).
\[ \text{ARNm (5' $\to$ 3')} = \texttt{5'-AUG GCU UAA-3'} \]
(Notez le codon STOP \texttt{UAA} en fin de séquence).

\item \textbf{Brins 2 = Brin transcrit (matrice).} Il est lu en 3' $\to$ 5' (soit \texttt{TAC CGA ATT}).
\begin{align*}
\text{Matrice (3' $\to$ 5')} &: \quad T \quad A \quad C \quad C \quad G \quad A \quad A \quad T \quad T \\
\text{ARNm (5' $\to$ 3')} &: \quad A \quad U \quad G \quad G \quad C \quad U \quad U \quad A \quad A \\
\text{ARNm} &= \texttt{5'-AUG GCU UAA-3'}
\end{align*}
\textbf{Conclusion :} Quel que soit le brin transcrit, l'ARNm est identique (copie du brin codant). C'est la définition du gène : l'information est redondante sur les deux brins.
\end{enumerate}
\end{exoresolu}

\begin{exercice}[Entraînement : Transcription et maturation]
\textbf{Énoncé :} On donne la séquence du brin transcrit (matrice) d'un gène eucaryote (lu en 3' $\to$ 5') :
\texttt{3'-TAC GTA AGG CTC ACT-5'}
Ce gène contient un intron situé entre le 6\textsuperscript{e} et le 7\textsuperscript{e} triplet de l'ARNm précurseur.
\begin{enumerate}
\item Écrire la séquence du pré-ARNm (5' $\to$ 3').
\item Écrire la séquence de l'ARNm mature après épiçage (on suppose que l'intron correspond exactement à la séquence codée par les triplets 7 à 9 de la matrice).
\item Quelle serait la conséquence d'une mutation dans le site d'épiçage 5' de l'intron (séquence consensus GU) ?
\end{enumerate}
\end{exercice}

\begin{corrige}[Exercice n]
\begin{enumerate}
\item \textbf{Pré-ARNm (copie complémentaire de la matrice) :}
Matrice : TAC GTA AGG CTC ACT
ARNm : \quad AUG CAU UCC GAG UGA
Séquence : \texttt{5'-AUG CAU UCC GAG UGA-3'}

\item \textbf{Maturation :} L'intron correspond aux triplets 7 à 9 de la matrice (\texttt{CTC ACT} $\to$ \texttt{GAG UGA} sur ARNm).
Les exons correspondent aux triplets 1-6 (\texttt{AUG CAU UCC}) et au reste (s'il y en a).
ARNm mature : \texttt{5'-AUG CAU UCC-3'} (suivi de la queue poly-A). \textit{Note : Ici, l'intron contient un codon STOP (UGA) ; son retrait est crucial pour produire une protéine fonctionnelle.}

\item \textbf{Mutation site d'épiçage 5' (GU $\to$ mutation) :} L'intron n'est pas excisé. L'ARNm mature conserve la séquence de l'intron. Cela décale le cadre de lecture (frameshift) ou introduit un codon STOP prématuré $\to$ protéine tronquée, non fonctionnelle $\to$ maladie génétique possible (ex : $\beta$-thalassémie au Cameroun liée à des mutations d'épiçage du gène $\beta$-globine).
\end{enumerate}
\end{corrige}

\begin{aretenir}[Bilan : La transcription]
\begin{itemize}
\item \textbf{Entrée :} Un gène (brin transcrit d'ADN) dans le noyau.
\item \textbf{Acteur clé :} ARN polymérase + facteurs de transcription.
\item \textbf{Sortie :} ARNm (copie complémentaire du brin transcrit, identique au brin codant sauf U/T).
\item \textbf{Règles :} Lecture 3'$\to$5' sur matrice, synthèse 5'$\to$3', complémentarité A-U, T-A, C-G.
\item \textbf{Spécificité eucaryote :} Maturation obligatoire (coiffage, poly-A, épiçage) avant export cytoplasmique.
\end{itemize}
\end{aretenir}

\coursec{IV. La traduction : de l'ARNm à la protéine}

Nous venons de voir que la transcription produit, dans le noyau, une molécule d'ARN messager (ARNm) qui porte la copie fidèle de l'information génétique contenue dans un gène. Mais une protéine n'est pas faite de nucléotides : elle est faite d'acides aminés. Comment la cellule passe-t-elle du « langage nucléique » (4 bases : A, U, C, G) au « langage protéique » (20 acides aminés) ? C'est précisément le rôle de la \cle{traduction}, véritable opération de conversion d'un texte écrit en lettres nucléiques vers un texte écrit en lettres aminées. Cette étape, qui se déroule dans le cytoplasme, achève l'expression du gène.

\courssub{Définition et localisation de la traduction}

\begin{definition}[Traduction]
La \cle{traduction} est la synthèse d'une chaîne \cle{polypeptidique} (enchaînement ordonné d'acides aminés reliés par des liaisons peptidiques) à partir de la séquence de nucléotides d'un \cle{ARN messager}. Elle constitue la deuxième étape de la biosynthèse des protéines, après la transcription.
\end{definition}

\begin{aretenir}[Localisation]
La traduction a lieu dans le \cle{cytoplasme}, au niveau des \cle{ribosomes}. On distingue :
\begin{itemize}
\item les \textbf{ribosomes libres}, dispersés dans le hyaloplasme : ils synthétisent les protéines destinées à rester dans la cellule (enzymes du cytoplasme, hémoglobine, etc.) ;
\item les \textbf{ribosomes fixés} sur le \textbf{réticulum endoplasmique granuleux (REG)} : ils synthétisent les protéines destinées à être exportées (hormones protéiques comme l'insuline, enzymes digestives), insérées dans les membranes ou stockées dans les lysosomes.
\end{itemize}
Chez les procaryotes (bactéries), qui n'ont pas de noyau, la traduction peut même commencer sur un ARNm encore en cours de transcription.
\end{aretenir}

\begin{experience}[Observation : où sont synthétisées les protéines ?]
\textbf{Protocole.} On injecte à un animal des acides aminés marqués au carbone radioactif ($^{14}$C). On prélève des cellules de son pancréas (glande qui sécrète de nombreuses enzymes protéiques) à différents temps et on localise la radioactivité par autoradiographie en microscopie électronique.

\textbf{Observations.}
\begin{center}
\begin{tabularx}{\linewidth}{|l|X|}\hline
\rowcolor{popblueL} \textbf{Temps après injection} & \textbf{Localisation de la radioactivité} \\ \hline
3 min & Réticulum endoplasmique granuleux (ribosomes fixés) \\ \hline
20 min & Appareil de Golgi \\ \hline
60 min & Vésicules de sécrétion, puis extérieur de la cellule \\ \hline
\end{tabularx}
\end{center}

\textbf{Interprétation.} Les acides aminés marqués sont d'abord incorporés au niveau des ribosomes du REG : c'est donc là que s'effectue l'assemblage des protéines. Celles-ci migrent ensuite vers l'appareil de Golgi (maturation, conditionnement) puis sont excrétées.

\textbf{Conclusion.} Les ribosomes sont bien le \emph{lieu} de la synthèse des protéines, et les acides aminés en sont les \emph{matériaux}.
\end{experience}

\courssub{Les acteurs de la traduction}

La traduction met en jeu cinq catégories d'acteurs, qu'il faut savoir citer et situer :

\begin{itemize}
\item l'\cle{ARNm} : il porte le \textbf{message} codé, lu par triplets de bases appelés \textbf{codons}, dans le sens $5' \rightarrow 3'$ ;
\item les \cle{ribosomes} : ce sont les « usines » d'assemblage. Chaque ribosome est constitué de deux sous-unités (une petite et une grande) formées d'\textbf{ARN ribosomiques (ARNr)} associés à des protéines. Le ribosome possède deux sites d'accueil des ARNt : le \textbf{site P} (qui porte la chaîne en construction) et le \textbf{site A} (qui accueille le nouvel ARNt chargé de son acide aminé) ;
\item les \cle{ARNt} (ARN de transfert) : ce sont les « adaptateurs ». Chaque ARNt transporte un \textbf{acide aminé} spécifique à son extrémité $3'$ et porte, sur une boucle opposée, un triplet de bases appelé \textbf{anticodon}, complémentaire du codon de l'ARNm ;
\item les \textbf{acides aminés} : les 20 acides aminés, matériaux de construction de la chaîne ;
\item les \textbf{enzymes} (notamment les aminoacyl-ARNt synthétases, qui accrochent chaque acide aminé au bon ARNt) et l'\textbf{énergie} fournie par l'\textbf{ATP} et le GTP.
\end{itemize}

\begin{definition}[Codon et anticodon]
Un \cle{codon} est un triplet de nucléotides de l'ARNm qui code un acide aminé (ou un signal d'arrêt). Un \cle{anticodon} est le triplet de nucléotides porté par un ARNt, complémentaire et antiparallèle d'un codon de l'ARNm. L'appariement codon--anticodon se fait selon les règles de complémentarité : A--U, U--A, C--G, G--C.
\end{definition}

\begin{exemplebox}[Appariement codon--anticodon]
Pour le codon $5'$-UUU-$3'$ de l'ARNm (phénylalanine), l'ARNt porte l'anticodon $3'$-AAA-$5'$ : les deux triplets sont \textbf{complémentaires} (U face à A) et \textbf{antiparallèles} (orientés en sens inverse), exactement comme les deux brins d'ADN. Cette reconnaissance base à base garantit que le \emph{bon} acide aminé est livré au \emph{bon} moment. De même, face au codon initiateur $5'$-AUG-$3'$ se place l'anticodon $3'$-UAC-$5'$.
\end{exemplebox}

\courssub{Les trois étapes de la traduction}

\textbf{a) L'initiation.}\par
La petite sous-unité du ribosome se fixe sur l'ARNm, puis la grande sous-unité s'assemble. Le premier codon lu est toujours le \textbf{codon initiateur} \cle{AUG}, qui code la \textbf{méthionine}. L'ARNt initiateur, porteur de la méthionine et muni de l'anticodon complémentaire UAC, se place sur le site P. Toute chaîne polypeptidique commence donc, au moment de sa synthèse, par une méthionine (souvent éliminée ensuite lors de la maturation).

\begin{definition}[Codon initiateur et codons stop]
Le \cle{codon initiateur} AUG marque le début de la traduction et code la méthionine. Les \cle{codons stop} (ou codons non-sens) UAA, UAG et UGA ne codent aucun acide aminé : ils signalent la fin de la traduction.
\end{definition}

\textbf{b) L'élongation.}\par
Le ribosome progresse le long de l'ARNm dans le sens $5' \rightarrow 3'$, codon par codon. À chaque étape :
\begin{enumerate}
\item un ARNt dont l'anticodon est complémentaire du codon exposé au \textbf{site A} s'y fixe, apportant son acide aminé ;
\item une \textbf{liaison peptidique} se forme (réaction catalysée par la grande sous-unité du ribosome) entre l'acide aminé nouvellement arrivé et l'extrémité de la chaîne polypeptidique naissante, qui est alors transférée sur l'ARNt du site A ;
\item le ribosome \textbf{transloque} d'un codon : l'ARNt libéré de son acide aminé quitte le ribosome, et l'ARNt porteur de la chaîne passe du site A au site P.
\end{enumerate}
La chaîne s'allonge ainsi d'un acide aminé par codon lu, toujours de l'extrémité N-terminale vers l'extrémité C-terminale.

\textbf{c) La terminaison.}\par
Lorsque le ribosome atteint un \textbf{codon stop} (UAA, UAG ou UGA), aucun ARNt ne peut s'y fixer. Des \textbf{facteurs de libération} se fixent à la place : la chaîne polypeptidique achevée est décrochée, le ribosome se dissocie en ses deux sous-unités et l'ARNm est libéré (puis dégradé ou réutilisé).

\textbf{d) Le produit.}\par
La chaîne polypeptidique libérée se \textbf{replie spontanément} dans l'espace (structure secondaire en hélice $\alpha$ ou feuillet $\beta$, puis structure tertiaire, parfois quaternaire) pour donner une \textbf{protéine fonctionnelle} : enzyme, hormone, anticorps, protéine de structure\ldots

\begin{popfigure}
\begin{center}
\begin{tikzpicture}[scale=0.9, every node/.style={font=\small}]
% ARNm
\draw[very thick, popblue] (0,0) -- (11,0);
\foreach \x/\c in {0.75/AUG, 2.25/UUU, 3.75/GGC, 5.25/AAA, 6.75/CCU, 8.25/UGA}{
  \draw[popblue, fill=popblueL, rounded corners=2pt] (\x-0.7,-0.35) rectangle (\x+0.7,0.35);
  \node[popdark] at (\x,0) {\c};
}
\node[popblue, left] at (-0.2,0) {$5'$};
\node[popblue, right] at (11.2,0) {$3'$};
\node[popblue, below=0.15cm] at (5.5,-0.5) {ARN messager};
% Ribosome
\draw[fill=poporangeL, draw=poporange, thick] (2.4,0.6) ellipse (1.9 and 0.8);
\draw[fill=poporange!60, draw=poporange, thick] (2.4,1.5) ellipse (1.5 and 0.65);
\node[popdark] at (2.4,1.5) {grande sous-unité};
\node[popdark, below=0.05cm] at (2.4,0.55) {petite sous-unité};
% ARNt 1 (site P)
\draw[thick, popgreen, fill=popgreenL] (1.5,0.4) -- (1.5,2.6) -- (1.9,2.9) -- (2.3,2.6) -- (2.3,0.4);
\node[popdark] at (1.9,0.15) {\scriptsize UAC};
\draw[fill=poppinkL, draw=poppink, thick] (1.9,3.3) circle (0.32);
\node at (1.9,3.3) {\scriptsize Met};
% ARNt 2 (site A)
\draw[thick, popgreen, fill=popgreenL] (3.1,0.4) -- (3.1,2.6) -- (3.5,2.9) -- (3.9,2.6) -- (3.9,0.4);
\node[popdark] at (3.5,0.15) {\scriptsize AAA};
\draw[fill=poppinkL, draw=poppink, thick] (3.5,3.3) circle (0.32);
\node at (3.5,3.3) {\scriptsize Phe};
% Chaîne polypeptidique naissante
\draw[very thick, poppurple] (1.9,3.62) .. controls (1.9,4.2) and (1.2,4.3) .. (0.9,4.7);
\draw[fill=poppurple!40, draw=poppurple] (0.7,4.9) circle (0.28);
\draw[fill=poppurple!40, draw=poppurple] (1.25,5.15) circle (0.28);
\draw[fill=poppurple!40, draw=poppurple] (1.8,5.25) circle (0.28);
\draw[very thick, poppurple] (0.98,4.98) -- (1.0,5.05);
\draw[very thick, poppurple] (1.53,5.2) -- (1.55,5.22);
\node[poppurple, align=left] at (0.4,5.6) {\scriptsize chaîne polypeptidique};
\node[poppurple, align=left] at (0.4,5.25) {\scriptsize naissante};
% Sens de lecture
\draw[-{Stealth[length=3mm]}, very thick, popgold] (4.6,2.6) -- (6.4,2.6);
\node[popgold!80!black, right] at (6.5,2.6) {sens de déplacement du ribosome $5' \rightarrow 3'$};
\end{tikzpicture}
\end{center}
\legende{Schéma d'un ribosome en cours de traduction. L'ARNm (bleu) est lu par codons (AUG, UUU, GGC\ldots). Les ARNt (verts) déposent leurs acides aminés (Met, Phe\ldots) grâce à l'appariement anticodon--codon ; la chaîne polypeptidique naissante (violette) s'allonge au fur et à mesure du déplacement du ribosome (orange) dans le sens $5' \rightarrow 3'$.}
\end{popfigure}

\courssub{Le polysome : une traduction en série}

\begin{experience}[Observation en microscopie électronique]
Sur des préparations de réticulocytes de lapin (cellules qui synthétisent massivement l'hémoglobine), on observe en microscopie électronique des « chapelets » : un même filament d'ARNm porte plusieurs ribosomes espacés, chacun tenant une chaîne polypeptidique de longueur croissante.

\textbf{Interprétation.} Dès qu'un ribosome s'est suffisamment éloigné du codon initiateur, un nouveau ribosome peut s'y fixer et commencer une nouvelle traduction. Plusieurs ribosomes traduisent donc \emph{simultanément} le même message. Le ribosome le plus avancé (le plus proche du codon stop) porte la chaîne la plus longue.

\textbf{Conclusion.} Cet ensemble ARNm + ribosomes multiples est appelé \cle{polysome} (ou polyribosome). Il permet une production rapide et abondante d'une même protéine à partir d'une seule molécule d'ARNm.
\end{experience}

\begin{popfigure}
\begin{center}
\begin{tikzpicture}[scale=0.95, every node/.style={font=\small}]
% ARNm
\draw[very thick, popblue] (0,0) .. controls (3,0.5) and (7,-0.5) .. (11,0);
\node[popblue, left] at (-0.1,0) {$5'$};
\node[popblue, right] at (11.2,0) {$3'$};
\node[popblue] at (5.5,-0.9) {ARNm};
% Ribosomes + chaînes de longueur croissante
\foreach \x/\n in {1.2/1, 4.0/3, 6.8/5, 9.4/7}{
  \draw[fill=poporangeL, draw=poporange, thick] (\x,0.75) circle (0.55);
  \foreach \i in {0,...,\n}{
    \draw[fill=poppurple!40, draw=poppurple] (\x-0.55+0.14*\i, 1.55+0.13*\i) circle (0.09);
  }
}
\node[popdark] at (1.2,0.75) {\scriptsize R};
\node[popdark] at (4.0,0.75) {\scriptsize R};
\node[popdark] at (6.8,0.75) {\scriptsize R};
\node[popdark] at (9.4,0.75) {\scriptsize R};
\draw[-{Stealth[length=3mm]}, very thick, popgold] (2.0,-1.5) -- (9.0,-1.5);
\node[popgold!80!black] at (5.5,-1.95) {sens de la traduction : chaînes de plus en plus longues};
\node[poppurple] at (10.6,2.6) {\scriptsize chaînes polypeptidiques};
\draw[-, poppurple] (9.75,2.15) -- (10.2,2.45);
\end{tikzpicture}
\end{center}
\legende{Schéma d'un polysome. Plusieurs ribosomes (R, orange) traduisent simultanément le même ARNm (bleu). Plus un ribosome est avancé vers l'extrémité $3'$, plus sa chaîne polypeptidique (violette) est longue.}
\end{popfigure}

\courssub{Déterminer une séquence peptidique : méthode et exemple}

\begin{methode}[De la séquence d'ARNm à la séquence polypeptidique]
\begin{enumerate}
\item Repérer le \textbf{codon initiateur AUG} : c'est à partir de lui que commence la lecture (le découpage en triplets est alors fixé).
\item Découper la suite de l'ARNm en \textbf{codons} (groupes de 3 bases), dans le sens $5' \rightarrow 3'$.
\item Pour chaque codon, chercher dans le \textbf{tableau du code génétique} l'acide aminé correspondant (1\textsuperscript{re} lettre à gauche, 2\textsuperscript{e} lettre en haut, 3\textsuperscript{e} lettre à droite).
\item S'arrêter au premier \textbf{codon stop} (UAA, UAG, UGA) : il ne code rien, il termine la chaîne.
\item Écrire la séquence peptidique : Met -- \ldots\ -- (dernier acide aminé avant le stop).
\end{enumerate}
\end{methode}

\begin{exoresolu}[Traduction d'un ARNm court]
\textbf{Énoncé.} On donne la séquence d'un ARNm : $5'$-AUG UUU GGC UAA-$3'$. Déterminer la chaîne polypeptidique synthétisée (on donne : AUG = Met ; UUU = Phe ; GGC = Gly ; UAA = stop).
\tcblower
\textbf{Solution détaillée.}
\begin{enumerate}
\item Le premier codon est AUG : codon initiateur $\rightarrow$ \textbf{Méthionine}.
\item Deuxième codon UUU $\rightarrow$ \textbf{Phénylalanine}.
\item Troisième codon GGC $\rightarrow$ \textbf{Glycine}.
\item Quatrième codon UAA : codon \textbf{stop} $\rightarrow$ arrêt de la traduction, aucun acide aminé ajouté.
\end{enumerate}
La chaîne polypeptidique obtenue est donc :
\[ \text{Méthionine} - \text{Phénylalanine} - \text{Glycine} \]
soit un tripeptide de 3 acides aminés, reliés par 2 liaisons peptidiques. Remarque : le codon stop n'est \emph{pas} traduit en acide aminé.
\end{exoresolu}

\begin{attention}[Erreur fréquente : lire le code génétique sur l'ADN]
Le tableau du code génétique s'utilise \textbf{uniquement avec les codons de l'ARNm}, écrits avec les bases A, U, C, G. Si l'énoncé donne une séquence d'\textbf{ADN} (bases A, T, C, G), il faut \emph{d'abord} transcrire :
\begin{itemize}
\item si c'est le \textbf{brin non transcrit} (brin codant) qui est donné, l'ARNm lui est identique, sauf que T est remplacé par U ;
\item si c'est le \textbf{brin transcrit} (brin matrice) qui est donné, l'ARNm est complémentaire : A$\rightarrow$U, T$\rightarrow$A, C$\rightarrow$G, G$\rightarrow$C.
\end{itemize}
Chercher directement un triplet d'ADN (par exemple ATG) dans le tableau du code est une faute : le tableau ne contient pas de T. Autre piège : oublier que le codon stop ne code \emph{aucun} acide aminé et compter un acide aminé de trop.
\end{attention}

\begin{savaistu}[Une protéine en moins d'une minute]
Chez la bactérie \emph{Escherichia coli}, un ribosome ajoute environ 15 à 20 acides aminés par seconde : une protéine moyenne de 300 acides aminés est assemblée en moins d'une minute ! Chez l'Homme, la vitesse est de 2 à 4 acides aminés par seconde. De plus, grâce aux polysomes, une cellule peut produire des centaines de copies d'une même protéine à partir d'un seul ARNm. Cette efficacité est mise à profit en biotechnologie : l'insuline utilisée pour traiter les diabétiques (y compris au Cameroun, où le diabète progresse) est produite industriellement par des bactéries ou des levures dont on a modifié le patrimoine génétique pour qu'elles traduisent le gène humain de l'insuline.
\end{savaistu}

\courssub{Bilan général : le dogme central de la biologie moléculaire}

Enchaînons les deux étapes étudiées dans cette leçon. L'information génétique, stockée dans l'ADN au niveau du \textbf{noyau}, est d'abord recopiée en ARNm (\textbf{transcription}), puis l'ARNm migre vers le \textbf{cytoplasme} où il est traduit en protéine (\textbf{traduction}). Ce flux unidirectionnel de l'information constitue le \cle{dogme central de la biologie moléculaire}, formulé par Francis Crick en 1958 :

\[ \text{ADN} \xrightarrow[\text{(noyau)}]{\text{transcription}} \text{ARNm} \xrightarrow[\text{(cytoplasme, ribosomes)}]{\text{traduction}} \text{Protéine} \]

\begin{popfigure}
\begin{center}
\begin{tikzpicture}[scale=0.9, every node/.style={font=\small}]
% Noyau
\draw[fill=popblueL!50, draw=popblue, thick] (2.2,0) ellipse (2.4 and 1.7);
\node[popblue] at (2.2,-1.35) {\textbf{Noyau}};
% ADN double hélice stylisée
\draw[very thick, popblue, decorate, decoration={coil, aspect=0.4, segment length=6pt, amplitude=4pt}] (0.9,0.9) -- (3.5,0.9);
\draw[very thick, popblue!60, decorate, decoration={coil, aspect=0.4, segment length=6pt, amplitude=4pt}] (0.9,0.55) -- (3.5,0.55);
\node[popdark] at (2.2,0.15) {\textbf{ADN} (gène)};
% Flèche transcription
\draw[-{Stealth[length=3mm]}, very thick, poporange] (4.9,0.7) -- (6.1,0.7);
\node[poporange!80!black, above] at (5.5,0.7) {transcription};
\node[poporange!80!black, below] at (5.5,0.7) {\scriptsize (noyau)};
% ARNm
\draw[very thick, popgreen] (6.3,0.7) -- (8.7,0.7);
\foreach \x in {6.7,7.3,7.9,8.5}{\draw[popgreen, fill=popgreenL] (\x-0.22,0.45) rectangle (\x+0.22,0.95);}
\node[popdark] at (7.5,0.1) {\textbf{ARNm}};
% Cytoplasme
\draw[fill=popcream, draw=popgold, thick, rounded corners=8pt] (5.6,-3.4) rectangle (12.4,-1.4);
\node[popgold!80!black] at (11.6,-3.05) {\textbf{Cytoplasme}};
% Flèche traduction
\draw[-{Stealth[length=3mm]}, very thick, poporange] (9.0,0.4) -- (10.2,-1.7);
\node[poporange!80!black, right, align=left] at (10.3,-0.7) {traduction\\ \scriptsize (ribosomes)};
% Ribosome + protéine
\draw[fill=poporangeL, draw=poporange, thick] (7.3,-2.4) circle (0.5);
\node[popdark] at (7.3,-2.4) {\scriptsize ribosome};
\draw[very thick, popgreen] (5.9,-2.75) -- (8.7,-2.75);
\foreach \i in {0,...,6}{
  \draw[fill=poppurple!40, draw=poppurple] (8.0+0.28*\i, -1.85+0.1*\i) circle (0.13);
}
\node[poppurple, right] at (9.9,-1.6) {\textbf{Protéine}};
\draw[-{Stealth[length=3mm]}, very thick, poppurple] (7.6,-1.95) -- (8.3,-1.75);
\end{tikzpicture}
\end{center}
\legende{Le dogme central de la biologie moléculaire. Dans le noyau, l'ADN est transcrit en ARNm ; dans le cytoplasme, l'ARNm est traduit par les ribosomes en une chaîne polypeptidique qui se replie en protéine fonctionnelle.}
\end{popfigure}

\begin{aretenir}[L'essentiel de la section]
\begin{itemize}
\item La \textbf{traduction} est la synthèse d'une chaîne polypeptidique à partir de l'ARNm ; elle a lieu dans le \textbf{cytoplasme}, sur les \textbf{ribosomes}.
\item Acteurs : ARNm (codons), ribosomes (ARNr + protéines), ARNt (anticodon + acide aminé), acides aminés, enzymes, ATP/GTP.
\item Trois étapes : \textbf{initiation} (codon AUG = méthionine), \textbf{élongation} (appariement codon--anticodon, liaisons peptidiques), \textbf{terminaison} (codons stop UAA, UAG, UGA).
\item Un \textbf{polysome} = plusieurs ribosomes traduisant simultanément le même ARNm.
\item Dogme central : ADN $\rightarrow$ ARNm $\rightarrow$ protéine. Le tableau du code génétique se lit \emph{uniquement} avec des codons d'ARNm (A, U, C, G).
\end{itemize}
\end{aretenir}

\begin{exercice}[À toi de jouer]
On donne la séquence du brin transcrit d'un fragment d'ADN : $3'$-TAC AAA CCG ATT-$5'$.
\begin{enumerate}
\item Écrire la séquence de l'ARNm correspondant (sens $5' \rightarrow 3'$).
\item En utilisant le tableau du code génétique, déterminer la séquence polypeptidique synthétisée.
\item Combien de liaisons peptidiques sont formées ? Combien d'ARNt ont participé à la traduction ?
\end{enumerate}
\end{exercice}

\begin{corrige}[Exercice : traduction d'un fragment d'ADN]
\begin{enumerate}
\item L'ARNm est complémentaire du brin transcrit (A$\rightarrow$U, T$\rightarrow$A, C$\rightarrow$G, G$\rightarrow$C) et antiparallèle : $5'$-AUG UUU GGC UAA-$3'$.
\item Lecture codon par codon : AUG = Méthionine ; UUU = Phénylalanine ; GGC = Glycine ; UAA = stop. La chaîne est : \textbf{Met -- Phe -- Gly}.
\item Trois acides aminés sont reliés par \textbf{2 liaisons peptidiques} ; \textbf{3 ARNt} ont participé (un par acide aminé ; le codon stop ne fait intervenir aucun ARNt).
\end{enumerate}
\end{corrige}

\coursec{V. Le code génétique et la relation gène – protéine – caractère}

Dans la section précédente, nous avons suivi pas à pas la traduction : le ribosome lit l'ARNm codon par codon, et chaque codon commande l'arrivée d'un acide aminé précis grâce aux ARNt. Mais une question fondamentale reste en suspens : \emph{qui décide} que tel codon correspond à tel acide aminé ? Autrement dit, quel est le \cle{dictionnaire} qui permet de passer du langage des nucléotides (4 lettres : A, U, G, C) au langage des protéines (20 acides aminés) ? Ce dictionnaire, c'est le \cle{code génétique}. Et une fois ce dictionnaire connu, nous pourrons répondre à la grande question de cette leçon : comment la séquence d'un gène détermine-t-elle un caractère visible de l'individu, par exemple la forme de ses globules rouges ?

\courssub{Le code génétique : le dictionnaire de la cellule}

\textbf{Le problème posé.} Il n'existe que 4 nucléotides différents dans l'ARNm, mais 20 acides aminés à coder. Un nucléotide seul ne pourrait coder que 4 acides aminés ; deux nucléotides ($4^2 = 16$ combinaisons) restent insuffisants. En revanche, des groupes de \textbf{trois nucléotides} offrent $4^3 = 64$ combinaisons possibles : c'est largement assez pour 20 acides aminés. C'est ce raisonnement, confirmé par les expériences de Nirenberg, Khorana et Ochoa (Prix Nobel 1968), qui a établi que le code est lu par \cle{triplets} de nucléotides appelés \cle{codons}.

\begin{definition}[Code génétique]
Le \cle{code génétique} est l'ensemble des correspondances entre les 64 codons de l'ARNm (triplets de nucléotides) et les 20 acides aminés (ou les signaux de terminaison). Il constitue le dictionnaire utilisé par le ribosome pour traduire la séquence d'un ARNm en séquence d'acides aminés d'une protéine.
\end{definition}

Le tableau ci-dessous donne les 64 correspondances. Pour lire un codon : première base dans la colonne de gauche (lignes), deuxième base en haut (colonnes), troisième base à droite (sous-colonnes).

\begin{center}
\small
\rowcolors{2}{popblueL}{white}
\begin{tabular}{|c|c|c|c|c|c|}
\hline
\rowcolor{popblue!40} \textbf{1\iere{} base} & \textbf{U} & \textbf{C} & \textbf{A} & \textbf{G} & \textbf{3\ieme{} base} \\
\hline
\multirow{4}{*}{\textbf{U}} & UUU Phé & UCU Ser & UAU Tyr & UGU Cys & U \\
 & UUC Phé & UCC Ser & UAC Tyr & UGC Cys & C \\
 & UUA Leu & UCA Ser & UAA \textbf{Stop} & UGA \textbf{Stop} & A \\
 & UUG Leu & UCG Ser & UAG \textbf{Stop} & UGG Trp & G \\
\hline
\multirow{4}{*}{\textbf{C}} & CUU Leu & CCU Pro & CAU His & CGU Arg & U \\
 & CUC Leu & CCC Pro & CAC His & CGC Arg & C \\
 & CUA Leu & CCA Pro & CAA Gln & CGA Arg & A \\
 & CUG Leu & CCG Pro & CAG Gln & CGG Arg & G \\
\hline
\multirow{4}{*}{\textbf{A}} & AUU Ile & ACU Thr & AAU Asn & AGU Ser & U \\
 & AUC Ile & ACC Thr & AAC Asn & AGC Ser & C \\
 & AUA Ile & ACA Thr & AAA Lys & AGA Arg & A \\
 & AUG \textbf{Met} (init.) & ACG Thr & AAG Lys & AGG Arg & G \\
\hline
\multirow{4}{*}{\textbf{G}} & GUU Val & GCU Ala & GAU Asp & GGU Gly & U \\
 & GUC Val & GCC Ala & GAC Asp & GGC Gly & C \\
 & GUA Val & GCA Ala & GAA Glu & GGA Gly & A \\
 & GUG Val & GCG Ala & GAG Glu & GGG Gly & G \\
\hline
\end{tabular}
\end{center}

\begin{center}
\emph{Tableau du code génétique. Abréviations : Phé = phénylalanine, Leu = leucine, Ile = isoleucine, Met = méthionine, Val = valine, Ser = sérine, Pro = proline, Thr = thréonine, Ala = alanine, Tyr = tyrosine, His = histidine, Gln = glutamine, Asn = asparagine, Lys = lysine, Asp = aspartate, Glu = glutamate, Cys = cystéine, Trp = tryptophane, Arg = arginine, Gly = glycine. Stop = codon de terminaison.}
\end{center}

\textbf{Lecture du tableau.} Quelques repères essentiels :
\begin{itemize}
\item Le codon \textbf{AUG} code la méthionine : c'est le \cle{codon d'initiation}, le point de départ de toute traduction.
\item Trois codons — \textbf{UAA, UAG, UGA} — ne codent aucun acide aminé : ce sont les \cle{codons stop}, qui signalent la fin de la chaîne polypeptidique.
\item Un même acide aminé est souvent codé par plusieurs codons : la leucine par 6 codons (UUA, UUG, CUU, CUC, CUA, CUG), la sérine par 6, etc.
\end{itemize}

\begin{propriete}[Propriétés du code génétique (admises)]
Le code génétique possède quatre propriétés fondamentales :
\begin{enumerate}
\item Il est \cle{universel} (ou quasi universel) : le même codon code le même acide aminé chez la bactérie, la levure, le palmier et l'Homme. C'est un argument puissant en faveur de la parenté de tous les êtres vivants.
\item Il est \cle{dégénéré} : plusieurs codons différents peuvent coder le même acide aminé (64 codons pour 20 acides aminés). Conséquence majeure : certaines substitutions de bases sont \textbf{silencieuses}, c'est-à-dire sans effet sur la protéine.
\item Il est \cle{sans chevauchement} : chaque nucléotide n'appartient qu'à un seul codon.
\item Il est \cle{sans ponctuation} : les codons se succèdent en continu, sans séparateur ; la lecture commence au codon d'initiation AUG et se poursuit par triplets jusqu'à un codon stop.
\end{enumerate}
\end{propriete}

\begin{attention}[Toute mutation ne change pas forcément la protéine]
Erreur très fréquente au Baccalauréat : affirmer qu'une mutation (changement d'une base) modifie toujours la protéine. C'est \textbf{faux}. Grâce à la dégénérescence du code, une substitution peut être \emph{silencieuse} : par exemple, si le codon \textbf{GAA} (glutamate) devient \textbf{GAG}, la protéine reçoit toujours du glutamate et reste inchangée. Avant de conclure à un effet d'une mutation, il faut \emph{toujours} consulter le tableau du code génétique.
\end{attention}

\courssub{Le gène : unité d'information génétique}

\begin{definition}[Gène]
Un \cle{gène} est un segment d'ADN (une portion de la molécule d'ADN d'un chromosome) qui porte l'information nécessaire à la synthèse d'une chaîne polypeptidique déterminée. La séquence des nucléotides du gène, transcrite en ARNm puis traduite, détermine la séquence des acides aminés de la protéine.
\end{definition}

On peut donc résumer le « dogme central » de la biologie moléculaire par la chaîne d'information suivante :

\[ \text{ADN (gène)} \xrightarrow{\ \text{transcription}\ } \text{ARNm} \xrightarrow{\ \text{traduction}\ } \text{protéine} \xrightarrow{\ \text{action}\ } \text{caractère} \]

Reste à démontrer le dernier maillon : la protéine détermine le caractère. C'est exactement ce que permet d'établir l'étude d'une maladie génétique fréquente en Afrique centrale : la \cle{drépanocytose}.

\courssub{Étude d'un cas concret : la drépanocytose, ou comment une base change un caractère}

\textbf{Observation.} Dans certaines familles camerounaises, des enfants souffrent d'anémie chronique sévère, de crises douloureuses et d'infections répétées : c'est la drépanocytose (ou anémie falciforme). Au microscope, leurs globules rouges ne présentent pas la forme normale de disque biconcave, mais une forme allongée en \emph{faucille} (d'où le nom « falciforme »).

\begin{popfigure}
\begin{center}
\begin{tikzpicture}[scale=1]
% Globule rouge normal
\draw[fill=popink!60, draw=popdark] (0,0) ellipse (1.1 and 0.85);
\draw[fill=popink!25, draw=popdark] (0,0) ellipse (0.45 and 0.32);
\node[below] at (0,-1.1) {\small Globule rouge normal};
\node[below, popmuted] at (0,-1.55) {\small (disque biconcave)};
% Globule rouge falciforme
\begin{scope}[xshift=5.5cm]
\draw[fill=popink!60, draw=popdark, rotate=-20] plot[smooth cycle, tension=1.1] coordinates {(0,0.15) (0.7,0.55) (1.5,0.45) (1.9,0.05) (1.45,-0.15) (0.75,0.05)};
\node[below] at (1,-1.1) {\small Globule rouge falciforme};
\node[below, popmuted] at (1,-1.55) {\small (en faucille, drépanocytose)};
\end{scope}
\end{tikzpicture}
\legende{Comparaison schématique d'un globule rouge normal (disque biconcave, souple, traversant facilement les capillaires) et d'un globule rouge falciforme (rigide, en forme de faucille, bloquant les petits vaisseaux et détruit prématurément, d'où l'anémie).}
\end{center}
\end{popfigure}

\textbf{Hypothèse.} La forme anormale des globules rouges pourrait provenir d'une anomalie de leur protéine principale, l'\cle{hémoglobine}.

\textbf{Vérification expérimentale (données historiques).} En 1949, Linus Pauling montre par électrophorèse que l'hémoglobine des malades (notée \cle{HbS}) migre différemment de l'hémoglobine normale (\cle{HbA}) : les deux protéines n'ont pas la même charge électrique. En 1956, Vernon Ingram découpe les deux hémoglobines en peptides et compare leurs séquences : sur les 146 acides aminés de la chaîne $\beta$, \textbf{une seule différence} existe, en position 6 :

\begin{center}
\begin{tabularx}{\linewidth}{|l|X|X|}
\hline
\rowcolor{popblueL} & \textbf{Hémoglobine normale (HbA)} & \textbf{Hémoglobine drépanocytaire (HbS)} \\
\hline
Acide aminé n°6 de la chaîne $\beta$ & Glutamate (Glu), acide aminé chargé négativement & Valine (Val), acide aminé non chargé \\
\hline
Codon correspondant sur l'ARNm & \textbf{GAA} (ou GAG) & \textbf{GUA} (ou GUG) \\
\hline
Triplet sur le \textbf{brin codant} de l'ADN & \textbf{GAA} (ou GAG) & \textbf{GTA} (ou GTG) \\
\hline
\end{tabularx}
\end{center}

\textbf{Interprétation.} Le remplacement d'\textbf{une seule base} sur le brin codant de l'ADN (A remplacé par T, donc A remplacé par U sur l'ARNm) suffit à changer le codon GAA en GUA, donc à substituer la valine au glutamate. Cette unique substitution modifie les propriétés de l'hémoglobine : en conditions de faible oxygénation, les molécules d'HbS s'agrègent en longues fibres qui déforment le globule rouge en faucille.

\begin{popfigure}
\begin{center}
\begin{tikzpicture}[font=\small, node distance=0.35cm]
% HbA (Brin codant)
\node[draw=popblue, fill=popblueL, rounded corners, minimum width=5.2cm, minimum height=0.7cm] (adnA) {Brin codant ADN : \quad ... G \textbf{A} A ...};
\node[draw=popgreen, fill=popgreenL, rounded corners, minimum width=5.2cm, minimum height=0.7cm, below=of adnA] (arnA) {ARNm : \quad\quad\quad... G \textbf{A} A ...};
\node[draw=poporange, fill=poporangeL, rounded corners, minimum width=5.2cm, minimum height=0.7cm, below=of arnA] (protA) {Protéine : \quad... \textbf{Glutamate} ... (HbA)};
\draw[->, thick, popdark] (adnA) -- (arnA) node[midway, right, popmuted]{\scriptsize transcription (T$\rightarrow$U)};
\draw[->, thick, popdark] (arnA) -- (protA) node[midway, right, popmuted]{\scriptsize traduction};
\node[left=0.4cm of adnA, popblue, font=\small\bfseries] {HbA};
% Flèche mutation
\draw[->, very thick, popink] (3.0,-0.35) -- (4.0,-0.35) node[midway, above, popink]{\scriptsize mutation A $\rightarrow$ T (brin codant)};
% HbS
\begin{scope}[xshift=8.0cm]
\node[draw=popblue, fill=popblueL, rounded corners, minimum width=5.2cm, minimum height=0.7cm] (adnS) {Brin codant ADN : \quad ... G \textbf{T} A ...};
\node[draw=popgreen, fill=popgreenL, rounded corners, minimum width=5.2cm, minimum height=0.7cm, below=of adnS] (arnS) {ARNm : \quad\quad\quad... G \textbf{U} A ...};
\node[draw=poporange, fill=poporangeL, rounded corners, minimum width=5.2cm, minimum height=0.7cm, below=of arnS] (protS) {Protéine : \quad... \textbf{Valine} ... (HbS)};
\draw[->, thick, popdark] (adnS) -- (arnS) node[midway, right, popmuted]{\scriptsize transcription (T$\rightarrow$U)};
\draw[->, thick, popdark] (arnS) -- (protS) node[midway, right, popmuted]{\scriptsize traduction};
\node[left=0.4cm of adnS, popink, font=\small\bfseries] {HbS};
\end{scope}
\end{tikzpicture}
\legende{Comparaison des séquences brin codant ADN, ARNm et protéine pour le gène de la chaîne $\beta$ de l'hémoglobine. La substitution d'une seule base (A $\rightarrow$ T sur le brin codant, soit A $\rightarrow$ U sur l'ARNm) transforme le codon GAA (glutamate) en GUA (valine) : l'hémoglobine HbA devient HbS.}
\end{center}
\end{popfigure}

\textbf{Conclusion : la chaîne de causalité complète.}

\[ \underbrace{\text{mutation du gène } (A \rightarrow T \text{ sur brin codant})}_{\text{génotype}} \rightarrow \underbrace{\text{protéine modifiée (HbS)}}_{\text{molécule}} \rightarrow \underbrace{\text{globules rouges en faucille}}_{\text{cellule}} \rightarrow \underbrace{\text{anémie : caractère modifié}}_{\text{phénotype}} \]

\begin{aretenir}[Relation gène – protéine – caractère]
Le caractère d'un individu dépend des protéines qu'il fabrique (enzymes, protéines de structure, hémoglobine...). Or la séquence en acides aminés de chaque protéine est déterminée par la séquence en nucléotides du gène qui la code, via le code génétique. Toute modification de la séquence du gène (mutation) peut donc modifier la protéine, et par conséquent le caractère. On résume : \textbf{gène $\rightarrow$ protéine $\rightarrow$ caractère}.
\end{aretenir}

\begin{savaistu}[Drépanocytose et paludisme au Cameroun]
Au Cameroun, environ 20 à 25 \% de la population est porteuse du \emph{trait drépanocytaire} : ces personnes hétérozygotes (génotype AS, possédant un allèle normal HbA et un allèle muté HbS) sont en bonne santé et produisent assez d'hémoglobine normale. Or le parasite du paludisme, \emph{Plasmodium falciparum}, se développe mal dans leurs globules rouges : le trait AS confère une \textbf{protection partielle contre le paludisme grave}. C'est pourquoi l'allèle S, bien que dangereux à l'état homozygote (SS, drépanocytose sévère), a été maintenu par la sélection naturelle dans les régions où le paludisme sévit, comme toute la zone tropicale camerounaise. C'est un exemple célèbre d'\emph{avantage sélectif de l'hétérozygote}.
\end{savaistu}

\courssub{Les mutations : origine des variations des gènes}

\begin{definition}[Mutation]
Une \cle{mutation} est une modification de la séquence des nucléotides de l'ADN. On distingue :
\begin{itemize}
\item la \cle{substitution} : remplacement d'une base par une autre (ex. A $\rightarrow$ T dans le gène de l'hémoglobine) ;
\item l'\cle{insertion} : ajout d'une ou plusieurs bases dans la séquence ;
\item la \cle{délétion} : perte d'une ou plusieurs bases.
\end{itemize}
\end{definition}

\textbf{Effets sur la protéine.} Une substitution ne concerne qu'un seul codon : elle peut être \emph{silencieuse} (même acide aminé, grâce à la dégénérescence du code), \emph{faux-sens} (un acide aminé remplacé par un autre, comme dans la drépanocytose) ou \emph{non-sens} (création d'un codon stop qui raccourcit la protéine). En revanche, l'insertion ou la délétion d'une base (non multiple de 3) provoque un \cle{décalage du cadre de lecture} : tous les codons situés en aval sont modifiés, et la protéine est généralement profondément altérée.

\begin{exemplebox}[Illustration du décalage du cadre de lecture]
Séquence normale lue par triplets : AUG -- GAA -- CUU -- GGC ... $\rightarrow$ Met -- Glu -- Leu -- Gly ...\\
Après \textbf{délétion} de la 4\ieme{} base (le G de GAA) : AUG -- AAC -- UUG -- GC... $\rightarrow$ Met -- Asn -- Leu -- ... : à partir du site de la délétion, \emph{tous} les acides aminés sont changés. On comprend pourquoi ce type de mutation est souvent plus grave qu'une simple substitution.
\end{exemplebox}

\begin{methode}[Comparer deux séquences pour identifier une mutation et prédire son effet]
\begin{enumerate}
\item \textbf{Aligner} les deux séquences d'ADN (ou d'ARNm), nucléotide par nucléotide, à partir du début.
\item \textbf{Repérer la première différence} : une base remplacée (substitution), une base en plus (insertion) ou une base en moins (délétion).
\item \textbf{Découper en codons} à partir du codon d'initiation AUG et identifier le ou les codons affectés.
\item \textbf{Consulter le tableau du code génétique} pour traduire le codon normal et le codon muté.
\item \textbf{Conclure} : mutation silencieuse (même acide aminé), faux-sens (acide aminé différent), non-sens (codon stop) ou décalage du cadre de lecture (toute la suite modifiée).
\end{enumerate}
\end{methode}

\begin{exoresolu}[Analyse d'une mutation du gène de l'hémoglobine]
\textbf{Énoncé.} Chez un patient, on séquence un fragment du \textbf{brin codant} de l'ADN du gène de la chaîne $\beta$ de l'hémoglobine :
\begin{itemize}
\item Sujet sain : ... TGG -- ACT -- CCT -- \textbf{GAG} -- GAG -- AAG ...
\item Patient : \quad\; ... TGG -- ACT -- CCT -- \textbf{GTG} -- GAG -- AAG ...
\end{itemize}
(1) Identifier le type de mutation. (2) Déterminer son effet sur la protéine à l'aide du code génétique. (3) Quelle maladie ce patient présente-t-il ?
\tcblower
\textbf{Solution détaillée.}
\begin{enumerate}
\item \textbf{Type de mutation.} L'alignement des deux séquences montre une seule différence : dans le 4\ieme{} triplet, la base \textbf{A} (dans GAG) est remplacée par un \textbf{T} (dans GTG). Il s'agit d'une \cle{substitution} (mutation ponctuelle).
\item \textbf{Effet sur la protéine.} Le brin donné est le \textbf{brin codant} : l'ARNm a la \textbf{même séquence} (avec U à la place de T).
      \begin{itemize}
      \item Codon normal (ARNm) : \textbf{GAG} $\rightarrow$ \textbf{Glutamate} (Glu).
      \item Codon muté (ARNm) : \textbf{GUG} $\rightarrow$ \textbf{Valine} (Val).
      \end{itemize}
      La mutation est donc \emph{faux-sens} : le glutamate en position 6 de la chaîne $\beta$ est remplacé par la valine.
\item \textbf{Conclusion.} Cette substitution Glu $\rightarrow$ Val en position 6 de la chaîne $\beta$ est exactement la signature de l'hémoglobine \textbf{HbS} : le patient est atteint de \cle{drépanocytose} (s'il est homozygote SS) ou porteur du trait drépanocytaire (s'il est hétérozygote AS).
\end{enumerate}
\end{exoresolu}

\begin{attention}[Pièges de rédaction à éviter]
\begin{itemize}
\item Ne pas confondre \textbf{brin transcrit} (brin matrice, lu par l'ARN polymérase) et \textbf{brin codant} : l'ARNm est \emph{complémentaire} du brin transcrit et \emph{identique} (sauf T $\rightarrow$ U) au brin codant. Lisez toujours l'énoncé pour savoir quel brin est donné.
\item Ne pas écrire que « la mutation crée toujours une maladie » : beaucoup de mutations sont silencieuses ou sans conséquence phénotypique.
\item Ne pas confondre \textbf{mutation} (modification de la séquence d'ADN) et \textbf{modification du phénotype} : le lien entre les deux passe obligatoirement par la protéine et doit être démontré par l'analyse des séquences, pas affirmé gratuitement.
\end{itemize}
\end{attention}

\begin{exercice}[Pour s'entraîner]
On donne un fragment d'ARNm normal : AUG -- UUU -- GGC -- UAA et le même fragment chez un mutant : AUG -- UUU -- GGU -- CUA -- A...
\begin{enumerate}
\item Identifier le type de mutation survenue entre le 2\ieme{} et le 3\ieme{} codon.
\item Traduire les deux séquences à l'aide du tableau du code génétique.
\item Expliquer pourquoi cette mutation est qualifiée de « décalage du cadre de lecture » et pourquoi la protéine mutée est probablement non fonctionnelle.
\end{enumerate}
\end{exercice}

\begin{corrige}[Exercice : décalage du cadre de lecture]
\begin{enumerate}
\item Le début est identique (AUG -- UUU -- GG...), puis le mutant possède une base \textbf{U} insérée après GGU : il s'agit d'une \textbf{insertion} d'un nucléotide.
\item Séquence normale : AUG (Met) -- UUU (Phé) -- GGC (Gly) -- UAA (Stop) : peptide Met--Phé--Gly, terminé normalement. Séquence mutée : AUG (Met) -- UUU (Phé) -- GGU (Gly) -- CUA (Leu) -- A... : la lecture se poursuit au-delà de l'ancien codon stop, qui a disparu.
\item L'insertion d'une base décale le regroupement en triplets à partir du site de la mutation : tous les codons en aval sont changés, l'ancien codon stop UAA est détruit, et la traduction se prolonge anormalement. La protéine obtenue a une séquence très différente et une longueur anormale : elle est vraisemblablement non fonctionnelle.
\end{enumerate}
\end{corrige}

\begin{aretenir}[L'essentiel de la section]
\begin{itemize}
\item Le \cle{code génétique} fait correspondre chacun des 64 codons à un acide aminé (ou à un signal stop) ; il est universel, dégénéré, sans chevauchement ni ponctuation, et lu par triplets à partir du codon AUG.
\item Un \cle{gène} est un segment d'ADN portant l'information de synthèse d'une chaîne polypeptidique.
\item La drépanocytose illustre la chaîne \textbf{gène $\rightarrow$ protéine $\rightarrow$ caractère} : une substitution A $\rightarrow$ T (brin codant) change le codon GAA en GUA, le glutamate en valine, l'hémoglobine HbA en HbS, et les globules rouges deviennent falciformes, provoquant l'anémie.
\item Les mutations (substitution, insertion, délétion) modifient la séquence d'ADN ; leur effet dépend du code génétique : silencieuses, faux-sens, non-sens ou décalage du cadre de lecture.
\end{itemize}
\end{aretenir}

\coursec{VI. Variation de la quantité d'ADN au cours de la méiose}

Dans les sections précédentes, nous avons vu comment l'information portée par le gène s'exprime en une protéine, donc en un caractère. Mais cette information génétique, pour être utile à long terme, doit aussi être \emph{transmise} : d'une cellule mère à ses cellules filles lors de la multiplication cellulaire, et d'un parent à sa descendance lors de la reproduction sexuée. Or la reproduction sexuée pose un problème arithmétique redoutable : si deux cellules reproductrices à $2n$ chromosomes fusionnaient telles quelles, la cellule-œuf posséderait $4n$ chromosomes, et ce nombre doublerait à chaque génération ! La nature a résolu ce problème grâce à une division cellulaire particulière, la \cle{méiose}, qui divise par deux le stock chromosomique avant la formation des gamètes. Comment la quantité d'ADN évolue-t-elle au cours de cette double division ? C'est l'objet de cette dernière section, qui fait la synthèse de tout le chapitre : réplication de l'ADN, divisions cellulaires et transmission de l'information génétique.

\courssub{Rappels sur la méiose : une double division réductrice}

\begin{definition}[La méiose]
La \cle{méiose} est une succession de \textbf{deux divisions cellulaires} précédées d'une \textbf{seule réplication de l'ADN}. Elle permet, à partir d'une \cle{cellule mère} diploïde à $2n$ chromosomes, d'obtenir \textbf{quatre cellules filles haploïdes} à $n$ chromosomes : les cellules reproductrices ou \cle{gamètes} (spermatozoïdes et ovules chez l'Homme). Elle a lieu dans les gonades (testicules, ovaires).
\end{definition}

L'originalité de la méiose tient en trois points qu'il faut bien avoir en tête :

\begin{itemize}
\item \textbf{L'interphase} qui précède la méiose comporte une phase S au cours de laquelle l'ADN est \textbf{répliqué} (réplication semi-conservative, étudiée à la section II) : chaque chromosome, d'abord formé d'une seule chromatide, devient un chromosome à \textbf{deux chromatides sœurs} identiques, réunies par le centromère.
\item \textbf{La méiose I} (première division) sépare les \textbf{chromosomes homologues} : c'est la \cle{division réductionnelle}, car elle fait passer le nombre de chromosomes de $2n$ à $n$.
\item \textbf{La méiose II} (seconde division) sépare les \textbf{chromatides sœurs} de chaque chromosome : c'est la \cle{division équationnelle}, car le nombre de chromosomes ($n$) est conservé.
\end{itemize}

\begin{definition}[Division réductionnelle et division équationnelle]
La \cle{division réductionnelle} (méiose I) est la division au cours de laquelle les deux chromosomes de chaque paire d'homologues migrent vers des pôles opposés : chaque cellule fille reçoit $n$ chromosomes (au lieu de $2n$), chacun encore formé de \textbf{deux chromatides}. La \cle{division équationnelle} (méiose II) est la division au cours de laquelle les deux chromatides sœurs de chaque chromosome se séparent : le nombre de chromosomes reste $n$, mais chaque chromosome n'est plus formé que d'\textbf{une seule chromatide}. Elle est formellement comparable à une mitose.
\end{definition}

\begin{attention}[Réduction = réduction du nombre de chromosomes]
Le terme « réductionnelle » ne signifie pas que la quantité d'ADN est divisée par quatre ! La méiose I réduit le \textbf{nombre de chromosomes} de moitié ($2n \rightarrow n$). Comme chaque chromosome comporte encore deux chromatides, la quantité d'ADN n'est divisée que par deux. C'est l'ensemble des deux divisions qui ramène finalement la quantité d'ADN à la moitié de celle de la cellule mère non répliquée.
\end{attention}

\courssub{La notation Q : mesurer la quantité d'ADN par cellule}

Pour suivre objectivement les variations du stock d'ADN, les biologistes utilisent une unité commode.

\begin{definition}[La quantité Q d'ADN]
On appelle \cle{Q} la \textbf{quantité d'ADN contenue dans une cellule haploïde} ($n$ chromosomes à une chromatide). Une cellule diploïde à $2n$ chromosomes à une chromatide (avant réplication) contient donc $2Q$ d'ADN ; après réplication, chaque chromosome possédant deux chromatides, elle contient $4Q$.
\end{definition}

\begin{exemplebox}[Ordre de grandeur chez l'Homme]
Chez l'Homme, $n = 23$ chromosomes et $2n = 46$. La quantité $Q$ correspond à environ \SI{3,2}{pg} d'ADN (picogrammes), soit environ 3,2 milliards de paires de bases. Une cellule mère en interphase avant la phase S contient $2Q \approx \SI{6,4}{pg}$ ; après réplication, $4Q \approx \SI{12,8}{pg}$. Ces quantités sont mesurables expérimentalement par cytométrie de fluorescence après coloration spécifique de l'ADN.
\end{exemplebox}

\begin{attention}[Ne pas confondre Q et n !]
C'est LA confusion classique du baccalauréat. Le nombre $n$ compte les \textbf{chromosomes} (repérés à leur centromère) ; la quantité $Q$ mesure la \textbf{masse d'ADN}, c'est-à-dire le nombre de \textbf{chromatides}. Un chromosome à deux chromatides compte pour \textbf{un} chromosome mais pour \textbf{deux} quantités d'ADN. Ainsi, après la méiose I, chaque cellule possède $n$ chromosomes mais $2Q$ d'ADN.
\end{attention}

\courssub{Évolution de la quantité d'ADN : analyse des paliers et des chutes}

Suivons pas à pas le destin de l'ADN dans une cellule mère qui s'engage en méiose.

\textbf{Étape 1 — Interphase, avant la phase S : $2Q$.} La cellule mère diploïde possède $2n$ chromosomes à une chromatide, soit $2Q$ d'ADN.

\textbf{Étape 2 — Interphase, phase S : passage de $2Q$ à $4Q$.} La réplication semi-conservative double la quantité d'ADN : chaque molécule d'ADN donne deux molécules filles identiques. Le nombre de chromosomes est inchangé ($2n$), mais chacun est désormais formé de deux chromatides sœurs. La quantité d'ADN atteint $4Q$.

\textbf{Étape 3 — Fin de la méiose I : passage de $4Q$ à $2Q$.} La division réductionnelle sépare les chromosomes homologues et répartit le cytoplasme en deux cellules filles. Chacune reçoit $n$ chromosomes à deux chromatides, soit $2Q$ d'ADN : la moitié du stock de la cellule mère répliquée.

\textbf{Étape 4 — Fin de la méiose II : passage de $2Q$ à $Q$.} La division équationnelle sépare les chromatides sœurs. Chacune des quatre cellules filles (gamètes ou cellules destinées à le devenir) reçoit $n$ chromosomes à une chromatide, soit $Q$ d'ADN.

\begin{aretenir}[La séquence fondamentale]
Au cours de la méiose, la quantité d'ADN par cellule suit la séquence :
\[ 2Q \xrightarrow{\text{réplication (phase S)}} 4Q \xrightarrow{\text{méiose I (réductionnelle)}} 2Q \xrightarrow{\text{méiose II (équationnelle)}} Q \]
Un seul doublement (la réplication), deux chutes (les deux divisions). Au bilan : \textbf{une réplication pour deux divisions}, d'où quatre cellules à $Q$ à partir d'une cellule à $2Q$.
\end{aretenir}

\begin{methode}[Tracer la courbe de la quantité d'ADN en quatre points]
Pour construire sans erreur la courbe $Q_{\text{ADN}} = f(t)$ au cours de la méiose :
\begin{enumerate}
\item Repérer les \textbf{valeurs extrêmes} : la courbe part de $2Q$ (cellule mère avant réplication) et se termine à $Q$ (gamètes). L'axe des ordonnées est gradué en $Q$, $2Q$, $4Q$.
\item Placer le \textbf{doublement} : une montée progressive (la phase S dure plusieurs heures) de $2Q$ à $4Q$ pendant l'interphase. Attention : c'est une pente, pas un saut vertical, car la réplication s'étale dans le temps.
\item Placer la \textbf{première chute} : de $4Q$ à $2Q$, brutale, à la fin de la méiose I (la cytodiérèse, séparation du cytoplasme, est quasi instantanée à l'échelle du graphique).
\item Placer la \textbf{seconde chute} : de $2Q$ à $Q$, à la fin de la méiose II. Entre les deux chutes, un court palier à $2Q$ (interkinèse, sans réplication d'ADN !).
\end{enumerate}
Vérification finale : la courbe comporte exactement \textbf{une montée} et \textbf{deux descentes}.
\end{methode}

\begin{popfigure}
\begin{center}
\begin{tikzpicture}
\begin{axis}[
width=0.95\linewidth, height=7cm,
xlabel={Temps (étapes de la méiose)},
ylabel={Quantité d'ADN par cellule},
xmin=0, xmax=10, ymin=0, ymax=5,
xtick={1,3,5,7,9},
xticklabels={Interphase (G1),Phase S,Méiose I,Interkinèse,Méiose II},
xticklabel style={font=\small, align=center},
ytick={1,2,4},
yticklabels={$Q$,$2Q$,$4Q$},
grid=major, grid style={popline!40},
axis lines=left,
]
\addplot[popcurve, very thick, color=popblue] coordinates {
(0,2) (2,2) (3,4) (4,4) (5.5,4) (5.7,2) (7,2) (8,2) (8.2,1) (10,1)
};
\node[color=popblue, font=\small, align=center] at (axis cs:2.5,3.1) {Réplication\\de l'ADN};
\node[color=popdark, font=\small, align=center] at (axis cs:5.6,4.6) {Division\\réductionnelle};
\node[color=popdark, font=\small, align=center] at (axis cs:8.1,2.6) {Division\\équationnelle};
\node[color=popgreen, font=\small] at (axis cs:1,1.6) {cellule mère $2n$};
\node[color=popgreen, font=\small] at (axis cs:9.4,0.6) {4 gamètes $n$};
\end{axis}
\end{tikzpicture}
\end{center}
\legende{Courbe de variation de la quantité d'ADN par cellule au cours de la méiose. La montée de $2Q$ à $4Q$ correspond à la réplication (phase S de l'interphase) ; la première chute ($4Q \rightarrow 2Q$) à la division réductionnelle (méiose I) ; la seconde chute ($2Q \rightarrow Q$) à la division équationnelle (méiose II).}
\end{popfigure}

\begin{attention}[Erreur fréquente : faire chuter la courbe dès la réplication]
Beaucoup d'élèves font \emph{baisser} la quantité d'ADN au moment de la réplication, ou font coïncider la première chute avec la phase S. C'est doublement faux : la réplication \textbf{double} la quantité d'ADN ($2Q \rightarrow 4Q$) et \textbf{ne change pas le nombre de chromosomes} (toujours $2n$). Les seules chutes de la courbe correspondent aux \textbf{divisions cellulaires} (séparation effective du cytoplasme en deux cellules), jamais à un événement chromosomique interne à une cellule unique.
\end{attention}

\courssub{Schéma récapitulatif des deux divisions}

\begin{popfigure}
\begin{center}
\begin{tikzpicture}[
cell/.style={ellipse, draw=popdark, thick, minimum width=2.4cm, minimum height=1.7cm, align=center, font=\small},
fleche/.style={-{Stealth[length=3mm]}, very thick, poporange}
]
\node[cell, fill=popblueL, align=center] (mere) at (0,0){Cellule mère\\$2n$ chromosomes\\$2Q$ ADN};
\node[cell, fill=poppurpleL, align=center] (repl) at (4.2,0){Après réplication\\$2n$ chr. à 2 chromatides\\$4Q$ ADN};
\node[cell, fill=popgreenL, align=center] (f1) at (9,-1.6){Cellule fille\\$n$ chr. à 2 chromatides\\$2Q$ ADN};
\node[cell, fill=popgreenL, align=center] (f2) at (9,1.6){Cellule fille\\$n$ chr. à 2 chromatides\\$2Q$ ADN};
\node[cell, fill=popgoldL, align=center] (g1) at (13.4,2.6){Gamète\\$n$, $Q$};
\node[cell, fill=popgoldL, align=center] (g2) at (13.4,0.9){Gamète\\$n$, $Q$};
\node[cell, fill=popgoldL, align=center] (g3) at (13.4,-0.9){Gamète\\$n$, $Q$};
\node[cell, fill=popgoldL, align=center] (g4) at (13.4,-2.6){Gamète\\$n$, $Q$};
\draw[fleche] (mere) -- (repl) node[midway, above, font=\small\itshape, popdark]{phase S};
\draw[fleche] (repl.east) -- (f1.west);
\draw[fleche] (repl.east) -- (f2.west);
\node[font=\small\itshape, popdark, align=center] at (7.2,0) {Méiose I\\réductionnelle};
\draw[fleche] (f2.east) -- (g1.west);
\draw[fleche] (f2.east) -- (g2.west);
\draw[fleche] (f1.east) -- (g3.west);
\draw[fleche] (f1.east) -- (g4.west);
\node[font=\small\itshape, popdark, align=center] at (11.6,0) {Méiose II\\équationnelle};
\end{tikzpicture}
\end{center}
\legende{Schéma fonctionnel de la méiose avec les quantités d'ADN. Une cellule mère à $2Q$ se réplique en $4Q$, puis deux divisions successives produisent quatre cellules haploïdes à $Q$. Noter qu'après la méiose I, les cellules sont haploïdes en chromosomes ($n$) mais contiennent encore $2Q$ d'ADN (chromosomes à deux chromatides).}
\end{popfigure}

\begin{exemplebox}[Application chiffrée : la spermatogenèse humaine]
Chez l'Homme, une spermatogonie (cellule mère) possède $2n = 46$ chromosomes et $2Q$ d'ADN.
\begin{itemize}
\item Après la phase S : toujours 46 chromosomes, mais chacun à deux chromatides : $4Q$ d'ADN (spermatocyte de premier ordre, ou spermatocyte I).
\item Après la méiose I : deux spermatocytes de second ordre (spermatocytes II) à $n = 23$ chromosomes (à deux chromatides), soit $2Q$ d'ADN chacun.
\item Après la méiose II : quatre spermatides à $n = 23$ chromosomes (à une chromatide), soit $Q$ d'ADN chacune, qui se transformeront en spermatozoïdes (spermiogenèse).
\end{itemize}
À la fécondation, la fusion d'un spermatozoïde ($Q$) et d'un ovule ($Q$) restaure la quantité $2Q$ et le nombre diploïde $2n = 46$ : la constance du stock chromosomique d'une génération à l'autre est assurée.
\end{exemplebox}

\courssub{Le lien avec la réplication semi-conservative}

Le passage de $2Q$ à $4Q$ n'est autre que la \textbf{réplication de l'ADN} étudiée à la section II. Chacune des 46 molécules d'ADN de la cellule mère est dédoublée selon le mode semi-conservatif : chaque molécule fille conserve un brin parental et porte un brin néosynthétisé. Les deux molécules filles issues d'un même chromosome restent associées au niveau du centromère : ce sont les \textbf{deux chromatides sœurs}. Ainsi, le doublement de la quantité d'ADN s'accompagne d'une \textbf{fidélité totale de l'information génétique} : les quatre cellules issues de la méiose reçoivent des copies conformes de l'ADN parental (à l'exception des recombinaisons dues au crossing-over en prophase I et des mutations, sources de diversité génétique).

\begin{experience}[Observation expérimentale : la cytométrie de flux]
\textbf{Principe.} On colore l'ADN des cellules d'un testicule (organe où la méiose est permanente) par un fluorochrome dont la fluorescence est proportionnelle à la quantité d'ADN, puis on mesure la fluorescence cellule par cellule en cytométrie de flux.

\textbf{Observations.} L'histogramme fait apparaître trois pics principaux : un pic à $2Q$ (cellules en G1, avant réplication), un pic à $4Q$ (cellules ayant terminé la phase S, en méiose I) et un pic à $Q$ (spermatides et spermatozoïdes). Entre les pics $2Q$ et $4Q$, des cellules de fluorescence intermédiaire correspondent aux cellules en cours de phase S.

\textbf{Interprétation.} L'existence discrète de ces trois populations confirme la séquence $2Q \rightarrow 4Q \rightarrow 2Q \rightarrow Q$ et l'absence de réplication entre les deux divisions (aucune cellule ne repasse de $2Q$ à $4Q$ pendant l'interkinèse).

\textbf{Conclusion.} La courbe théorique de variation de la quantité d'ADN est validée expérimentalement : le doublement précède les deux divisions, et il n'y a qu'une seule phase S pour deux divisions.
\end{experience}

\courssub{Comparaison avec la mitose : ne pas confondre les deux courbes}

La mitose, division de multiplication cellulaire, conserve le stock chromosomique : une cellule mère à $2n$ donne deux cellules filles à $2n$. Sa courbe de quantité d'ADN est donc différente : $2Q \rightarrow 4Q \rightarrow 2Q$, avec \textbf{une seule chute}.

\begin{center}
\rowcolors{1}{popcream}{white}
\begin{tabularx}{\linewidth}{|l|X|X|}
\hline
\rowcolor{popblueL} \textbf{Critère} & \textbf{Mitose} & \textbf{Méiose} \\
\hline
Nombre de divisions & 1 & 2 \\
\hline
Nombre de réplications & 1 & 1 \\
\hline
Séquence des quantités d'ADN & $2Q \rightarrow 4Q \rightarrow 2Q$ & $2Q \rightarrow 4Q \rightarrow 2Q \rightarrow Q$ \\
\hline
Nombre de cellules filles & 2 & 4 \\
\hline
Formule chromosomique des cellules filles & $2n$ (diploïdes) & $n$ (haploïdes) \\
\hline
Rôle biologique & Croissance, renouvellement cellulaire & Production des gamètes \\
\hline
\end{tabularx}
\end{center}

\begin{popfigure}
\begin{center}
\begin{tikzpicture}
\begin{axis}[
width=0.95\linewidth, height=6.5cm,
xlabel={Temps}, ylabel={Quantité d'ADN par cellule},
xmin=0, xmax=10, ymin=0, ymax=5,
xtick=\empty,
ytick={1,2,4},
yticklabels={$Q$,$2Q$,$4Q$},
grid=major, grid style={popline!40},
axis lines=left,
legend style={at={(0.02,0.98)}, anchor=north west, font=\small},
]
\addplot[very thick, color=popblue] coordinates {
(0,2) (1.5,2) (2.5,4) (4.5,4) (4.7,2) (7,2) (7.2,1) (10,1)
};
\addlegendentry{Méiose : $2Q \rightarrow 4Q \rightarrow 2Q \rightarrow Q$}
\addplot[very thick, color=poporange, dashed] coordinates {
(0,2) (1.5,2) (2.5,4) (4.5,4) (4.7,2) (10,2)
};
\addlegendentry{Mitose : $2Q \rightarrow 4Q \rightarrow 2Q$}
\end{axis}
\end{tikzpicture}
\end{center}
\legende{Comparaison des courbes de quantité d'ADN au cours de la méiose (bleu) et de la mitose (orange). Les deux courbes partagent la même montée ($2Q \rightarrow 4Q$, réplication) et la même première chute (division), mais la méiose comporte une \textbf{seconde chute} ($2Q \rightarrow Q$) correspondant à la division équationnelle, absente de la mitose.}
\end{popfigure}

\begin{attention}[Piège classique du Bac]
Si l'énoncé montre une courbe avec \textbf{une seule chute} après le doublement, il s'agit d'une \textbf{mitose} (ou d'une division équationnelle isolée) ; s'il y a \textbf{deux chutes}, il s'agit d'une \textbf{méiose complète}. Autre piège : affirmer que « la méiose I divise la quantité d'ADN par quatre ». Faux : elle la divise par deux ($4Q \rightarrow 2Q$) ; c'est par rapport à la cellule mère \emph{répliquée} que le bilan final est un quart, mais par rapport à la cellule mère \emph{initiale} ($2Q$), le bilan est une simple réduction de moitié ($2Q \rightarrow Q$).
\end{attention}

\begin{exoresolu}[Exploitation d'une courbe expérimentale]
\textbf{Énoncé.} On mesure la quantité d'ADN par cellule dans les gonades d'un animal au cours du temps. On obtient la séquence suivante : palier à \SI{6,4}{pg}, montée progressive jusqu'à \SI{12,8}{pg}, palier, chute brutale à \SI{6,4}{pg}, court palier, nouvelle chute à \SI{3,2}{pg}, palier final. (1) Identifier le phénomène étudié et justifier. (2) Donner la valeur de $Q$ pour cette espèce. (3) Indiquer à quels événements correspondent la montée et chacune des deux chutes. (4) Préciser la formule chromosomique des cellules obtenues à la fin, sachant que la cellule mère possède $2n = 16$ chromosomes.
\tcblower
\textbf{Solution détaillée.}
\begin{enumerate}
\item La séquence comporte \textbf{un doublement} suivi de \textbf{deux chutes} : c'est la signature de la \textbf{méiose} (une mitose ne présenterait qu'une seule chute).
\item Le palier final correspond aux cellules haploïdes : $Q = \SI{3,2}{pg}$. La cellule mère contenait $2Q = \SI{6,4}{pg}$ et la cellule répliquée $4Q = \SI{12,8}{pg}$.
\item La montée de \SI{6,4}{pg} à \SI{12,8}{pg} correspond à la \textbf{réplication semi-conservative de l'ADN} pendant la phase S de l'interphase. La première chute (\SI{12,8}{pg} $\rightarrow$ \SI{6,4}{pg}) correspond à la \textbf{méiose I} (division réductionnelle : séparation des chromosomes homologues). La seconde chute (\SI{6,4}{pg} $\rightarrow$ \SI{3,2}{pg}) correspond à la \textbf{méiose II} (division équationnelle : séparation des chromatides sœurs).
\item Les quatre cellules filles sont haploïdes : $n = 8$ chromosomes à une chromatide, contenant chacune $Q = \SI{3,2}{pg}$ d'ADN.
\end{enumerate}
\end{exoresolu}

\begin{exercice}[À toi de jouer]
Chez le chimpanzé, $2n = 48$ chromosomes. (1) Donner la quantité d'ADN (en $Q$) et le nombre de chromosomes d'une cellule germinale mâle : a) avant la phase S ; b) à la fin de la phase S ; c) à la fin de la méiose I ; d) à la fin de la méiose II. (2) Tracer l'allure de la courbe de variation de la quantité d'ADN en fonction du temps. (3) Comparer avec la courbe obtenue lors d'une mitose d'une cellule de la peau du même animal.
\end{exercice}

\begin{corrige}[Exercice : le chimpanzé]
(1) a) Avant la phase S : $2n = 48$ chromosomes à une chromatide, $2Q$ d'ADN. b) Fin de la phase S : $48$ chromosomes à deux chromatides, $4Q$ d'ADN. c) Fin de la méiose I : $n = 24$ chromosomes à deux chromatides, $2Q$ d'ADN (deux cellules). d) Fin de la méiose II : $n = 24$ chromosomes à une chromatide, $Q$ d'ADN (quatre cellules). (2) La courbe part de $2Q$, monte progressivement à $4Q$ (phase S), chute à $2Q$ (fin méiose I), bref palier, puis chute à $Q$ (fin méiose II). (3) La mitose présente la même montée $2Q \rightarrow 4Q$ mais une seule chute, de $4Q$ à $2Q$ : les deux cellules filles restent diploïdes à $2n = 48$ chromosomes.
\end{corrige}

\begin{aretenir}[L'essentiel de la section]
\begin{itemize}
\item La méiose = \textbf{une réplication} ($2Q \rightarrow 4Q$) suivie de \textbf{deux divisions} : la méiose I, réductionnelle ($4Q \rightarrow 2Q$, séparation des homologues), et la méiose II, équationnelle ($2Q \rightarrow Q$, séparation des chromatides).
\item Une cellule mère à $2n$ chromosomes et $2Q$ d'ADN produit \textbf{quatre cellules haploïdes} à $n$ chromosomes et $Q$ d'ADN.
\item Après la méiose I, les cellules sont à $n$ chromosomes mais à $2Q$ d'ADN : ne jamais confondre nombre de chromosomes et quantité d'ADN.
\item La courbe de la méiose comporte \textbf{une montée et deux chutes} ; celle de la mitose, \textbf{une montée et une seule chute}.
\item La fécondation ($Q + Q$) restaure la diploïdie ($2n$, $2Q$), assurant la constance du patrimoine génétique au fil des générations.
\end{itemize}
\end{aretenir}

\begin{savaistu}[Et quand la méiose déraille...]
Des erreurs de séparation des chromosomes au cours de la méiose (les \textbf{non-disjonctions}) produisent des gamètes anormaux, porteurs d'un chromosome en plus ou en moins. Après fécondation, l'embryon présente alors une \textbf{anomalie chromosomique} : c'est l'origine de la trisomie 21 (trois exemplaires du chromosome 21, soit $2n+1 = 47$ chromosomes), la cause génétique la plus fréquente de handicap intellectuel, diagnostiquée notamment par caryotype dans les laboratoires d'analyses médicales, y compris au Cameroun. Comprendre la méiose, c'est donc aussi comprendre la base de la génétique médicale moderne.
\end{savaistu}

\newpage

\coursec{Activité d'intégration}

\coursec{Activité d'intégration : Enquête génétique sur la drépanocytose}

\courssub{Objectifs de l'activité}

À l'issue de cette enquête génétique, tu devras être capable de :
\begin{itemize}
\item comparer deux séquences d'ADN et identifier une mutation ponctuelle ;
\item transcrire une séquence d'ADN (brin matrice) en ARNm en respectant la polarité 3' $\rightarrow$ 5' du brin matrice ;
\item traduire un ARNm en chaîne polypeptidique à l'aide du code génétique ;
\item établir la chaîne de causalité \cle{gène} $\rightarrow$ \cle{protéine} $\rightarrow$ \cle{caractère} ;
\item interpréter la courbe de variation de la quantité d'ADN au cours de la méiose ;
\item formuler un conseil génétique simple, fondé sur des arguments scientifiques.
\end{itemize}

\courssub{Situation}

À l'hôpital régional de Bafoussam (région de l'Ouest), le couple M. et M\textsuperscript{me} Kamdem consulte le service de pédiatrie. Sur leurs quatre enfants, deux (le fils aîné et la fille cadette) présentent des crises douloureuses répétées, une anémie chronique et des infections fréquentes. Le diagnostic biologique est posé : il s'agit de la \cle{drépanocytose} (anémie falciforme), une maladie héréditaire dans laquelle les globules rouges prennent une forme de faucille en conditions d'hypoxie. Au Cameroun, environ \num{20} à \num{25}\,\% de la population est porteuse saine de l'allèle muté, et près de \num{2}\,\% des nouveau-nés sont atteints.

Le médecin généticien te confie le dossier de la famille. Ta mission : expliquer l'origine moléculaire de la maladie et conseiller les parents.

\textbf{Document 1 : arbre généalogique de la famille Kamdem.}

\begin{popfigure}
\begin{center}
\begin{tikzpicture}[scale=1.1, every node/.style={font=\small}]
% Génération I
\draw (0,2) rectangle (0.6,2.6);
\node at (0.3,1.6) {I-1};
\draw (1.4,2.3) circle (0.3);
\node at (1.4,1.6) {I-2};
% Lien union
\draw (0.6,2.3) -- (1.1,2.3);
% Lien descendance
\draw (0.85,2.3) -- (0.85,1.2);
\draw (-0.7,1.2) -- (2.4,1.2);
% Génération II : 4 enfants
\draw[fill=black] (-0.7,0.6) rectangle (-0.1,1.2);
\node at (-0.4,0.2) {II-1};
\draw (0.7,0.6) rectangle (1.3,1.2);
\node at (1.0,0.2) {II-2};
\draw[fill=black] (1.8,0.9) circle (0.3);
\node at (1.8,0.2) {II-3};
\draw (2.4,0.9) circle (0.3);
\node at (2.4,0.2) {II-4};
% Légende interne
\node[align=left, anchor=north west] at (2.8,2.6) {\textbf{Légende :}\\ $\square$ = Homme \quad $\blacksquare$ = Homme malade\\ $\bigcirc$ = Femme \quad $\bullet$ = Femme malade};
\end{tikzpicture}
\end{center}
\legende{Arbre généalogique de la famille Kamdem. Carré : homme ; cercle : femme ; symbole noir : individu drépanocytaire (homozygote $S/S$) ; symbole blanc : individu sain (phénotype normal). Les parents I-1 et I-2 sont sains, mais les enfants II-1 (garçon) et II-3 (fille) sont malades.}
\end{popfigure}

\textbf{Document 2 : séquences d'un fragment du gène de la $\beta$-globine (brin matrice / transcrit).}

On a séquencé le \emph{brin transcrit} (brin matrice, lu 3' $\rightarrow$ 5' par l'ARN polymérase) du gène codant la chaîne $\beta$ de l'hémoglobine, au niveau de la région correspondant aux acides aminés n\textsuperscript{o} 4 à 8 de la protéine (codons 4 à 8 de l'ARNm mature).

\begin{center}
\begin{tabularx}{\linewidth}{|l|X|}\hline
\rowcolor{popblueL} Individu & Séquence du brin matrice (5' $\rightarrow$ 3') \\ \hline
Sain (allèle $A$) & 5'-ACC ACT CTC CTC \textbf{GAG} AGT-3' \\ \hline
Malade (allèle $S$) & 5'-ACC ACT CTC CTC \textbf{GTG} AGT-3' \\ \hline
\end{tabularx}
\end{center}
\emph{Note :} La séquence codante (brin sens) complémentaire du brin matrice sain est 3'-TGG TGA GAG GAG \textbf{CTC} TCA-5', ce qui donne l'ARNm 5'-ACC ACU CUC CUC \textbf{GAG} GAG-3' (Thr-Thr-Leu-Leu-\textbf{Glu}-Glu). La mutation modifie le 5\textsuperscript{e} triplet du brin matrice présenté ici (position du codon 6 de la $\beta$-globine).

\textbf{Document 3 : extrait du code génétique (codons de l'ARNm).}

\begin{center}
\begin{tabular}{|c|c||c|c|}\hline
\rowcolor{popblueL} Codon (ARNm) & Acide aminé & Codon (ARNm) & Acide aminé \\ \hline
ACC & Thréonine (Thr) & GAG & Glutamate (Glu) \\ \hline
ACU & Thréonine (Thr) & GUG & Valine (Val) \\ \hline
CUC & Leucine (Leu) & UCC & Sérine (Ser) \\ \hline
CUG & Leucine (Leu) & UGA & Stop \\ \hline
\end{tabular}
\end{center}

\textbf{Document 4 : variation de la quantité d'ADN au cours de la spermatogenèse chez le père.}

\begin{popfigure}
\begin{center}
\begin{tikzpicture}
\begin{axis}[width=0.9\linewidth,height=5.5cm,
xlabel={Étapes de la spermatogenèse}, ylabel={Quantité d'ADN par cellule},
xtick={1,2,3,4,5}, xticklabels={$G_1$, $S$, $G_2$, Méiose I, Méiose II},
ytick={1,2,4}, yticklabels={$Q$,$2Q$,$4Q$},
ymin=0, ymax=4.5, xmin=0.5, xmax=5.5,
axis lines=left, thick, tick label style={font=\small}, label style={font=\small}]
\addplot[popcurve, very thick] coordinates {(1,2) (1.5,2) (2.5,4) (3,4) (3.5,2) (4,2) (4.5,1) (5.3,1)};
\addplot[dashed, popmuted] coordinates {(0.5,4) (5.5,4)};
\addplot[dashed, popmuted] coordinates {(0.5,2) (5.5,2)};
\addplot[dashed, popmuted] coordinates {(0.5,1) (5.5,1)};
\node[popdark, anchor=south] at (axis cs:1.5,2.2) {$\times 2$};
\node[popdark, anchor=north] at (axis cs:3.5,3) {$\div 2$};
\node[popdark, anchor=north] at (axis cs:4.5,1.5) {$\div 2$};
\end{axis}
\end{tikzpicture}
\end{center}
\legende{Variation de la quantité d'ADN par cellule au cours de la spermatogenèse. $Q$ : quantité d'ADN d'une cellule haploïde ($n$ chromosomes simples, 1 chromatide). Phase $S$ : réplication (synthèse d'ADN) ; Méiose I : division réductionnelle (séparation des chromosomes homologues) ; Méiose II : division équationnelle (séparation des chromatides sœurs).}
\end{popfigure}

\courssub{Consignes}

\begin{enumerate}
\item À l'aide du Doc 1, détermine si l'allèle responsable de la drépanocytose est dominant ou récessif. Justifie ta réponse en précisant le type de transmission (autosomique / gonosomique).
\item Compare les deux séquences du Doc 2 et identifie précisément la mutation (position dans le fragment, nature du changement de nucléotide sur le brin matrice).
\item Écris la séquence des deux ARNm (allèle $A$ et allèle $S$) obtenus par transcription du brin matrice. Montre l'étape intermédiaire (brin matrice lu 3' $\rightarrow$ 5').
\item Traduis les deux ARNm en chaînes polypeptidiques à l'aide du Doc 3 et identifie l'acide aminé substitué (numéro et nature).
\item Rédige la chaîne de causalité complète reliant la mutation au phénotype maladie (gène $\rightarrow$ ARNm $\rightarrow$ protéine $\rightarrow$ caractère).
\item À l'aide du Doc 4, indique à quelle étape la quantité d'ADN passe de $4Q$ à $2Q$, puis de $2Q$ à $Q$, et explique le mécanisme chromosomique correspondant.
\item Rédige un court texte (5 à 8 lignes) de conseil génétique à destination des parents : génotypes probables, risque pour une future grossesse, évocation du dépistage.
\end{enumerate}

\courssub{Ressources à mobiliser}

\begin{aretenir}[Outils de la leçon]
\begin{itemize}
\item \cle{Complémentarité des bases} : A--U (ARN) / A--T (ADN), G--C. L'ARNm est synthétisé \textbf{5' $\rightarrow$ 3'} sur le brin matrice lu \textbf{3' $\rightarrow$ 5'}. L'ARNm a la même séquence que le brin non transcrit (brin sens), l'uracile (U) remplaçant la thymine (T).
\item \cle{Code génétique} : chaque codon (triplet de nucléotides de l'ARNm) correspond à un acide aminé ; le code est universel, dégénéré et sans chevauchement. La lecture se fait 5' $\rightarrow$ 3' sans chevauchement.
\item \cle{Réplication semi-conservative} : en phase $S$, la quantité d'ADN double ($2Q \rightarrow 4Q$) ; chaque molécule fille conserve un brin parental.
\item \cle{Méiose} : la division réductionnelle (méiose I) sépare les chromosomes homologues ($4Q \rightarrow 2Q$, passage de $2n$ chromosomes à deux chromatides à $n$ chromosomes à deux chromatides) ; la division équationnelle (méiose II) sépare les chromatides sœurs ($2Q \rightarrow Q$, passage à $n$ chromosomes simples).
\item Transmission autosomique récessive : deux parents sains hétérozygotes $[A//S]$ ont, à chaque grossesse, \num{25}\,\% de risque d'avoir un enfant malade $[S//S]$, \num{50}\,\% porteur sain $[A//S]$, \num{25}\,\% indemne $[A//A]$.
\end{itemize}
\end{aretenir}

\begin{methode}[Transcription rigoureuse (ADN $\rightarrow$ ARNm)]
\begin{enumerate}
\item Écrire le brin matrice (donné 5' $\rightarrow$ 3') dans le sens de lecture par l'ARN polymérase : \textbf{3' $\rightarrow$ 5'} (inverser la séquence).
\item Appliquer la complémentarité (A$\rightarrow$U, T$\rightarrow$A, G$\rightarrow$C, C$\rightarrow$G) pour obtenir l'ARNm \textbf{5' $\rightarrow$ 3'}.
\item Vérifier : l'ARNm obtenu doit être identique au brin sens (non transcrit) 5' $\rightarrow$ 3', T remplacé par U.
\end{enumerate}
\end{methode}

\begin{methode}[Traduction (ARNm $\rightarrow$ Protéine)]
\begin{enumerate}
\item Repérer le codon d'initiation (AUG) si la séquence est complète, sinon commencer au premier codon de la séquence fournie.
\item Lire les codons par triplets non chevauchants dans le sens 5' $\rightarrow$ 3'.
\item Utiliser la table du code génétique pour écrire la séquence d'acides aminés (notation à trois lettres ou une lettre).
\item S'arrêter au codon Stop (UAA, UAG, UGA).
\end{enumerate}
\end{methode}

\courssub{Démarche et corrigé commenté}

\begin{exoresolu}[Enquête génétique : la drépanocytose dans la famille Kamdem]
Résoudre les sept tâches de l'enquête à partir des quatre documents du dossier.
\tcblower
\textbf{1. Mode de transmission (Doc 1).}
Les parents I-1 et I-2 sont phénotypiquement sains, mais ils ont engendré des enfants malades (II-1 et II-3). L'allèle responsable ($S$) peut donc être présent à l'état hétérozygote sans s'exprimer : il est \cle{récessif}.
La transmission touche les deux sexes (un garçon II-1 et une fille II-3 sont malades) : elle est \cle{autosomique} (non liée au chromosome X). Si l'allèle était récessif lié à l'X, le père d'une fille malade serait obligatoirement malade.
\textbf{Conclusion :} Transmission \textbf{autosomique récessive}. Les parents sont tous deux hétérozygotes \textbf{$[A//S]$} (porteurs sains).

\textbf{2. Identification de la mutation (Doc 2).}
Alignement des séquences du brin matrice (5' $\rightarrow$ 3') :
\begin{center}
\begin{tabular}{l l}
Sain ($A$) : & 5'-ACC ACT CTC CTC \textbf{GAG} AGT-3' \\
Malade ($S$) : & 5'-ACC ACT CTC CTC \textbf{GTG} AGT-3' \\
\end{tabular}
\end{center}
Le changement porte sur le \textbf{5\textsuperscript{e} triplet} du fragment (correspondant au codon 6 de la $\beta$-globine) : le nucléotide central \textbf{A} (adénine) est remplacé par un \textbf{T} (thymine) sur le brin matrice (\textbf{GAG $\rightarrow$ GTG}). C'est une \cle{mutation ponctuelle par substitution} (transition de type transversion purine $\rightarrow$ pyrimidine) d'une seule paire de bases. Le nombre de nucléotides est inchangé.

\textbf{3. Transcription (Doc 2 $\rightarrow$ ARNm).}
On applique la méthode : brin matrice lu 3' $\rightarrow$ 5' $\rightarrow$ ARNm 5' $\rightarrow$ 3'.
\begin{itemize}
\item \textbf{Brin matrice sain (5' $\rightarrow$ 3')} : ACC ACT CTC CTC GAG AGT
\item \textbf{Brin matrice sain lu 3' $\rightarrow$ 5'} : \textbf{TGA GAG GAG GAG TCA GGT} (Attention : complémentarité inverse. On inverse l'ordre : T G A G A G G A G G A G T C A G G T ? Non. Séquence 5'->3' : A C C A C T C T C C T C G A G A G T. Inversée 3'->5' : T G A G A G C T C C T C A G T C C A. Complémentaire pour ARNm : A C U C U C G A G G A G U C A G G U. Vérifions avec le brin sens.)
\end{itemize}
\emph{Calcul rigoureux :}
Brin matrice (5'$\rightarrow$3') : \texttt{ACC ACT CTC CTC GAG AGT}
Brin sens (3'$\rightarrow$5') : \texttt{TGG TGA GAG GAG CTC TCA}
Brin sens (5'$\rightarrow$3') : \texttt{ACT CAG GAG GAG TCA GGT} $\rightarrow$ ARNm (U/T) : \textbf{5'-ACU CAG GAG GAG UCA GGU-3'} ? Non, le fragment est court.
Reprenons sur le fragment centré sur la mutation (codons 4-8 $\beta$-globine).
Brin matrice sain 5' : \texttt{... CTC \textbf{GAG} AGT ...}
Brin matrice sain 3' : \texttt{... \textbf{CTC} GAG TCA ...} (Inverse + Complément : GAG -> CTC, AGT -> TCA).
ARNm 5' : \texttt{... \textbf{GAG} CUC AGU ...} (Complément de brin matrice 3' : CTC->GAG, GAG->CUC, TCA->AGU).
Codons ARNm : ... \textbf{GAG} (Glu) | CUC (Leu) | AGU (Ser) ...
Le codon 6 est bien \textbf{GAG} (Glu).

Brin matrice malade 5' : \texttt{... CTC \textbf{GTG} AGT ...}
Brin matrice malade 3' : \texttt{... \textbf{CAC} GAG TCA ...} (GTG -> CAC).
ARNm 5' : \texttt{... \textbf{GUG} CUC AGU ...} (CAC -> GUG).
Codons ARNm : ... \textbf{GUG} (Val) | CUC (Leu) | AGU (Ser) ...

\textbf{Résultat final pour le fragment fourni (6 codons) :}
\begin{center}
\begin{tabular}{l l}
\textbf{ARNm allèle $A$ (sain)} & 5'-\textbf{ACC ACU CUC CUC GAG GAG}-3' \\
\textbf{ARNm allèle $S$ (malade)} & 5'-\textbf{ACC ACU CUC CUC GUG GAG}-3' \\
\end{tabular}
\end{center}
(Les 4 premiers codons ACC, ACU, CUC, CUC codent Thr, Thr, Leu, Leu. Le 5\textsuperscript{e} codon change : GAG $\rightarrow$ GUG).

\textbf{4. Traduction (Doc 3).}
\begin{center}
\begin{tabular}{l c c c c c c}
Codon n\textsuperscript{o} & 1 & 2 & 3 & 4 & \textbf{5 (Codon 6 $\beta$)} & 6 \\
\hline
ARNm $A$ & ACC & ACU & CUC & CUC & \textbf{GAG} & GAG \\
Acide aminé & Thr & Thr & Leu & Leu & \textbf{Glu} & Glu \\
\hline
ARNm $S$ & ACC & ACU & CUC & CUC & \textbf{GUG} & GAG \\
Acide aminé & Thr & Thr & Leu & Leu & \textbf{Val} & Glu \\
\end{tabular}
\end{center}
La substitution est : \textbf{Glutamate (Glu, acide, hydrophile) $\rightarrow$ Valine (Val, hydrophobe)} en \textbf{position 6} de la chaîne $\beta$-globine.
D'après Doc 3 : GAG $\rightarrow$ Glu ; GUG $\rightarrow$ Val. Confirmé.

\textbf{5. Chaîne de causalité Gène $\rightarrow$ Protéine $\rightarrow$ Caractère.}
\begin{enumerate}
\item \textbf{Niveau ADN (Gène) :} Substitution d'une paire de bases (A$\leftrightarrow$T sur brin sens, soit G$\leftrightarrow$C sur brin matrice) au codon 6 du gène $\beta$-globine ($HBB$ sur chr 11). Allèle $A$ (normal) $\rightarrow$ Allèle $S$ (muté).
\item \textbf{Niveau ARNm (Transcription) :} Le codon 6 de l'ARNm passe de \textbf{GAG} à \textbf{GUG}.
\item \textbf{Niveau Protéine (Traduction) :} L'acide aminé en position 6 de la $\beta$-globine passe de \textbf{Glutamate (Glu)} à \textbf{Valine (Val)}. L'hémoglobine formée ($\alpha_2\beta_2^S$) est l'hémoglobine S (HbS).
\item \textbf{Niveau Caractère (Phénotype) :} La Valine hydrophobe en position 6$\beta$ crée une \emph{tache hydrophobe} à la surface de la $\beta$-globine. En conditions de basse pression partielle en $O_2$ (hypoxie), l'HbS polymérise en longues fibres rigides qui déforment le globule rouge en \emph{faucille} (drépanocyte).
\item \textbf{Conséquences cliniques :} Les drépanocytes sont rigides $\rightarrow$ obstruction microvasculaire (crises vaso-occlusives douloureuses, infarctus tissulaires) + destruction hémolytique précoce (anémie hémolytique chronique, ictère) + susceptibilité aux infections (hyposplénisme fonctionnel).
\end{enumerate}
\textbf{Un seul nucléotide modifié suffit à expliquer tout le phénotype maladie.}

\textbf{6. Interprétation de la courbe (Doc 4).}
\begin{itemize}
\item \textbf{Phase $S$ (Interphase) :} Passage de $2Q$ à $4Q$. \textbf{Réplication de l'ADN} (synthèse). Chaque chromosome se duplique en deux chromatides sœurs. La cellule passe de $2n$ chromosomes simples à $2n$ chromosomes doubles.
\item \textbf{Méiose I (Division réductionnelle) :} Passage de $4Q$ à $2Q$. \textbf{Séparation des chromosomes homologues} (paires de chromosomes paternels/maternels). Chaque cellule fille reçoit $n$ chromosomes, chacun encore constitué de \textbf{deux chromatides sœurs}. C'est ici que les allèles $A$ et $S$ du gène $\beta$-globine (portés par les chromosomes 11 homologues) se séparent.
\item \textbf{Méiose II (Division équationnelle) :} Passage de $2Q$ à $Q$. \textbf{Séparation des chromatides sœurs}. Chaque spermatozoïde reçoit $n$ chromosomes \textbf{simples} (une chromatide = une molécule d'ADN). Quantité finale $Q$.
\end{itemize}
Le père $[A//S]$ produit donc \textbf{50\,\% de spermatozoïdes porteurs de $A$} et \textbf{50\,\% porteurs de $S$}.

\textbf{7. Conseil génétique (rédaction).}
\begin{quote}
\textit{« Monsieur, Madame Kamdem, l'analyse génétique confirme que vous êtes tous les deux \textbf{porteurs sains (hétérozygotes $[A//S]$)} de l'allèle drépanocytaire. Vous ne présentez pas la maladie, mais vous pouvez la transmettre. À chaque grossesse, les risques sont identiques et indépendants du sexe de l'enfant : \textbf{25\,\% de risque d'avoir un enfant drépanocytaire $[S//S]$}, 50\,\% de chances d'avoir un enfant porteur sain comme vous $[A//S]$, et 25\,\% de chances d'avoir un enfant indemne $[A//A]$. Ce statut de porteur vous confère d'ailleurs une résistance partielle au paludisme grave. Nous vous recommandons un suivi régulier de vos enfants atteints au centre de drépanocytose (hydratation, prophylaxie anti-infectieuse, vaccination, hydroxyurée si indiquée) et de proposer le dépistage à vos proches. »}
\end{quote}
\end{exoresolu}

\begin{attention}[Pièges fréquents et points de vigilance]
\begin{itemize}
\item \textbf{Sens de transcription :} Ne jamais transcrire directement le brin matrice 5' $\rightarrow$ 3'. Il faut \textbf{l'inverser (le lire 3' $\rightarrow$ 5')} pour obtenir l'ARNm 5' $\rightarrow$ 3'. L'erreur classique donne une séquence ARNm inversée et fausse.
\item \textbf{Brin matrice vs Brin sens :} L'ARNm est \textbf{complémentaire} du brin matrice (A$\leftrightarrow$U, T$\leftrightarrow$A, G$\leftrightarrow$C) et \textbf{identique} au brin sens (sauf U/T). Vérifie toujours ton ARNm en le comparant au brin sens déduit.
\item \textbf{Mutation drépanocytose :} C'est une \textbf{substitution} (missense), pas une délétion ni une insertion. Le cadre de lecture est conservé.
\item \textbf{Méiose I vs II :} $4Q \rightarrow 2Q$ = Méiose I (séparation homologues, \textbf{réductionnelle}). $2Q \rightarrow Q$ = Méiose II (séparation chromatides, \textbf{équationnelle}). Ne pas inverser.
\item \textbf{Portage sain :} Un hétérozygote $[A//S]$ n'est \textbf{pas malade} (pas d'anémie, pas de crises). Ne pas écrire « porteur malade » ou « un peu malade ». Le terme exact est \emph{porteur sain} ou \emph{hétérozygote asymptomatique}.
\item \textbf{Conseil génétique :} Le risque (25\,\%) est \textbf{par grossesse}, pas « un enfant sur quatre au total ». Chaque grossesse est un événement indépendant.
\end{itemize}
\end{attention}

\courssub{Bilan de l'activité}

Cette enquête t'a permis de mobiliser l'ensemble des savoir-faire de la leçon « Du gène à la protéine » dans un cas réel et fréquent au Cameroun. Tu as vérifié que :
\begin{itemize}
\item L'information génétique circule selon le schéma \textbf{ADN $\xrightarrow{\text{transcription}}$ ARNm $\xrightarrow{\text{traduction}}$ Protéine}, avec des règles de complémentarité et de lecture strictes (polarité, cadre de lecture).
\item Une \textbf{mutation ponctuelle} (un seul nucléotide) peut modifier un acide aminé (substitution Glu$\rightarrow$Val) et bouleverser la structure de la protéine (polymérisation HbS), entraînant un phénotype pathologique majeur (drépanocytose).
\item La \textbf{méiose}, par ses deux divisions successives, réduit la quantité d'ADN de $4Q$ à $Q$ et assure la ségrégation des allèles (loi de Mendel), base de la transmission héréditaire.
\end{itemize}
Tu as enfin produit un \textbf{message de santé publique} : le dépistage du statut drépanocytaire (électrophorèse d'hémoglobine) avant le mariage ou en début de grossesse est un enjeu majeur de prévention au Cameroun, pays à haute prévalence. La connaissance de son génotype ($[A//A]$, $[A//S]$, $[S//S]$) permet des choix reproductifs éclairés.

\begin{savaistu}[La drépanocytose au Cameroun : un enjeu de santé publique]
Au Cameroun, la prévalence du trait drépanocytaire (hétérozygotes $[A//S]$) est de \textbf{20 à 25\,\%} de la population générale, avec des pics dans les zones forestières du Sud et de l'Ouest (région de Bafoussam). La forme homozygote $[S//S]$ (maladie) touche environ \textbf{1,5 à 2\,\% des naissances} (soit 6 000 à 8 000 nouveaux cas/an). Le programme national de lutte contre la drépanocytose (PNLCD) promeut le \textbf{dépistage néonatal systématique} (test de Guthrie ou électrophorèse sur sang de cordon) et la \textbf{prise en charge précoce} (pneumocoque vaccine, acide folique, pénicilline prophylactique, hydroxyurée) qui a fait chuter la mortalité infantile. Le conseil génétique prénuptial est recommandé par le Ministère de la Santé Publique.
\end{savaistu}

\newpage

\coursec{Méthodes et exercices résolus}

\coursec{I. Structure architecturale des acides nucléiques}

\begin{definition}[Nucléotide, monomère des acides nucléiques]
Un \cle{nucléotide} est l'unité monomère des acides nucléiques. Il est constitué de trois éléments covalemment liés :
\begin{enumerate}
\item un \cle{sucre} pentose : le \cle{désoxyribose} dans l'ADN (absence d'un groupement OH sur le carbone 2'), le \cle{ribose} dans l'ARN ;
\item un \cle{groupement phosphate} (acide phosphorique) estérifié sur le carbone 5' du sucre ;
\item une \cle{base azotée} : \cle{puriques} (double cycle) : \cle{adénine} (A) et \cle{guanine} (G) ; \cle{pyrimidiques} (cycle simple) : \cle{cytosine} (C), \cle{thymine} (T) pour l'ADN, \cle{uracile} (U) pour l'ARN.
\end{enumerate}
\end{definition}

\begin{propriete}[Structure de l'ADN : le modèle de Watson et Crick (1953)]
La molécule d'ADN est une \cle{double hélice} droite (forme B) formée de deux brins \cle{antiparallèles} enroulés autour d'un même axe.
\begin{itemize}
\item \textbf{Squelette sucre-phosphate} : les nucléotides sont liés par des \cle{liaisons phosphodiester} entre le phosphate du carbone 5' d'un nucléotide et l'hydroxyle du carbone 3' du suivant. Chaque brin a une polarité $5' \rightarrow 3'$.
\item \textbf{Appariement des bases} : les bases des deux brins sont tournées vers l'intérieur et s'apparient spécifiquement par \cle{liaisons hydrogène} selon la \cle{règle de complémentarité} (Chargaff) : A = T (2 liaisons H) et C $\equiv$ G (3 liaisons H).
\item \textbf{Géométrie} : diamètre $\approx$ \SI{2}{nm}, pas d'hélice $\approx$ \SI{3,4}{nm} (10 paires de bases par tour), sillon majeur et sillon mineur.
\end{itemize}
La complémentarité des bases implique que la séquence d'un brin dicte celle de l'autre : c'est le fondement de la réplication et de la transcription.
\end{propriete}

\begin{propriete}[Structure de l'ARN]
L'ARN est généralement \textbf{monocaténaire} (un seul brin). Il diffère de l'ADN par le sucre (ribose) et la base uracile (U) qui remplace la thymine. Le brin unique peut s'apparier avec lui-même (structures secondaires : tiges-boucles, structure en trèfle de l'ARNt) ou avec un brin d'ADN (hybride ADN-ARN lors de la transcription).
\end{propriete}

\begin{popfigure}
\begin{tikzpicture}[scale=0.9, every node/.style={font=\small}]
% Nucleotide
\begin{scope}[xshift=-7cm]
\node[draw, rounded corners, fill=popblueL, inner sep=2pt] at (0,1.5) {Nucléotide d'ADN};
\node[circle, draw, fill=poporangeL, minimum size=1.2cm] (base) at (0,0) {Base azotée (A,T,G,C)};
\node[rectangle, draw, fill=popgreenL, minimum width=1.5cm, minimum height=0.7cm] (sucre) at (0,-1.2) {Désoxyribose};
\node[draw, fill=poppinkL, minimum width=1cm, minimum height=0.5cm] (phos) at (0,-2.3) {Phosphate};
\draw[thick] (base.south) -- (sucre.north);
\draw[thick] (sucre.south) -- (phos.north);
\node[left] at (-1.8,0) {C1'};
\node[left] at (-1.8,-1.2) {C3' / C5'};
\end{scope}

% DNA Double Helix schematic
\begin{scope}[xshift=2cm]
\node[draw, rounded corners, fill=popblueL, inner sep=2pt] at (0,3.5) {Double hélice ADN (schématique)};
% Backbones
\draw[poporange, thick] (-1.5,3) .. controls (-1.5,2) and (-1.5,1) .. (-1.5,0) .. controls (-1.5,-1) and (-1.5,-2) .. (-1.5,-2.5);
\draw[poporange, thick] (1.5,3) .. controls (1.5,2) and (1.5,1) .. (1.5,0) .. controls (1.5,-1) and (1.5,-2) .. (1.5,-2.5);
% Base pairs
\foreach \y/\bases/\col in {2.5/{A=T}/popgreen, 1.5/{C$\equiv$G}/poppurple, 0.5/{T=A}/popgreen, -0.5/{G$\equiv$C}/poppurple, -1.5/{A=T}/popgreen} {
    \draw[\col, thick] (-1.5,\y) -- (1.5,\y) node[midway, above, sloped, font=\tiny] {\bases};
}
% Arrows for antiparallel
\draw[->, thick, popdark] (-1.8,2.8) -- (-1.8,1.5) node[left, midway] {$5' \rightarrow 3'$};
\draw[->, thick, popdark] (1.8,-1.5) -- (1.8,-2.8) node[right, midway] {$5' \rightarrow 3'$};
\draw[->, thick, popdark] (-1.8,-1.5) -- (-1.8,-2.8) node[left, midway] {$3' \rightarrow 5'$};
\draw[->, thick, popdark] (1.8,2.8) -- (1.8,1.5) node[right, midway] {$3' \rightarrow 5'$};
\node[align=center, font=\tiny] at (0,-3.2) {Squelette sucre-phosphate\\Liaisons H entre bases};
\end{scope}
\end{tikzpicture}
\legende{Organisation architecturale des acides nucléiques. À gauche : composition d'un nucléotide d'ADN. À droite : schéma simplifié de la double hélice montrant l'antiparallélisme des brins (flèches), l'appariement complémentaire des bases (A=T, C$\equiv$G) et le squelette sucre-phosphate.}
\end{popfigure}

\courssub{Méthode 1 : Décrire l'organisation architecturale des acides nucléiques}

\begin{methode}[Décrire la structure de l'ADN et de l'ARN]
Pour décrire rigoureusement un acide nucléique au Bac, procède toujours du plus petit au plus grand niveau d'organisation :
\begin{enumerate}
\item \textbf{Identifier le monomère} : le \cle{nucléotide}, constitué de trois éléments : un \cle{sucre} (désoxyribose pour l'ADN, ribose pour l'ARN), un \cle{groupement phosphate} et une \cle{base azotée}.
\item \textbf{Citer les bases azotées} : pour l'ADN : adénine (A), guanine (G), cytosine (C), thymine (T) ; pour l'ARN : l'uracile (U) remplace la thymine. Préciser : puriques (A, G) et pyrimidiques (C, T, U).
\item \textbf{Décrire l'assemblage} : les nucléotides sont liés par des \cle{liaisons phosphodiester} entre le groupement phosphate d'un nucléotide et le sucre du suivant, formant un brin orienté $5' \rightarrow 3'$.
\item \textbf{Préciser la structure spatiale} : pour l'ADN, \textbf{deux} brins antiparallèles enroulés en \cle{double hélice}, unis par des liaisons hydrogène selon la \cle{complémentarité des bases} : A = T (2 liaisons H) et C $\equiv$ G (3 liaisons H). L'ARN est généralement monocaténaire (un seul brin).
\item \textbf{Conclure par la conséquence} : la complémentarité des bases est le fondement de la réplication et de la transcription.
\end{enumerate}
\textbf{Réflexe de rédaction} : toute description doit être accompagnée d'un schéma légendé si l'énoncé le demande, avec un vocabulaire exact (jamais « branche » pour « brin »).
\end{methode}

\begin{exoresolu}[Analyse d'un fragment d'ADN]
\textbf{Énoncé.} On analyse un fragment d'ADN double brin extrait de cellules du foie d'un patient au CHU de Yaoundé. L'un des brins porte la séquence suivante, lue dans le sens $5' \rightarrow 3'$ :
\[ 5'-\text{A T G C C A T T G}-3' \]
\begin{enumerate}
\item Donner la composition d'un nucléotide d'ADN.
\item Écrire la séquence du brin complémentaire en précisant son orientation.
\item Calculer le nombre total de liaisons hydrogène assurant la cohésion des deux brins de ce fragment.
\end{enumerate}
\tcblower
\textbf{Solution détaillée.}
\begin{enumerate}
\item Un nucléotide d'ADN est constitué de trois éléments : un \textbf{désoxyribose} (sucre à 5 atomes de carbone), un \textbf{groupement phosphate} et une \textbf{base azotée} parmi quatre : adénine (A), thymine (T), guanine (G) ou cytosine (C).
\item Les deux brins d'ADN sont \textbf{antiparallèles} et unis par la complémentarité des bases (A avec T, C avec G). On écrit donc le brin complémentaire en le lisant $3' \rightarrow 5'$ sous le premier :
\begin{align*}
\text{Brin donné : } & 5'-\text{A T G C C A T T G}-3'\\
\text{Brin complémentaire : } & 3'-\text{T A C G G T A A C}-5'
\end{align*}
Soit, lu dans le sens $5' \rightarrow 3'$ : $5'-\text{C A A T G G C A T}-3'$.
\item Les paires A--T comportent \textbf{2 liaisons hydrogène} et les paires C--G en comportent \textbf{3}. Dénombrons les paires du fragment (9 paires de bases) :
\begin{itemize}
\item Paires A--T : positions 1 (A-T), 2 (T-A), 6 (A-T), 7 (T-A), 8 (T-A) $\rightarrow$ \textbf{5 paires} A--T ;
\item Paires C--G : positions 3 (G-C), 4 (C-G), 5 (C-G), 9 (G-C) $\rightarrow$ \textbf{4 paires} C--G.
\end{itemize}
Nombre total de liaisons hydrogène :
\[ N = 5 \times 2 + 4 \times 3 = 10 + 12 = 22 \text{ liaisons hydrogène.} \]
\end{enumerate}
\textbf{Conclusion} : le brin complémentaire est $3'-\text{T A C G G T A A C}-5'$, et la cohésion du fragment est assurée par \textbf{22 liaisons hydrogène}. On remarque que les régions riches en C--G sont plus stables thermiquement, car elles exigent plus d'énergie pour être séparées.
\end{exoresolu}

\courssub{Méthode 2 : Réaliser une maquette de molécule d'ADN}

\begin{methode}[Construire une maquette d'ADN]
La réalisation d'une maquette (TP fréquent au lycée, exploitable dans une question de synthèse) suit une démarche ordonnée :
\begin{enumerate}
\item \textbf{Choisir le matériel} : perles ou papiers de 4 couleurs différentes pour les 4 bases azotées (ex. rouge = A, vert = T, bleu = C, jaune = G), bâtonnets ou trombones pour les groupements phosphate, cure-dents ou fils pour le sucre.
\item \textbf{Assembler les nucléotides} : fixer à chaque « sucre » un « phosphate » et une « base » ; chaque nucléotide est donc un triplet sucre--phosphate--base.
\item \textbf{Former un premier brin} : relier les nucléotides par liaisons phosphate--sucre (squelette sucre-phosphate), les bases pointant vers l'intérieur.
\item \textbf{Former le second brin en respectant la complémentarité} : A face à T, C face à G, et en orientant le brin en sens inverse (antiparallélisme).
\item \textbf{Torsader} l'ensemble pour obtenir la double hélice : un tour complet tous les 10 nucléotides environ (pas d'hélice $\approx$ \SI{3,4}{nm}, diamètre $\approx$ \SI{2}{nm}).
\item \textbf{Vérifier} : compter les paires de bases, contrôler que A est toujours face à T et C face à G.
\end{enumerate}
\textbf{Intérêt pédagogique} : la maquette matérialise la complémentarité des bases, clé de la réplication et de la transcription.
\end{methode}

\begin{exoresolu}[Conception et exploitation d'une maquette]
\textbf{Énoncé.} Un élève de Terminale C du lycée de Ngaoundéré réalise une maquette d'un fragment d'ADN de 12 paires de bases. Il dispose de 15 bases adénine, 10 bases thymine, 9 bases cytosine et 8 bases guanine.
\begin{enumerate}
\item Citer deux règles qu'il doit respecter pour assembler les deux brins.
\item Vérifier si son stock de bases est suffisant pour construire un fragment contenant 7 paires A--T et 5 paires C--G.
\item Expliquer en quoi cette maquette permet de comprendre le mécanisme de la réplication.
\end{enumerate}
\tcblower
\textbf{Solution détaillée.}
\begin{enumerate}
\item Deux règles de construction :
\begin{itemize}
\item \textbf{Complémentarité des bases} : une adénine fait toujours face à une thymine (2 liaisons H) et une cytosine toujours face à une guanine (3 liaisons H) ;
\item \textbf{Antiparallélisme} : les deux brins sont orientés en sens inverse, l'un lu $5' \rightarrow 3'$, l'autre $3' \rightarrow 5'$.
\end{itemize}
\item Le fragment comporte 12 paires de bases, soit 24 nucléotides au total.
\begin{itemize}
\item 7 paires A--T nécessitent $7 \times 1 = 7$ adénines et $7$ thymines ;
\item 5 paires C--G nécessitent $5$ cytosines et $5$ guanines.
\end{itemize}
Comparaison avec le stock :
\begin{center}
\begin{tabularx}{\linewidth}{|l|X|X|X|}\hline
\rowcolor{popblueL} \textbf{Base} & \textbf{Besoin} & \textbf{Stock} & \textbf{Suffisant ?} \\ \hline
A & 7 & 15 & Oui (reste 8) \\ \hline
T & 7 & 10 & Oui (reste 3) \\ \hline
C & 5 & 9 & Oui (reste 4) \\ \hline
G & 5 & 8 & Oui (reste 3) \\ \hline
\end{tabularx}
\end{center}
Le stock est donc \textbf{largement suffisant}.
\item La maquette montre que chaque brin détermine obligatoirement la séquence du brin complémentaire. Lors de la réplication, les deux brins se séparent (rupture des liaisons H) et chacun sert de \textbf{matrice} : les nucléotides libres viennent s'apparier selon la complémentarité. On obtient ainsi deux molécules d'ADN identiques à la molécule initiale, chacune formée d'un brin ancien et d'un brin nouveau : c'est la \cle{réplication semi-conservative}.
\end{enumerate}
\textbf{Conclusion} : la maquette, construite selon les règles de complémentarité et d'antiparallélisme, matérialise le principe de la réplication semi-conservative : chaque brin parental sert de matrice à la synthèse d'un brin nouveau.
\end{exoresolu}

\coursec{II. La duplication de l'ADN : la réplication semi-conservative}

\begin{definition}[Réplication de l'ADN]
La \cle{réplication} (ou duplication) de l'ADN est le processus par lequel une molécule d'ADN double brin donne naissance à deux molécules filles identiques à la molécule mère, avant la division cellulaire (phase S de l'interphase). Elle est \cle{semi-conservative} : chaque molécule fille conserve un brin parental et reçoit un brin néosynthétisé.
\end{definition}

\begin{propriete}[Acteurs et étapes de la réplication (procaryotes/eucaryotes)]
\begin{itemize}
\item \textbf{Hélicase} : ouvre la double hélice en brisant les liaisons H $\rightarrow$ \cle{fourche de réplication}.
\item \textbf{Protéines SSB} : stabilisent les brins simples.
\item \textbf{Primase (ARN polymérase)} : synthétise un \cle{amorce d'ARN} (court fragment d'ARN) pour fournir un groupement 3'-OH libre.
\item \textbf{ADN polymérase} : ajoute des dNTP (désoxyribonucléotides triphosphates) complémentaires au brin matrice \textbf{uniquement dans le sens $5' \rightarrow 3'$} (lecture matrice $3' \rightarrow 5'$). Hydrolyse les 2 phosphates terminaux pour l'énergie.
\item \textbf{Brin conducteur (leading strand)} : synthèse continue dans le sens de progression de la fourche.
\item \textbf{Brin retardé (lagging strand)} : synthèse discontinue par fragments d'\cle{Okazaki} (sens opposé à la fourche).
\item \textbf{ADN polymérase I (exonucléase 5'$\rightarrow$3')} : retire les amorces d'ARN et les remplace par de l'ADN.
\item \textbf{ADN ligase} : scelle les entailles (liaisons phosphodiester) entre fragments d'Okazaki et après remplacement des amorces.
\item \textbf{Topoisomérase (Gyrase)} : relâche les superenroulements en amont de la fourche.
\end{itemize}
\end{propriete}

\begin{experience}[Expérience de Meselson et Stahl (1958) : preuve de la réplication semi-conservative]
\textbf{Protocole} :
\begin{enumerate}
\item Culture de bactéries (\emph{E. coli}) pendant plusieurs générations en milieu $^{15}$N (azote lourd) $\rightarrow$ ADN « lourd » (densité élevée).
\item Transfert en milieu $^{14}$N (azote léger) et prélèvement à chaque génération.
\item Centrifugation en gradient de densité de chlorure de césium (CsCl) : l'ADN se bande à sa densité de flottement.
\end{enumerate}
\textbf{Résultats} :
\begin{itemize}
\item Génération 0 : 1 bande « lourde ».
\item Génération 1 : 1 bande « intermédiaire » (hybride $^{15}$N-$^{14}$N).
\item Génération 2 : 1 bande « intermédiaire » + 1 bande « légère » ($^{14}$N-$^{14}$N).
\end{itemize}
\textbf{Interprétation} : L'apparition d'une bande unique intermédiaire à G1 exclut le mode conservatif (qui donnerait 2 bandes : lourde + légère) et le mode dispersif (bande unique mais densité changeante). Seule la réplication semi-conservative prédit ces résultats.
\end{experience}

\courssub{Méthode 5 : Résoudre un exercice sur la réplication semi-conservative}

\begin{methode}[Exploiter les données de la réplication]
\begin{enumerate}
\item \textbf{Rappeler le principe} : avant la division cellulaire (phase S de l'interphase), chaque molécule d'ADN produit deux molécules filles identiques, chacune composée d'\textbf{un brin parental} et d'\textbf{un brin néosynthétisé} : réplication \cle{semi-conservative}.
\item \textbf{Citer les acteurs} : ADN parental (matrice), \cle{ADN polymérase}, hélicase (ouverture de la double hélice), primase, nucléotides libres (dNTP), énergie.
\item \textbf{Formules à retenir} : après $n$ cycles de réplication à partir d'une molécule d'ADN :
\[ \text{Nombre de molécules d'ADN} = 2^n \]
\[ \text{Nombre de molécules contenant encore un brin parental} = 2 \quad (\text{pour } n \geq 1) \]
\item \textbf{Exploiter un marquage isotopique} (type Meselson-Stahl, $^{15}$N / $^{14}$N) : après 1 génération en milieu léger, tout l'ADN est \textbf{hybride} (intermédiaire) ; après 2 générations, moitié hybride, moitié léger.
\item \textbf{Conclure toujours} en reliant le résultat au caractère semi-conservatif.
\end{enumerate}
\end{methode}

\begin{exoresolu}[Expérience de Meselson et Stahl]
\textbf{Énoncé.} Des bactéries sont cultivées pendant plusieurs générations dans un milieu contenant de l'azote lourd $^{15}$N : tout leur ADN est « lourd ». On les transfère ensuite dans un milieu à azote léger $^{14}$N et on prélève l'ADN après chaque génération ; on le sépare par centrifugation en gradient de densité. Résultats :
\begin{center}
\begin{tabularx}{\linewidth}{|l|X|X|X|}\hline
\rowcolor{popblueL} \textbf{Génération} & \textbf{0} & \textbf{1} & \textbf{2} \\ \hline
Bandes observées & 1 bande lourde & 1 bande intermédiaire & 1 bande intermédiaire + 1 bande légère \\ \hline
\end{tabularx}
\end{center}
\begin{enumerate}
\item Interpréter les résultats des générations 1 et 2.
\item Montrer que ces résultats permettent de rejeter l'hypothèse d'une réplication conservative.
\item Calculer la proportion de molécules d'ADN totalement « légères » à la 4\ieme{} génération.
\end{enumerate}
\tcblower
\textbf{Solution détaillée.}
\begin{enumerate}
\item \textbf{Génération 1} : chaque molécule d'ADN est formée d'un brin lourd parental ($^{15}$N) et d'un brin léger néosynthétisé ($^{14}$N) : tout l'ADN est \textbf{hybride}, d'où une unique bande intermédiaire. \textbf{Génération 2} : les molécules hybrides se répliquent à nouveau ; le brin lourd parental engendre une molécule hybride, le brin léger parental engendre une molécule totalement légère : on obtient \textbf{50 \% d'ADN hybride et 50 \% d'ADN léger}, d'où les deux bandes.
\item Dans l'hypothèse d'une réplication \textbf{conservative}, la molécule parentale resterait intacte et une molécule fille entièrement nouvelle serait synthétisée. À la génération 1, on observerait alors \textbf{deux bandes} : une bande lourde (ADN parental conservé) et une bande légère (ADN néoformé). Or on observe une \textbf{unique bande intermédiaire} : l'hypothèse conservative est donc \textbf{rejetée}. Les résultats sont en revanche exactement ceux prévus par la réplication \textbf{semi-conservative} : chaque molécule fille conserve un brin parental.
\item À chaque génération, le nombre de molécules double : à la génération $n$, il y a $2^n$ molécules. À la 4\ieme{} génération : $2^4 = 16$ molécules. Seules \textbf{2 molécules} conservent un brin parental lourd (elles restent hybrides) ; les $16 - 2 = 14$ autres sont totalement légères.
\[ \text{Proportion d'ADN léger} = \frac{14}{16} = 0{,}875 = 87{,}5\,\%. \]
\end{enumerate}
\textbf{Conclusion} : l'apparition d'une bande intermédiaire unique à la génération 1, puis de deux bandes (intermédiaire et légère) à la génération 2, démontre que la réplication de l'ADN est \textbf{semi-conservative} ; à la 4\ieme{} génération, 87,5 \% des molécules sont totalement légères et seules 2 molécules sur 16 portent encore un brin parental.
\end{exoresolu}

\begin{popfigure}
\begin{tikzpicture}[scale=0.8, every node/.style={font=\small}]
% Replication Fork
\draw[poporange, thick] (-3,2) -- (-1,2) node[midway, above] {Brin conducteur $5'\rightarrow3'$};
\draw[poporange, thick] (-3,-2) -- (-1,-2) node[midway, below] {Brin retardé $3'\leftarrow5'$};
\draw[popblue, thick, ->] (-3,2) .. controls (-2,2.5) and (-1,2.5) .. (0,2) node[above right] {Synthèse continue};
\draw[popblue, thick, ->] (-3,-2) .. controls (-2,-2.5) and (-1,-2.5) .. (0,-2) node[below right] {Fragments d'Okazaki};
% Helicase
\draw[poppurple, fill=poppurpleL, thick] (-3.5,0) circle (0.5) node {Hélicase};
\draw[popdark, thick, ->] (-3.5,0.5) -- (-3.5,1.5) node[above] {Ouverture};
\draw[popdark, thick, ->] (-3.5,-0.5) -- (-3.5,-1.5) node[below] {Fourche};
% Primase/Pol
\draw[popgreen, fill=popgreenL, thick] (0.5,2) rectangle (1.5,2.6) node[midway] {ADN Pol III};
\draw[popgreen, fill=popgreenL, thick] (0.5,-2.6) rectangle (1.5,-2) node[midway] {ADN Pol III};
\draw[poppink, fill=poppinkL, thick] (1.5,2.3) rectangle (2.2,2.6) node[midway, font=\tiny] {Primase};
% Ligase
\draw[popgold, fill=popgoldL, thick] (2.5,-2.3) rectangle (3.2,-2) node[midway, font=\tiny] {Ligase};
% Parental strands
\draw[popdark, very thick] (-4,2.5) -- (4,2.5) node[right] {Brin parental $3'\rightarrow5'$};
\draw[popdark, very thick] (-4,-2.5) -- (4,-2.5) node[right] {Brin parental $5'\rightarrow3'$};
% New strands
\draw[popblue, very thick, dashed] (-4,2) -- (4,2) node[right] {Brins néo (ADN)};
\draw[popblue, very thick, dashed] (-4,-2) -- (4,-2) node[right] {Brins néo (ADN)};
% RNA primers
\draw[poppink, thick] (0.5,-2.3) -- (1.5,-2.3) node[midway, above, font=\tiny] {Amorce ARN};
\draw[poppink, thick] (1.8,-2.3) -- (2.5,-2.3) node[midway, above, font=\tiny] {Amorce ARN};
\end{tikzpicture}
\legende{Schéma d'une fourche de réplication. L'hélicase ouvre l'hélice. Le brin conducteur est synthétisé en continu par l'ADN polymérase III. Le brin retardé est synthétisé par fragments (Okazaki) nécessitant des amorces d'ARN (primase) et la ligase pour les souder.}
\end{popfigure}

\coursec{III. La transcription : de l'ADN à l'ARN messager}

\begin{definition}[Transcription]
La \cle{transcription} est la synthèse d'une molécule d'ARN sur un brin d'ADN matrice (ou brin transcrit), catalysée par l'\cle{ARN polymérase}. Chez les eucaryotes, elle a lieu dans le \cle{noyau} et produit un \cle{pré-ARNm} qui subit une \cle{maturation} (capping 5', polyadénylation 3', \cle{épissage} / \cle{splicing} : suppression des \cle{introns}, jonction des \cle{exons}) pour devenir l'ARNm mature exporté vers le cytoplasme.
\end{definition}

\begin{propriete}[Étapes de la transcription]
\begin{enumerate}
\item \textbf{Initiation} : l'ARN polymérase se fixe sur le \cle{promoteur} (séquence spécifique en amont du gène, ex. boîte TATA chez les eucaryotes), ouvre localement l'hélice (bulle de transcription).
\item \textbf{Élargation} : lecture du brin matrice dans le sens $3' \rightarrow 5'$, incorporation de rNTP (A, U, C, G) selon complémentarité (A$\rightarrow$U, T$\rightarrow$A, C$\rightarrow$G, G$\rightarrow$C), synthèse ARNm $5' \rightarrow 3'$.
\item \textbf{Termination} : rencontre d'un signal de terminaison (séquence terminateur), libération de l'ARN polymérase et de l'ARN néoformé.
\end{enumerate}
\end{propriete}

\begin{popfigure}
\begin{tikzpicture}[scale=0.85, every node/.style={font=\small}]
% DNA
\draw[popdark, thick] (-4,1) -- (4,1) node[right] {Brin codant (non transcrit)};
\draw[popdark, thick] (-4,-1) -- (4,-1) node[right] {Brin matrice (transcrit) $3'\rightarrow5'$};
% Promoter
\draw[poppurple, fill=poppurpleL, thick] (-2.5,1.3) rectangle (-1.5,0.7) node[midway] {Promoteur};
% RNA Pol
\draw[popblue, fill=popblueL, thick, rounded corners] (-0.5,1.5) rectangle (1.5,-0.5) node[midway, align=center] {ARN\\Polymérase};
% Transcription bubble
\draw[poporange, thick] (-0.5,1) .. controls (0,1.5) and (1,1.5) .. (1.5,1);
\draw[poporange, thick] (-0.5,-1) .. controls (0,-1.5) and (1,-1.5) .. (1.5,-1);
% mRNA
\draw[popgreen, very thick, ->] (1.5,0) -- (3.5,0) node[right] {ARNm $5'\rightarrow3'$};
% Base pairing in bubble
\foreach \x/\baseD/\baseR in {0.2/T/A, 0.6/A/U, 1.0/C/G, 1.4/G/C} {
    \node[font=\tiny, popdark] at (\x, 1.2) {\baseD};
    \node[font=\tiny, popgreen] at (\x, -0.8) {\baseR};
    \draw[popline, thin] (\x, 1.1) -- (\x, -0.7);
}
% Maturation symbols (eukaryotes)
\node[draw, fill=poppinkL, rounded corners, font=\tiny] at (4.5, 1.5) {Capping 5'};
\node[draw, fill=poppinkL, rounded corners, font=\tiny] at (4.5, 0) {Épissage (Introns $\downarrow$)};
\node[draw, fill=poppinkL, rounded corners, font=\tiny] at (4.5, -1.5) {Poly-A 3'};
\end{tikzpicture}
\legende{Transcription chez les eucaryotes. L'ARN polymérase se lie au promoteur, ouvre l'ADN (bulle) et synthétise l'ARNm complémentaire du brin matrice. Le pré-ARNm subit une maturation (capping, épissage, polyadénylation) dans le noyau avant export.}
\end{popfigure}

\coursec{IV. La traduction : de l'ARNm à la protéine}

\begin{definition}[Traduction]
La \cle{traduction} est la synthèse d'une chaîne polypeptidique (protéine) dirigée par la séquence de l'ARNm, réalisée par les \cle{ribosomes} dans le \cle{cytoplasme} (libres ou fixés au \cle{réticulum endoplasmique rugueux}). Elle utilise les \cle{ARN de transfert} (ARNt) comme adaptateurs entre le code nucléotidique (codons) et le code amino-acidique.
\end{definition}

\begin{propriete}[Acteurs et étapes de la traduction]
\begin{itemize}
\item \textbf{ARNm} : lu par triplets non chevauchants = \cle{codons} (64 codons : 61 sens + 3 stop). Codon d'initiation : \cle{AUG} (Méthionine).
\item \textbf{ARNt} : structure en trèfle (secondaire) / L (tertiaire). Portent un acide aminé spécifique (fixé sur l'ACC 3' terminal CCA par une \cle{aminoacyl-ARNt synthétase}, coût 1 ATP $\rightarrow$ AMP + 2Pi). Possèdent un \cle{anticodon} complémentaire du codon (appariement \cle{wobble} possible en 3ème position).
\item \textbf{Ribosome} : deux sous-unités (40S + 60S = 80S eucaryotes ; 30S + 50S = 70S procaryotes). Sites : A (aminoacyl), P (peptidyl), E (exit). Activité \cle{peptidyl-transférase} (ribozyme de l'ARNr 28S/23S) pour la liaison peptidique.
\item \textbf{Facteurs} : d'initiation (eIF), d'élongation (eEF), de terminaison (eRF), GTPases.
\end{itemize}
\textbf{Étapes} :
\begin{enumerate}
\item \textbf{Initiation} : assemblage sous-unité petite + ARNm (codon AUG) + ARNt initiateur (Met) + sous-unité grande. Coût : GTP.
\item \textbf{Élargation} (cycle) : 1. Entrée ARNtaa au site A (codon-anticodon, eEF1A-GTP). 2. Transfert peptidique (site P $\rightarrow$ site A, peptidyl-transférase). 3. Translocation (ribosome avance d'un codon, eEF2-GTP) : ARNt déacylé sort site E, peptidyl-ARNt passe P $\leftarrow$ A. Coût : 2 GTP/aa.
\item \textbf{Termination} : codon STOP (UAA, UAG, UGA) reconnu par facteur de libération (eRF1) $\rightarrow$ hydrolyse liaison ester peptidyl-ARNt $\rightarrow$ libération polypeptide + dissociation ribosome.
\end{enumerate}
\end{propriete}

\courssub{Méthode 3 : Identifier et localiser les étapes et les acteurs de la biosynthèse des protéines}

\begin{methode}[Décrire transcription et traduction]
Face à un document (micrographie, schéma, expérience de marquage radioactif), identifie chaque étape en répondant à trois questions : \textbf{où ? avec quels acteurs ? quel produit ?}
\begin{enumerate}
\item \textbf{Transcription} : lieu = \cle{noyau} (cellule eucaryote) ; acteurs = ADN (brin matrice ou brin transcrit), \cle{ARN polymérase}, nucléotides libres (A, U, C, G) ; produit = \cle{ARN messager} (ARNm), copie complémentaire du gène. Règle d'appariement : A--U, T--A, C--G, G--C.
\item \textbf{Traduction} : lieu = \cle{cytoplasme} (ribosomes, libres ou fixés sur le réticulum endoplasmique rugueux) ; acteurs = ARNm, \cle{ribosomes}, \cle{ARN de transfert} (ARNt) porteurs d'acides aminés et munis d'un \cle{anticodon}, enzymes et énergie (ATP, GTP) ; produit = \cle{chaîne polypeptidique}.
\item \textbf{Lecture du code} : l'ARNm est lu par \cle{codons} (triplets de nucléotides) ; chaque codon correspond à un acide aminé (ou à un signal stop). Le codon d'initiation est AUG (méthionine).
\item \textbf{Ordre logique} : ADN $\xrightarrow{\text{transcription}}$ ARNm $\xrightarrow{\text{traduction}}$ protéine. Ne jamais inverser.
\item \textbf{Exploiter un marquage} : si des acides aminés radioactifs apparaissent d'abord dans le cytoplasme et des nucléotides marqués dans le noyau, cela confirme les localisations respectives.
\end{enumerate}
\end{methode}

\begin{exoresolu}[Localisation cellulaire de la synthèse des protéines]
\textbf{Énoncé.} Des chercheurs cultivent des cellules hépatiques de rat en présence de leucine radioactive pendant 5 minutes, puis suivent la radioactivité par autoradiographie au cours du temps. Résultats :
\begin{center}
\begin{tabularx}{\linewidth}{|l|X|X|X|}\hline
\rowcolor{popblueL} \textbf{Temps après marquage} & \textbf{5 min} & \textbf{20 min} & \textbf{60 min} \\ \hline
Compartiment le plus radioactif & Ribosomes du RE rugueux & Appareil de Golgi & Vésicules de sécrétion \\ \hline
\end{tabularx}
\end{center}
\begin{enumerate}
\item Nommer l'étape de la biosynthèse des protéines mise en évidence par le marquage à la leucine et préciser son lieu.
\item Expliquer le trajet de la radioactivité observé entre 5 et 60 minutes.
\item Citer les acteurs moléculaires de l'étape identifiée en 1.
\end{enumerate}
\tcblower
\textbf{Solution détaillée.}
\begin{enumerate}
\item La leucine est un \textbf{acide aminé}, c'est-à-dire un monomère des protéines. Son incorporation dans des molécules radioactives au niveau des ribosomes traduit la \textbf{traduction} : l'assemblage des acides aminés en chaîne polypeptidique. Lieu : le \textbf{cytoplasme}, ici sur les ribosomes fixés au réticulum endoplasmique rugueux (RER).
\item Interprétation du trajet :
\begin{itemize}
\item À \textbf{5 min}, la radioactivité est localisée sur les ribosomes du RER : c'est le lieu d'assemblage de la chaîne polypeptidique (traduction de l'ARNm) ;
\item à \textbf{20 min}, elle se retrouve dans l'\textbf{appareil de Golgi} : la protéine néosynthétisée y est transportée pour subir des maturations (repliement, glycosylation) ;
\item à \textbf{60 min}, elle se situe dans les \textbf{vésicules de sécrétion} : la protéine mature est conditionnée pour être exportée hors de la cellule par exocytose.
\end{itemize}
Ce trajet illustre la \textbf{voie de sécrétion} : RER $\rightarrow$ Golgi $\rightarrow$ vésicules $\rightarrow$ exocytose.
\item Les acteurs de la traduction sont :
\begin{itemize}
\item l'\textbf{ARN messager}, support de l'information lu par codons ;
\item les \textbf{ribosomes}, « usines » d'assemblage ;
\item les \textbf{ARN de transfert} (ARNt), chacun portant un acide aminé spécifique et un anticodon complémentaire du codon de l'ARNm ;
\item les \textbf{acides aminés} libres (dont la leucine radioactive ici) ;
\item des \textbf{enzymes} (aminoacyl-ARNt synthétases, peptidyl-transférase) et de l'énergie fournie par l'ATP et le GTP.
\end{itemize}
\end{enumerate}
\textbf{Conclusion} : le marquage à la leucine radioactive démontre que la traduction s'effectue sur les ribosomes du cytoplasme, et que la protéine synthétisée suit ensuite la voie de sécrétion (RER, appareil de Golgi, vésicules) avant son exocytose.
\end{exoresolu}

\begin{popfigure}
\begin{tikzpicture}[scale=0.8, every node/.style={font=\small}]
% Ribosome
\draw[popblue, fill=popblueL, thick] (0,0) ellipse (2.5 and 1.2) node[midway, align=center] {Grande sous-unité (60S)};
\draw[popblue, fill=popblue!50, thick] (0,-1.2) ellipse (2 and 0.8) node[midway, align=center] {Petite sous-unité (40S)};
% mRNA
\draw[popgreen, thick, ->] (-3.5,-0.6) -- (3.5,-0.6) node[right] {ARNm $5'\rightarrow3'$};
% Codons
\foreach \x/\codon in {-2.5/AUG, -1.5/CCU, -0.5/UUA, 0.5/AGG, 1.5/GUU, 2.5/UAA} {
    \node[draw, fill=popgreenL, rounded corners, font=\tiny] at (\x, -0.6) {\codon};
}
% tRNAs
\foreach \x/\aa/\anticodon in {-2.5/Met/UAC, -1.5/Pro/GGA, -0.5/Leu/AAU, 0.5/Arg/UCC, 1.5/Val/CAA} {
    \draw[poporange, thick] (\x, 0.2) -- (\x, 1.2) node[above, font=\tiny] {\aa};
    \node[font=\tiny, popdark] at (\x, 0.6) {anticodon: \anticodon};
    \draw[popline, thin] (\x, -0.2) -- (\x, 0.2);
}
% Peptide chain
\draw[poppink, very thick] (-2.5, 1.5) -- (1.5, 1.5) node[right] {Chaîne polypeptidique naissante};
\foreach \x in {-2.5, -1.5, -0.5, 0.5, 1.5} {
    \draw[poppink, thick] (\x, 1.2) -- (\x, 1.5);
}
% Sites A, P, E
\node[font=\small, popdark] at (-3.2, 1.0) {Site A};
\node[font=\small, popdark] at (-3.2, 0.0) {Site P};
\node[font=\small, popdark] at (-3.2, -1.0) {Site E};
% Peptidyl transferase center
\node[draw, fill=poppurpleL, rounded corners, font=\tiny, align=center] at (3.5, 1.0) {Centre\\peptidyl-\\transférase\\(ARNr)};
\end{tikzpicture}
\legende{Traduction sur le ribosome. L'ARNm est lu codon par codon. Les ARNt apportent les acides aminés (site A), la liaison peptidique se forme (centre peptidyl-transférase), le ribosome translocé (sites A$\rightarrow$P$\rightarrow$E).}
\end{popfigure}

\coursec{V. Le code génétique et la relation gène – protéine – caractère}

\begin{propriete}[Propriétés du code génétique]
Le code génétique est l'ensemble des correspondances entre les 64 codons de l'ARNm et les 20 acides aminés (ou signaux stop). Il est :
\begin{itemize}
\item \textbf{Universel} : quasi identique chez tous les êtres vivants (exceptions mineures : mitochondries, certains protozoaires) $\rightarrow$ preuve d'origine commune, permet le génie génétique (ex. insuline humaine dans bactérie).
\item \textbf{Dégénéré (redondant)} : la plupart des acides aminés sont codés par plusieurs codons (jusqu'à 6 pour Leu, Ser, Arg). La dégénérescence porte souvent sur la 3ème base (hypothèse \cle{wobble} de Crick).
\item \textbf{Non chevauchant} : les codons se lisent séquentiellement, sans partage de nucléotides.
\item \textbf{Sans ponctuation interne} : pas de séparateur entre codons ; le cadre de lecture est fixé par le codon d'initiation AUG.
\item \textbf{Colinéaire} : la séquence des codons dans l'ARNm correspond à la séquence des acides aminés dans la protéine (sauf maturation post-traductionnelle).
\end{itemize}
\end{propriete}

\courssub{Méthode 4 : Exploiter le code génétique (traduction d'une séquence)}

\begin{methode}[Traduire une séquence nucléotidique en protéine]
\begin{enumerate}
\item \textbf{Identifier la molécule donnée} : s'agit-il du brin matrice d'ADN, du brin non transcrit ou de l'ARNm ? C'est le point le plus piégeux de l'énoncé.
\item \textbf{Obtenir l'ARNm} :
\begin{itemize}
\item Si brin \textbf{matrice} (transcrit) donné : écrire le brin complémentaire en remplaçant T par U (A$\rightarrow$U, T$\rightarrow$A, C$\rightarrow$G, G$\rightarrow$C). Orientation : ARNm $5' \rightarrow 3'$ antiparallèle à matrice $3' \rightarrow 5'$.
\item Si brin \textbf{non transcrit} (codant) donné : l'ARNm est \emph{identique} en séquence (sauf T $\rightarrow$ U), même orientation $5' \rightarrow 3'$.
\end{itemize}
\item \textbf{Découper en codons} à partir du codon d'initiation AUG, par groupes de 3 nucléotides, sans chevauchement.
\item \textbf{Utiliser le tableau du code génétique} (fourni à l'examen) : chaque codon donne un acide aminé ; s'arrêter au premier codon stop (UAA, UAG, UGA).
\item \textbf{Écrire la chaîne polypeptidique} dans l'ordre : acide aminé 1 -- acide aminé 2 -- \dots
\item \textbf{Vérifier} : le nombre d'acides aminés = nombre de codons traduits ; un gène de $n$ nucléotides (brin matrice) code au maximum $n/3$ acides aminés.
\end{enumerate}
\textbf{Propriétés du code à citer} : le code génétique est \cle{universel}, \cle{dégénéré} (plusieurs codons pour un même acide aminé) et \cle{non chevauchant}.
\end{methode}

\begin{exoresolu}[Du gène au polypeptide]
\textbf{Énoncé.} On donne la séquence du \textbf{brin matrice} d'un fragment de gène :
\[ 3'-\text{T A C G G A A A T T C C C A A}-5' \]
Extrait du code génétique (codons de l'ARNm) : AUG = méthionine (Met) ; CCU = proline (Pro) ; UUA = leucine (Leu) ; AGG = arginine (Arg) ; GUU = valine (Val) ; UAA, UAG, UGA = stop ; UCA = sérine (Ser).
\begin{enumerate}
\item Écrire la séquence de l'ARNm transcrit.
\item Déterminer la séquence d'acides aminés du polypeptide correspondant.
\item Une mutation remplace le 8\ieme{} nucléotide du brin matrice (A) par un G. Quelle en est la conséquence sur le polypeptide ? Commenter.
\end{enumerate}
\tcblower
\textbf{Solution détaillée.}
\begin{enumerate}
\item L'ARNm est complémentaire du brin matrice, avec l'uracile à la place de la thymine. Appariement : T$\rightarrow$A, A$\rightarrow$U, C$\rightarrow$G, G$\rightarrow$C. Lecture matrice $3'\rightarrow5'$ $\rightarrow$ ARNm $5'\rightarrow3'$.
\begin{align*}
\text{Brin matrice : } & 3'-\text{T A C G G A A A T T C C C A A}-5'\\
\text{ARNm : } & 5'-\text{A U G C C U U U A A G G G U U}-3'
\end{align*}
\item Découpage en codons à partir de AUG :
\[ \text{AUG} \;|\; \text{CCU} \;|\; \text{UUA} \;|\; \text{AGG} \;|\; \text{GUU} \]
Traduction à l'aide du tableau :
\begin{itemize}
\item AUG $\rightarrow$ Met (codon d'initiation) ;
\item CCU $\rightarrow$ Pro ;
\item UUA $\rightarrow$ Leu ;
\item AGG $\rightarrow$ Arg ;
\item GUU $\rightarrow$ Val.
\end{itemize}
Le polypeptide obtenu est donc : \textbf{Met -- Pro -- Leu -- Arg -- Val}.
\item Le 8\ieme{} nucléotide du brin matrice est le second A du triplet $3'$-AAT-$5'$ (positions 7,8,9 : A-A-T). Il devient $3'$-AGT-$5'$. Le codon correspondant de l'ARNm passe de \textbf{UUA} (Leu) à \textbf{UCA} (Ser).
Il s'agit d'une \textbf{mutation par substitution} entraînant le remplacement d'un acide aminé par un autre (mutation \textbf{faux-sens}) : Met -- Pro -- \textbf{Ser} -- Arg -- Val. Un tel changement peut modifier le repliement et donc la fonction de la protéine (c'est le principe de la drépanocytose, où une substitution unique remplace l'acide glutamique par la valine en position 6 de la chaîne $\beta$-globine).
\end{enumerate}
\textbf{Conclusion} : l'ARNm est $5'$-AUG CCU UUA AGG GUU-$3'$ et code le polypeptide Met--Pro--Leu--Arg--Val ; la substitution d'un seul nucléotide peut suffire à changer un acide aminé, illustrant la relation directe entre la séquence du gène et la structure (donc la fonction) de la protéine.
\end{exoresolu}

\begin{savaistu}[La drépanocytose au Cameroun : un exemple de mutation faux-sens]
La drépanocytose (anémie falciforme) est une maladie génétique héréditaire autosomique récessive très fréquente au Cameroun (prévalence du trait drépanocytaire $\approx$ 20-30\% dans certaines zones). Elle est due à une mutation ponctuelle sur le gène $\beta$-globine (chromosome 11) : codon GAG (Glu) $\rightarrow$ GUG (Val) sur l'ARNm (substitution A $\rightarrow$ T sur le brin codant de l'ADN). L'hémoglobine S (HbS) ainsi formée polymérise en conditions d'hypoxie, déformant les globules rouges en « faucilles » (drepanocytes) $\rightarrow$ anémie hémolytique, crises vaso-occlusives, infections. Le dépistage néonatal et la prise en charge précoce (acide folique, vaccination, hydroxyurée) sont des enjeux majeurs de santé publique au Cameroun.
\end{savaistu}

\begin{aretenir}[Relation Gène $\rightarrow$ Protéine $\rightarrow$ Caractère]
\begin{itemize}
\item \textbf{Gène} : portion d'ADN (suite de nucléotides) codant pour une protéine (ou ARN fonctionnel).
\item \textbf{Protéine} : polypeptide plié, structure 3D déterminée par la séquence d'acides aminés (structure primaire).
\item \textbf{Caractère} : phénotype observable (ex. groupe sanguin, couleur des yeux, résistance au paludisme, pathologie).
\item \textbf{Chaîne de causalité} : Mutation ADN $\rightarrow$ modification ARNm $\rightarrow$ modification séquence acides aminés $\rightarrow$ modification structure/fonction protéine $\rightarrow$ modification du caractère.
\item \textbf{Allèle} : version alternative d'un gène (différence de séquence). L'allèle « normal » et l'allèle « mutant » coexistent dans la population.
\end{itemize}
\end{aretenir}

\coursec{VI. Variation de la quantité d'ADN au cours de la méiose}

\begin{propriete}[Variation de la quantité d'ADN et nombre de chromosomes]
Soit une cellule diploïde $2n$ en G1 avec une quantité d'ADN $Q$ (ex. humain : $2n=46$, $Q \approx \SI{6,6}{pg}$).
\begin{center}
\begin{tabularx}{\linewidth}{|l|X|X|X|}\hline
\rowcolor{popblueL} \textbf{Étape} & \textbf{Quantité d'ADN} & \textbf{Nombre de chromosomes} & \textbf{État des chromosomes} \\ \hline
G1 (avant réplication) & $Q$ & $2n$ & Simples (1 chromatide) \\ \hline
Fin Phase S (après réplication) & $2Q$ & $2n$ & Doubles (2 chromatides sœurs) \\ \hline
Fin Méiose I (Télophase I) & $Q$ & $n$ & Doubles (2 chromatides sœurs) \\ \hline
Fin Méiose II (Télophase II) & $Q/2$ & $n$ & Simples (1 chromatide) \\ \hline
\end{tabularx}
\end{center}
\textbf{Rappel} : La méiose comprend une réplication (phase S) suivie de deux divisions successives sans réplication intermédiaire : Méiose I (réductionnelle, séparation des homologues) et Méiose II (équationnelle, séparation des chromatides).
\end{propriete}

\courssub{Méthode 6 : Tracer et interpréter la courbe de variation de la quantité d'ADN au cours de la méiose}

\begin{methode}[Construire la courbe $Q_{\text{ADN}} = f(t)$ au cours de la méiose]
\begin{enumerate}
\item \textbf{Fixer les repères} : on note $Q$ la quantité d'ADN d'une cellule en G1 (avant réplication) et $n$ le nombre de chromosomes simples. Une cellule diploïde au départ possède $2n$ chromosomes et une quantité $Q$ d'ADN.
\item \textbf{Identifier les étapes clés} :
\begin{itemize}
\item \textbf{Interphase (phase S)} : réplication $\Rightarrow$ la quantité passe de $Q$ à $2Q$ (chromosomes à 2 chromatides) ;
\item \textbf{Fin de la méiose I} (division réductionnelle) : chaque cellule fille reçoit $2Q/2 = Q$ (chromosomes doubles, $n$ chromosomes à 2 chromatides) ;
\item \textbf{Fin de la méiose II} (division équationnelle) : chaque cellule reçoit $Q/2$ (chromosomes simples).
\end{itemize}
\item \textbf{Tracer} : axe des abscisses = temps (étapes), axe des ordonnées = quantité d'ADN par cellule (en unités $Q$). La courbe monte progressivement de $Q$ à $2Q$ (phase S), reste à $2Q$ pendant la méiose I, chute à $Q$ (télophase I), puis chute à $Q/2$ (télophase II).
\item \textbf{Interpréter} : la première chute correspond à la séparation des chromosomes homologues (réduction du nombre de chromosomes) ; la seconde à la séparation des chromatides sœurs.
\item \textbf{Conclure} : la méiose produit 4 cellules haploïdes ($n$ chromosomes, $Q/2$ d'ADN) à partir d'une cellule diploïde.
\end{enumerate}
\end{methode}

\begin{exoresolu}[Quantité d'ADN au cours de la spermatogenèse]
\textbf{Énoncé.} Chez un mammifère, une cellule germinale mâle en phase G1 contient une quantité d'ADN $Q = \SI{6,6}{pg}$ (picogrammes) et $2n = 46$ chromosomes (on se placera dans le cas humain).
\begin{enumerate}
\item Donner la quantité d'ADN et le nombre de chromosomes de la cellule à la fin de la phase S, puis dans chaque cellule à l'issue de la méiose I et de la méiose II.
\item Tracer la courbe de variation de la quantité d'ADN par cellule au cours du temps.
\item Expliquer l'origine des deux chutes observées.
\end{enumerate}
\tcblower
\textbf{Solution détaillée.}
\begin{enumerate}
\item \textbf{Fin de la phase S} : la réplication double la quantité d'ADN sans changer le nombre de chromosomes (chaque chromosome est alors formé de 2 chromatides) : $2Q = \SI{13,2}{pg}$, $2n = 46$ chromosomes doubles. \textbf{Fin de la méiose I} : les chromosomes homologues se sont séparés : chaque cellule fille contient $n = 23$ chromosomes (toujours à 2 chromatides) et $Q = \SI{6,6}{pg}$ d'ADN. \textbf{Fin de la méiose II} : les chromatides sœurs se sont séparées : chaque spermatide contient $n = 23$ chromosomes simples et $Q/2 = \SI{3,3}{pg}$ d'ADN.
\begin{center}
\begin{tabularx}{\linewidth}{|l|X|X|X|X|}\hline
\rowcolor{popblueL} \textbf{Étape} & \textbf{G1} & \textbf{Fin S} & \textbf{Fin méiose I} & \textbf{Fin méiose II} \\ \hline
Quantité d'ADN & $Q = \SI{6,6}{pg}$ & $2Q = \SI{13,2}{pg}$ & $Q = \SI{6,6}{pg}$ & $Q/2 = \SI{3,3}{pg}$ \\ \hline
Chromosomes & $2n = 46$ & $2n = 46$ (doubles) & $n = 23$ (doubles) & $n = 23$ (simples) \\ \hline
\end{tabularx}
\end{center}
\item Courbe de variation de la quantité d'ADN par cellule :
\begin{popfigure}
\begin{tikzpicture}
\begin{axis}[
width=0.95\linewidth, height=7cm,
xlabel={Temps (étapes successives)},
ylabel={Quantité d'ADN par cellule (pg)},
xmin=0, xmax=5.5, ymin=0, ymax=14.5,
xtick={0.5,1.5,2.5,3.5,4.5},
xticklabels={G1, Phase S, Méiose I, Méiose II, 4 Spermatides},
xticklabel style={font=\small, rotate=15, anchor=east},
ytick={3.3,6.6,13.2},
yticklabels={$Q/2 = \SI{3{,}3}{pg}$, $Q = \SI{6{,}6}{pg}$, $2Q = \SI{13{,}2}{pg}$},
grid=major, grid style={popline!40},
legend style={at={(0.5,-0.15)}, anchor=north, legend columns=1, font=\small},
]
\addplot[popcurve, mark=*, mark size=2pt, thick] coordinates {
(0,6.6) (0.5,6.6) 
(1.5,13.2) (2.5,13.2) (3.0,13.2)
(3.4,6.6) (3.8,6.6) (4.2,6.6)
(4.6,3.3) (5.2,3.3)
};
\addlegendentry{Quantité d'ADN}
% Annotations
\node[popdark, font=\small, align=center] at (axis cs:1,14) {Réplication (Phase S)\\+$Q \rightarrow 2Q$};
\node[popdark, font=\small, align=center] at (axis cs:3.2,10) {Méiose I\\Séparation homologues};
\node[popdark, font=\small, align=center] at (axis cs:4.8,5) {Méiose II\\Séparation chromatides};
\end{axis}
\end{tikzpicture}
\legende{Variation de la quantité d'ADN par cellule au cours de la méiose. La montée progressive de $Q$ à $2Q$ correspond à la réplication (phase S) ; le palier à $2Q$ correspond à la prophase I / métaphase I ; la première chute brutale ($2Q \rightarrow Q$) marque la fin de la méiose I (anaphase I/télophase I) ; la seconde chute ($Q \rightarrow Q/2$) marque la fin de la méiose II (anaphase II/télophase II).}
\end{popfigure}
\item \textbf{Première chute} ($2Q \rightarrow Q$) : à l'anaphase I, les \textbf{chromosomes homologues} de chaque paire se séparent et migrent vers les pôles opposés ; la cytodiérèse répartit le stock d'ADN en deux cellules à $n$ chromosomes doubles. C'est la \textbf{division réductionnelle} : elle divise par deux le nombre de chromosomes. \textbf{Seconde chute} ($Q \rightarrow Q/2$) : à l'anaphase II, les \textbf{chromatides sœurs} de chaque chromosome se séparent ; chaque cellule fille ne reçoit qu'un exemplaire de chaque chromosome. C'est la \textbf{division équationnelle}.
\end{enumerate}
\textbf{Conclusion} : grâce à une réplication unique suivie de deux divisions successives, la méiose fait passer la quantité d'ADN de $2Q$ à $Q/2$ et le stock chromosomique de $2n$ à $n$ : elle produit 4 cellules haploïdes génétiquement distinctes, condition du maintien de la diploïdie à la fécondation.
\end{exoresolu}

\begin{attention}[Les 5 erreurs qui coûtent des points au Bac]
\begin{enumerate}
\item \textbf{Confondre brin matrice et brin non transcrit.} Si l'énoncé donne le brin \emph{non transcrit} (ou « brin codant »), l'ARNm lui est \emph{identique} (sauf T $\rightarrow$ U) ; s'il donne le brin \emph{matrice}, l'ARNm est \emph{complémentaire}. Lis l'énoncé deux fois avant d'écrire la moindre base.
\item \textbf{Oublier l'orientation des brins.} Écrire le brin complémentaire sans indiquer les extrémités $5'$ et $3'$, ou sans inverser le sens (antiparallélisme), fait perdre des points de rigueur. Toute séquence écrite doit être encadrée par ses polarités.
\item \textbf{Inverser transcription et traduction.} La transcription (ADN $\rightarrow$ ARNm) a lieu dans le \textbf{noyau} et met en jeu l'ARN polymérase ; la traduction (ARNm $\rightarrow$ protéine) a lieu dans le \textbf{cytoplasme} sur les ribosomes. Attribuer la traduction au noyau est une erreur éliminatoire.
\item \textbf{Mal compter les liaisons hydrogène ou les molécules d'ADN.} Rappel : A--T = 2 liaisons H, C--G = 3 liaisons H ; après $n$ réplications, $2^n$ molécules dont seulement 2 conservent un brin parental. Refais les dénombrements posément, paire par paire.
\item \textbf{Tracer la courbe de la méiose avec des chutes au mauvais moment.} La quantité d'ADN passe de $Q$ à $2Q$ \emph{avant} la méiose (phase S, montée progressive), chute à $Q$ à la fin de la méiose I (séparation des homologues) puis à $Q/2$ à la fin de la méiose II (séparation des chromatides). Ne confonds pas avec la courbe du cycle cellulaire mitotique (qui redescend à $Q$ en une seule fois). Et n'oublie jamais de légender les axes avec la grandeur et l'unité.
\end{enumerate}
\end{attention}

\newpage

\coursec{Exercices}

\courssub{Je vérifie mes connaissances}

\begin{exercice}[QCM : les bases de la leçon]
Pour chaque question, choisis la (ou les) bonne(s) réponse(s) et justifie brièvement ton choix.
\begin{enumerate}
\item La transcription d'un gène a lieu :
\begin{enumerate}
\item dans le cytoplasme ;
\item dans le noyau (chez une cellule eucaryote) ;
\item sur les ribosomes ;
\item dans les mitochondries uniquement.
\end{enumerate}
\item La traduction de l'ARNm en protéine a lieu :
\begin{enumerate}
\item dans le noyau ;
\item sur les ribosomes du cytoplasme ;
\item dans le nucléole ;
\item sur la membrane plasmique.
\end{enumerate}
\item L'ARNt (ARN de transfert) a pour rôle :
\begin{enumerate}
\item de porter l'information génétique du noyau vers le cytoplasme ;
\item d'apporter les acides aminés au ribosome grâce à son anticodon ;
\item de constituer la structure du ribosome ;
\item de servir de matrice à la synthèse de l'ARNm.
\end{enumerate}
\item Le code génétique est :
\begin{enumerate}
\item redondant (dégénéré) ;
\item universel (ou quasi universel) ;
\item composé de codons de 3 nucléotides ;
\item chevauchant (les codons se recouvrent).
\end{enumerate}
\item Au cours de la méiose, la quantité d'ADN $Q$ d'une cellule :
\begin{enumerate}
\item double une seule fois, avant la première division ;
\item double avant chacune des deux divisions ;
\item est divisée par deux à chaque division ;
\item reste constante du début à la fin.
\end{enumerate}
\end{enumerate}
\end{exercice}

\begin{exercice}[Vrai ou faux justifié]
Pour chacune des affirmations suivantes, indique si elle est vraie ou fausse, puis justifie ta réponse en une ou deux phrases.
\begin{enumerate}
\item « La réplication de l'ADN est dite semi-conservative parce que chaque molécule fille conserve l'un des deux brins de la molécule mère. »
\item « L'ARNm est synthétisé à partir des deux brins d'ADN simultanément. »
\item « Une protéine de 150 acides aminés est codée par un ARNm comportant au moins 450 nucléotides. »
\item « L'adénine s'apparie avec la thymine par trois liaisons hydrogène. »
\item « Au cours de la méiose, la quantité d'ADN par cellule passe directement de $2Q$ à $Q$ lors de la première division. »
\end{enumerate}
\end{exercice}

\begin{exercice}[Définitions à rédiger]
Définis précisément, en une ou deux phrases chacune, les termes suivants :
\begin{enumerate}
\item nucléotide ;
\item réplication semi-conservative ;
\item transcription ;
\item traduction ;
\item codon et anticodon ;
\item code génétique.
\end{enumerate}
\end{exercice}

\begin{exercice}[Calcul de composition d'un fragment d'ADN]
L'analyse chimique d'un fragment d'ADN double brin extrait de cellules de maïs cultivé dans la région de l'Ouest du Cameroun montre qu'il contient 22 \% d'adénine (A).
\begin{enumerate}
\item Détermine les pourcentages des trois autres bases azotées (T, G, C) dans ce fragment.
\item Justifie tes résultats en t'appuyant sur les règles de complémentarité des bases et sur la structure en double hélice de l'ADN.
\item Calcule le rapport $\dfrac{A + T}{G + C}$ de ce fragment. Que peut-on dire de la richesse de cet ADN en liaisons hydrogène ?
\end{enumerate}
\end{exercice}

\courssub{Je m'entraîne}

\begin{exercice}[D'un brin d'ADN à la protéine]
On donne la séquence suivante du brin non transcrit (brin codant) d'un fragment de gène, lue dans le sens 5' $\rightarrow$ 3' :
\begin{center}
5' -- ATG GCT TAC GGA TTT CAA TGA -- 3'
\end{center}
\begin{enumerate}
\item Écris la séquence du brin transcrit (brin matrice) complémentaire, en précisant sa polarité.
\item Écris la séquence de l'ARNm synthétisé, en précisant sa polarité.
\item À l'aide du tableau du code génétique fourni en annexe de ton manuel, détermine la séquence peptidique traduite.
\item Combien de codons cet ARNm comporte-t-il ? Combien codent effectivement un acide aminé ?
\end{enumerate}
\end{exercice}

\begin{exercice}[Rôles des trois types d'ARN]
Lors d'un TP au lycée, des élèves observent au microscope électronique des cellules de pancréas de rat en pleine synthèse d'insuline.
\begin{enumerate}
\item Nomme les trois types d'ARN impliqués dans la biosynthèse de l'insuline et précise le rôle exact de chacun.
\item Indique le lieu de synthèse de chacun de ces ARN dans la cellule.
\item Schématise par un organigramme fléché le trajet de l'information génétique, du gène de l'insuline jusqu'à la protéine fonctionnelle, en précisant la localisation de chaque étape.
\end{enumerate}
\end{exercice}

\begin{exercice}[Quantité d'ADN au cours du cycle cellulaire]
Chez une espèce animale, la quantité d'ADN contenue dans le noyau d'une cellule somatique en phase G1 est $Q = \SI{6,6}{pg}$ (picogrammes).
\begin{enumerate}
\item Quelle est la quantité d'ADN d'une cellule de cette espèce en phase G2 de l'interphase ? Justifie.
\item Quelle est la quantité d'ADN de chaque cellule fille issue de la mitose ?
\item Quelle est la quantité d'ADN de chaque gamète produit par méiose ?
\item Représente sous forme de tableau les valeurs de la quantité d'ADN ($Q$, $2Q$ ou $Q/2$) et du nombre de chromosomes ($n$ ou $2n$) aux grandes étapes de la méiose, pour une espèce à $2n = 8$.
\end{enumerate}
\end{exercice}

\begin{exercice}[Analyse d'une molécule d'ADN]
Un fragment d'ADN double brin comporte 1\,000 paires de bases.
\begin{enumerate}
\item Combien de nucléotides ce fragment contient-il au total ?
\item Ce fragment contient 300 molécules de guanine. Calcule le nombre de molécules de cytosine, d'adénine et de thymine.
\item Sachant que la distance entre deux paires de bases successives est de \SI{0,34}{nm}, calcule la longueur de ce fragment d'ADN en nanomètres, puis en micromètres.
\item Après trois cycles de réplication \emph{in vitro}, combien de molécules d'ADN obtient-on à partir d'une seule molécule initiale ? Combien de ces molécules contiennent encore un brin de la molécule d'origine ?
\end{enumerate}
\end{exercice}

\courssub{Je me prépare au Bac}

\begin{exercice}[Type Bac — L'expérience de Meselson et Stahl (1958)]
Des chercheurs font croître des bactéries \emph{Escherichia coli} pendant plusieurs générations dans un milieu de culture dont l'unique source d'azote est le chlorure d'ammonium marqué à l'azote lourd $^{15}$N. Tout l'ADN des bactéries est alors « lourd ». Les bactéries sont ensuite transférées dans un milieu contenant uniquement de l'azote léger $^{14}$N. On prélève des échantillons à différents temps de culture, on extrait l'ADN et on le soumet à une centrifugation en gradient de chlorure de césium : l'ADN migre et forme des bandes dont la position dépend de sa densité (bande basse = ADN lourd $^{15}$N/$^{15}$N ; bande intermédiaire = ADN hybride $^{15}$N/$^{14}$N ; bande haute = ADN léger $^{14}$N/$^{14}$N). Les résultats obtenus sont résumés dans le tableau ci-dessous.

\begin{center}
\begin{tabularx}{\linewidth}{|l|X|X|X|}
\hline
\rowcolor{popblueL}
\textbf{Génération} & \textbf{Bande basse (lourde)} & \textbf{Bande intermédiaire (hybride)} & \textbf{Bande haute (légère)} \\
\hline
Génération 0 & 1 bande intense & aucune & aucune \\
\hline
Génération 1 & aucune & 1 bande intense & aucune \\
\hline
Génération 2 & aucune & 1 bande (50 \% de l'ADN) & 1 bande (50 \% de l'ADN) \\
\hline
\end{tabularx}
\end{center}

Trois hypothèses de duplication de l'ADN étaient envisagées à l'époque :
\begin{itemize}
\item \textbf{hypothèse conservative} : la molécule mère reste intacte et une molécule fille entièrement nouvelle est synthétisée ;
\item \textbf{hypothèse semi-conservative} : chaque molécule fille est formée d'un brin parental et d'un brin néosynthétisé ;
\item \textbf{hypothèse dispersive} : chaque brin fils est un assemblage de fragments anciens et de fragments neufs répartis le long de la molécule.
\end{itemize}

\begin{enumerate}
\item Explique le principe de la centrifugation en gradient de densité et précise pourquoi elle permet de distinguer les trois types d'ADN.
\item Pour chacune des trois hypothèses, prédis les bandes attendues à la génération 1 et à la génération 2 (tu pourras t'aider de schémas légendés des molécules d'ADN, en distinguant par deux couleurs les brins lourds et les brins légers).
\item Confronte ces prédictions aux résultats expérimentaux et réfute successivement l'hypothèse conservative et l'hypothèse dispersive.
\item Conclus en démontrant que seule l'hypothèse semi-conservative rend compte de l'ensemble des résultats. Prédis alors le résultat attendu à la génération 3.
\item Décris brièvement le déroulement de la réplication au niveau d'un œil de réplication (enzymes principales, sens de synthèse).
\end{enumerate}
\end{exercice}

\begin{exercice}[Type Bac — Exploitation de séquences et effet des mutations]
Chez l'Homme, un court fragment du brin transcrit (brin matrice) d'un gène présente la séquence suivante de 21 nucléotides, lue dans le sens 3' $\rightarrow$ 5' :
\begin{center}
3' -- TAC CGA ATG CCT AAA GTT ACT -- 5'
\end{center}
On rappelle que la lecture de l'ARNm s'effectue codon par codon, sans chevauchement, à partir du codon d'initiation.

\begin{enumerate}
\item Écris la séquence de l'ARNm correspondant à ce fragment, en précisant sa polarité.
\item À l'aide du tableau du code génétique, détermine la séquence peptidique obtenue par traduction de cet ARNm. Que remarques-tu à la fin de la traduction ?
\item Une première mutation remplace le 7\up{e} nucléotide du brin transcrit (A) par une guanine (G). Écris le nouvel ARNm et la nouvelle séquence peptidique. Ce type de mutation est-il synonyme (silencieux) ou non synonyme ? Justifie.
\item Une seconde mutation correspond à la délétion (perte) du 4\up{e} nucléotide du brin transcrit. Écris la nouvelle séquence de l'ARNm et la séquence peptidique correspondante. Compare l'ampleur des conséquences de cette délétion à celles de la substitution de la question 3, et explique la différence.
\item Une troisième mutation transforme le codon n\textsuperscript{o} 3 de l'ARNm en codon UAA. Quelle en est la conséquence sur la protéine synthétisée ?
\item En t'appuyant sur cet exercice, dégage deux propriétés du code génétique et explique la relation de cause à effet qui relie le gène, la protéine et le caractère.
\end{enumerate}
\end{exercice}

\begin{exercice}[Type Bac — Quantité d'ADN au cours de la méiose et albinisme]
\textbf{Partie A.} Chez une espèce animale dont les cellules somatiques possèdent $2n = 8$ chromosomes, on mesure la quantité d'ADN contenue dans le noyau d'une cellule germinale au cours du temps. On note $Q$ la quantité d'ADN d'une cellule en phase G1. Les mesures sont consignées dans le tableau suivant :

\begin{center}
\begin{tabularx}{\linewidth}{|l|X|X|X|X|X|X|X|}
\hline
\rowcolor{popblueL}
\textbf{Temps (h)} & 0 & 4 & 8 & 12 & 16 & 20 & 24 \\
\hline
\textbf{Quantité d'ADN} & $Q$ & $1{,}5\,Q$ & $2Q$ & $2Q$ & $Q$ & $Q$ & $Q/2$ \\
\hline
\end{tabularx}
\end{center}

\begin{enumerate}
\item Trace la courbe de variation de la quantité d'ADN en fonction du temps (échelle : 1 cm pour 4 h en abscisse ; 1 cm pour $Q/2$ en ordonnée).
\item Identifie sur la courbe, en les justifiant, les phases correspondant : à la réplication de l'ADN, à la première division de la méiose (division de réduction) et à la seconde division (division équationnelle).
\item Indique, pour chacune des quatre cellules obtenues à l'issue de la méiose, la quantité d'ADN et le nombre de chromosomes.
\item Compare cette évolution à celle observée au cours d'une mitose : construis le tableau comparatif des quantités d'ADN et des nombres de chromosomes avant division et après division dans les deux cas, puis dégage la particularité fondamentale de la méiose.
\end{enumerate}

\textbf{Partie B.} L'albinisme est un caractère héréditaire rencontré dans plusieurs populations camerounaises : les individus atteints présentent une dépigmentation de la peau, des poils et des yeux. On sait que la mélanine, pigment responsable de la coloration, est synthétisée à partir de la tyrosine (un acide aminé) grâce à une enzyme, la tyrosinase. Chez les individus albinos, le gène codant la tyrosinase porte une mutation et l'enzyme produite est non fonctionnelle.

\begin{enumerate}
\setcounter{enumi}{4}
\item En mobilisant la relation gène -- protéine -- caractère, explique comment une mutation ponctuelle du gène de la tyrosinase peut conduire au phénotype albinos (tu préciseras la chaîne de causalité, de la séquence nucléotidique jusqu'au caractère observable).
\item Propose une explication au fait que deux mutations différentes du même gène puissent toutes deux aboutir à une tyrosinase non fonctionnelle.
\end{enumerate}
\end{exercice}

\newpage

\coursec{Corrigés des exercices}

\begin{corrige}[Exercice 1 — QCM : les bases de la leçon]
\begin{enumerate}
\item \textbf{Réponse b.} Chez une cellule eucaryote, la transcription (synthèse de l'ARNm à partir de l'ADN) a lieu dans le \cle{noyau}, là où se trouve l'ADN. Le cytoplasme et les ribosomes sont le lieu de la traduction ; les mitochondries possèdent bien leur propre ADN, mais ce n'est pas le lieu de transcription des gènes nucléaires.
\item \textbf{Réponse b.} La traduction s'effectue sur les \cle{ribosomes} du cytoplasme (libres ou fixés au réticulum endoplasmique rugueux).
\item \textbf{Réponse b.} L'ARNt transporte un acide aminé spécifique et le positionne sur le ribosome grâce à son \cle{anticodon}, complémentaire du codon de l'ARNm. C'est l'ARNm qui porte l'information (a) et l'ARNr qui constitue le ribosome (c).
\item \textbf{Réponses a, b et c.} Le code génétique est \cle{redondant} (plusieurs codons pour un même acide aminé), \cle{quasi universel} et lu par triplets non chevauchants de 3 nucléotides. La réponse d est fausse : la lecture est \emph{sans chevauchement}.
\item \textbf{Réponses a et c.} La quantité d'ADN double une seule fois, pendant l'interphase (avant la première division) : $Q \rightarrow 2Q$. Puis elle est divisée par deux à chaque division : $2Q \rightarrow Q$ (méiose I), puis $Q \rightarrow Q/2$ (méiose II). Il n'y a \emph{pas} de réplication entre les deux divisions.
\end{enumerate}
\end{corrige}

\begin{corrige}[Exercice 2 — Vrai ou faux justifié]
\begin{enumerate}
\item \textbf{Vrai.} C'est précisément la définition : chaque molécule d'ADN fille est constituée d'un brin parental conservé et d'un brin néosynthétisé, comme l'ont démontré Meselson et Stahl (1958).
\item \textbf{Faux.} Un seul brin, le \cle{brin transcrit} (ou brin matrice), sert de modèle à la synthèse de l'ARNm ; l'autre brin est le brin non transcrit (codant).
\item \textbf{Vrai.} Chaque acide aminé est codé par un codon de 3 nucléotides : $150 \times 3 = 450$ nucléotides. « Au moins » car il faut ajouter le codon stop (3 nucléotides) et les régions non traduites.
\item \textbf{Faux.} A et T sont unies par \textbf{deux} liaisons hydrogène ; c'est le couple G--C qui en comporte trois.
\item \textbf{Faux.} Avant la méiose, la réplication fait passer la quantité de $Q$ à $2Q$. La première division ramène chaque cellule fille de $2Q$ à $Q$ : c'est exact. Attention toutefois à la formulation : la cellule mère est à $2Q$ \emph{après réplication} ; la première division produit deux cellules à $Q$ (chromosomes à deux chromatides). L'affirmation est donc \textbf{vraie} si l'on considère le passage $2Q \rightarrow Q$ lors de la méiose I. (Rédaction attendue : méiose I : $2Q \rightarrow Q$ ; méiose II : $Q \rightarrow Q/2$.)
\end{enumerate}
\begin{attention}[Point de vigilance]
Ne pas confondre « quantité d'ADN » et « nombre de chromosomes » : après la méiose I, chaque cellule contient $n$ chromosomes mais encore $Q$ ADN (chromosomes à deux chromatides).
\end{attention}
\end{corrige}

\begin{corrige}[Exercice 3 — Définitions à rédiger]
\begin{enumerate}
\item \textbf{Nucléotide} : monomère constitutif des acides nucléiques, formé de l'association d'un phosphate, d'un sucre à 5 atomes de carbone (désoxyribose dans l'ADN, ribose dans l'ARN) et d'une base azotée (A, T/U, G, C).
\item \textbf{Réplication semi-conservative} : duplication de l'ADN précédant toute division cellulaire, au cours de laquelle chaque molécule fille reçoit un brin de la molécule mère (brin conservé) et un brin nouvellement synthétisé.
\item \textbf{Transcription} : synthèse d'une molécule d'ARNm complémentaire de l'un des brins d'ADN d'un gène, catalysée par l'ARN polymérase ; elle a lieu dans le noyau chez les eucaryotes.
\item \textbf{Traduction} : synthèse d'une chaîne polypeptidique sur les ribosomes du cytoplasme, par lecture de l'ARNm codon par codon, chaque codon étant reconnu par l'anticodon d'un ARNt porteur de l'acide aminé correspondant.
\item \textbf{Codon} : triplet de 3 nucléotides de l'ARNm correspondant à un acide aminé (ou à un signal stop). \textbf{Anticodon} : triplet de 3 nucléotides porté par l'ARNt, complémentaire et antiparallèle d'un codon, qui assure la reconnaissance codon--ARNt.
\item \textbf{Code génétique} : ensemble des correspondances entre les 64 codons de l'ARNm et les 20 acides aminés (ou les signaux stop) ; il est redondant, non chevauchant et quasi universel.
\end{enumerate}
\end{corrige}

\begin{corrige}[Exercice 4 — Calcul de composition d'un fragment d'ADN]
\begin{enumerate}
\item D'après les règles de Chargaff (complémentarité A--T et G--C dans un ADN double brin) :
\begin{align*}
\%T &= \%A = 22~\% \\
\%G + \%C &= 100~\% - (22 + 22)~\% = 56~\% \\
\%G &= \%C = \frac{56}{2} = 28~\%
\end{align*}
Le fragment contient donc : A $= 22~\%$, T $= 22~\%$, G $= 28~\%$, C $= 28~\%$.
\item \textbf{Justification.} Dans la double hélice, chaque base d'un brin s'apparie avec sa base complémentaire sur l'autre brin : A avec T (2 liaisons hydrogène) et G avec C (3 liaisons hydrogène). Il y a donc autant de A que de T, et autant de G que de C, d'où $\%A = \%T$ et $\%G = \%C$.
\item Rapport :
\[ \frac{A + T}{G + C} = \frac{22 + 22}{28 + 28} = \frac{44}{56} \approx 0{,}79. \]
Le rapport est inférieur à 1 : l'ADN est plus riche en paires G--C qu'en paires A--T. Or chaque paire G--C comporte \textbf{3} liaisons hydrogène contre 2 pour A--T : cet ADN est donc relativement \cle{riche en liaisons hydrogène}, ce qui lui confère une bonne stabilité thermique de la double hélice.
\end{enumerate}
\end{corrige}

\begin{corrige}[Exercice 5 — D'un brin d'ADN à la protéine]
\begin{enumerate}
\item Le brin transcrit (matrice) est complémentaire et antiparallèle du brin codant :
\begin{center}
Brin codant (5' $\rightarrow$ 3') : 5' -- ATG GCT TAC GGA TTT CAA TGA -- 3'\\
Brin matrice (3' $\rightarrow$ 5') : 3' -- TAC CGA ATG CCT AAA GTT ACT -- 5'
\end{center}
\item L'ARNm est complémentaire du brin matrice, avec U à la place de T ; il est synthétisé dans le sens 5' $\rightarrow$ 3' et a donc la même séquence que le brin codant (sauf U/T) :
\begin{center}
ARNm : 5' -- AUG GCU UAC GGA UUU CAA UGA -- 3'
\end{center}
\item Traduction codon par codon :
\begin{center}
\begin{tabular}{|c|c|c|c|c|c|c|c|}
\hline
\rowcolor{popblueL}
\textbf{Codon} & AUG & GCU & UAC & GGA & UUU & CAA & UGA \\
\hline
\textbf{Acide aminé} & Met & Ala & Tyr & Gly & Phe & Gln & stop \\
\hline
\end{tabular}
\end{center}
Séquence peptidique : Met -- Ala -- Tyr -- Gly -- Phe -- Gln (6 acides aminés) ; la traduction s'arrête au codon UGA.
\item L'ARNm comporte $21 / 3 = 7$ codons ; 6 codent un acide aminé, le 7\up{e} (UGA) est un codon stop qui ne code aucun acide aminé mais provoque l'arrêt de la traduction.
\end{enumerate}
\begin{attention}[Erreur fréquente]
Ne jamais écrire de T dans l'ARNm ni de U dans l'ADN ; et toujours préciser la polarité 5'/3' des séquences.
\end{attention}
\end{corrige}

\begin{corrige}[Exercice 6 — Rôles des trois types d'ARN]
\begin{enumerate}
\item \textbf{Les trois ARN et leurs rôles :}
\begin{itemize}
\item \textbf{ARNm} (messager) : copie de la séquence du gène de l'insuline ; il transporte l'information codante du noyau vers les ribosomes et sert de matrice à la traduction.
\item \textbf{ARNt} (transfert) : chaque ARNt fixe un acide aminé spécifique et l'apporte au ribosome ; son anticodon reconnaît le codon correspondant de l'ARNm.
\item \textbf{ARNr} (ribosomique) : avec des protéines, il constitue les ribosomes, « usines » d'assemblage des acides aminés ; il catalyse la formation des liaisons peptidiques.
\end{itemize}
\item \textbf{Lieux de synthèse :} l'ARNm est transcrit dans le \cle{noyau} (sur l'ADN du gène) ; les ARNt sont également transcrits dans le noyau ; les ARNr sont transcrits dans le \cle{nucléole}, où s'assemble aussi le ribosome. La traduction, elle, a lieu dans le cytoplasme (ribosomes du réticulum endoplasmique rugueux pour l'insuline, protéine sécrétée).
\item \textbf{Organigramme du trajet de l'information :}
\begin{popfigure}
\begin{tikzpicture}[node distance=0.55cm, every node/.style={font=\small},
box/.style={draw=popblue, fill=popblueL, rounded corners, align=center, inner sep=4pt}]
\node[box, align=center] (gene){Gène de l'insuline (ADN)\\ \emph{noyau}};
\node[box, below=of gene, align=center] (arnm){ARNm\\ \emph{transcription, noyau}};
\node[box, below=of arnm, align=center] (poly){Chaîne polypeptidique\\ \emph{traduction, ribosomes du RE rugueux}};
\node[box, below=of poly, align=center] (prot){Insuline fonctionnelle\\ \emph{maturation, RE + Golgi, sécrétion}};
\draw[-{Stealth}, thick, poporange] (gene) -- (arnm);
\draw[-{Stealth}, thick, poporange] (arnm) -- (poly);
\draw[-{Stealth}, thick, poporange] (poly) -- (prot);
\end{tikzpicture}
\legende{Trajet de l'information génétique : ADN $\rightarrow$ ARNm $\rightarrow$ polypeptide $\rightarrow$ insuline fonctionnelle, avec la localisation cellulaire de chaque étape.}
\end{popfigure}
\end{enumerate}
\end{corrige}

\begin{corrige}[Exercice 7 — Quantité d'ADN au cours du cycle cellulaire]
\begin{enumerate}
\item En phase G2, la réplication (phase S) a eu lieu : la quantité d'ADN a doublé :
\[ 2Q = 2 \times 6{,}6 = \SI{13,2}{pg}. \]
\item La mitose répartit équitablement les $2Q$ entre les deux cellules filles : chacune reçoit $Q = \SI{6,6}{pg}$, identique à la cellule mère en G1.
\item La méiose comporte une seule réplication ($Q \rightarrow 2Q$) suivie de deux divisions : $2Q \rightarrow Q \rightarrow Q/2$. Chaque gamète contient donc :
\[ \frac{Q}{2} = \frac{6{,}6}{2} = \SI{3,3}{pg}. \]
\item Tableau récapitulatif pour $2n = 8$ :
\begin{center}
\begin{tabularx}{\linewidth}{|l|X|X|X|X|}
\hline
\rowcolor{popblueL}
\textbf{Étape} & G1 (avant réplication) & Après interphase / début méiose I & Après méiose I & Après méiose II \\
\hline
Quantité d'ADN & $Q$ & $2Q$ & $Q$ & $Q/2$ \\
\hline
Chromosomes & $2n = 8$ & $2n = 8$ (à 2 chromatides) & $n = 4$ (à 2 chromatides) & $n = 4$ (à 1 chromatide) \\
\hline
\end{tabularx}
\end{center}
\end{enumerate}
\end{corrige}

\begin{corrige}[Exercice 8 — Analyse d'une molécule d'ADN]
\begin{enumerate}
\item Une paire de bases = 2 nucléotides. Pour 1\,000 paires de bases : $1\,000 \times 2 = 2\,000$ nucléotides.
\item Complémentarité des bases :
\begin{itemize}
\item $G = C$, donc C $= 300$ molécules ;
\item total des bases $= 2\,000$ ; $A + T = 2\,000 - (300 + 300) = 1\,400$ ;
\item $A = T$, donc A $= 700$ et T $= 700$ molécules.
\end{itemize}
Vérification : $300 + 300 + 700 + 700 = 2\,000$. \checkmark
\item Longueur du fragment :
\[ L = 1\,000 \times 0{,}34 = \SI{340}{nm} = \SI{0,34}{\micro m} \quad (1~\mu\text{m} = 1\,000~\text{nm}). \]
\item Chaque cycle de réplication double le nombre de molécules : après 3 cycles, $2^3 = 8$ molécules d'ADN. La réplication étant semi-conservative, les deux brins de la molécule initiale sont conservés chacun dans une molécule fille : seules \textbf{2 molécules} sur 8 contiennent encore un brin d'origine.
\end{enumerate}
\end{corrige}

\begin{corrige}[Exercice 9 — Type Bac : l'expérience de Meselson et Stahl (1958)]
\textbf{Barème indicatif (sur 20) :} question 1 : 3 pts ; question 2 : 6 pts (2 pts par hypothèse) ; question 3 : 4 pts ; question 4 : 4 pts ; question 5 : 3 pts.

\begin{enumerate}
\item \textbf{Principe.} La centrifugation en gradient de chlorure de césium soumet l'ADN à une force centrifuge intense : chaque molécule migre jusqu'à la zone du gradient dont la densité est égale à la sienne, où elle forme une bande. Or l'azote $^{15}$N est plus lourd que $^{14}$N : l'ADN $^{15}$N/$^{15}$N (deux brins lourds) est le plus dense (bande basse), l'ADN hybride $^{15}$N/$^{14}$N a une densité intermédiaire, et l'ADN $^{14}$N/$^{14}$N est le moins dense (bande haute). La position des bandes distingue donc les trois types de molécules.
\item \textbf{Prédictions.}
\begin{itemize}
\item \emph{Hypothèse conservative :} G1 : une bande lourde (molécule mère intacte $^{15}$N/$^{15}$N) + une bande légère (molécule fille $^{14}$N/$^{14}$N), en proportions égales. G2 : 1/4 d'ADN lourd et 3/4 d'ADN léger, jamais de bande intermédiaire.
\item \emph{Hypothèse semi-conservative :} G1 : 100 \% d'ADN hybride (chaque molécule = 1 brin $^{15}$N + 1 brin $^{14}$N), donc une seule bande intermédiaire. G2 : 50 \% hybride + 50 \% léger (les molécules hybrides se séparent : le brin $^{15}$N s'associe à un brin neuf $^{14}$N, le brin $^{14}$N ancien à un brin neuf $^{14}$N).
\item \emph{Hypothèse dispersive :} G1 : une seule bande intermédiaire (chaque brin contient moitié $^{15}$N, moitié $^{14}$N en fragments). G2 : une seule bande, plus proche de la position légère (densité intermédiaire entre hybride et léger), mais jamais deux bandes distinctes.
\end{itemize}
\item \textbf{Réfutations.}
\begin{itemize}
\item À la génération 1, on observe une \emph{unique} bande intermédiaire et \emph{aucune} bande lourde : l'hypothèse conservative, qui prédit la persistance d'une bande lourde (molécule mère intacte), est \textbf{réfutée}.
\item À la génération 2, on observe \emph{deux} bandes distinctes (50 \% hybride, 50 \% légère) : l'hypothèse dispersive, qui prédit une bande unique de densité progressivement décroissante, est \textbf{réfutée}.
\end{itemize}
\item \textbf{Conclusion.} Seule l'hypothèse semi-conservative prédit exactement les résultats observés : 100 \% hybride en G1, puis 50 \% hybride + 50 \% léger en G2. Chaque molécule fille conserve donc un brin parental et un brin néosynthétisé.\\
\emph{Prédiction pour G3 :} les molécules hybrides donnent à nouveau 50 \% hybride + 50 \% léger, et les molécules légères donnent 100 \% léger : on attend \textbf{25 \% d'ADN hybride et 75 \% d'ADN léger} (la bande intermédiaire devient deux fois moins intense que la bande haute).
\item \textbf{Déroulement au niveau d'un œil de réplication.} L'\cle{hélicase} ouvre la double hélice en rompant les liaisons hydrogène, créant deux fourches de réplication qui progressent en sens opposés. Des protéines SSB stabilisent les brins séparés ; l'ARN primase dépose une amorce. L'\cle{ADN polymérase} synthèse les brins néoformés toujours dans le sens 5' $\rightarrow$ 3', en fixant les désoxyribonucléotides complémentaires (A face à T, G face à C) : la synthèse est continue sur le brin 3' $\rightarrow$ 5' (brin précoce) et discontinue, par fragments d'Okazaki, sur l'autre brin (brin tardif) ; la ligase soude ensuite les fragments. La réplication est bidirectionnelle à partir de chaque origine.
\end{enumerate}
\begin{attention}[Points de vigilance]
\begin{itemize}
\item Justifier chaque réfutation par un \emph{résultat expérimental précis} (présence/absence d'une bande), pas par un argument théorique.
\item Ne pas écrire que l'ADN hybride « disparaît » en G2 : il représente encore 50 \% de l'ADN total.
\item Citer le sens de synthèse 5' $\rightarrow$ 3' de l'ADN polymérase : c'est une condition d'application souvent exigée.
\end{itemize}
\end{attention}
\end{corrige}

\begin{corrige}[Exercice 10 — Type Bac : exploitation de séquences et effet des mutations]
\textbf{Barème indicatif (sur 20) :} q.1 : 2 pts ; q.2 : 3 pts ; q.3 : 4 pts ; q.4 : 4 pts ; q.5 : 2 pts ; q.6 : 5 pts.

\begin{enumerate}
\item L'ARNm est complémentaire et antiparallèle du brin matrice (U remplace T) :
\begin{center}
Brin matrice : 3' -- TAC CGA ATG CCT AAA GTT ACT -- 5'\\
ARNm : 5' -- AUG GCU UAC GGA UUU CAA UGA -- 3'
\end{center}
\item Traduction : AUG = Met ; GCU = Ala ; UAC = Tyr ; GGA = Gly ; UUU = Phe ; CAA = Gln ; UGA = stop.\\
Séquence peptidique : Met -- Ala -- Tyr -- Gly -- Phe -- Gln. \textbf{Remarque :} la traduction s'achève sur un codon stop (UGA), qui ne code aucun acide aminé : le polypeptide compte donc 6 acides aminés pour 7 codons.
\item Le 6\up{e} nucléotide du brin matrice est le A du triplet CGA (positions 4--6 : C-G-A). Le brin matrice devient : 3' -- TAC CG\textbf{G} ATG CCT AAA GTT ACT -- 5'.\\
Nouvel ARNm : 5' -- AUG G\textbf{C}C UAC GGA UUU CAA UGA -- 3'.\\
Le codon GCU devient GCC : or GCU et GCC codent \emph{tous deux} l'alanine. La séquence peptidique est \textbf{inchangée} : Met -- Ala -- Tyr -- Gly -- Phe -- Gln. C'est une mutation \textbf{synonyme (silencieuse)}, rendue possible par la redondance (dégénérescence) du code génétique.
\item La délétion du 4\up{e} nucléotide (le C de CGA) donne un brin matrice : 3' -- TAC GAA TGC CTA AAG TTA CT -- 5' (20 nucléotides).\\
Nouvel ARNm : 5' -- AUG GCU UAC GGA UUU CAA UGA -- 3' devient, en regroupant par triplets depuis le décalage : 5' -- AUG CUU ACG GAU UUC AAU GA -- 3', soit les codons : AUG -- CUU -- ACG -- GAU -- UUC -- AAU (le dernier nucléotide reste seul).\\
Séquence peptidique : Met -- Leu -- Thr -- Asp -- Phe -- Asn (et absence de codon stop dans le fragment lu).\\
\textbf{Comparaison :} la substitution (q.3) ne modifie qu'un seul codon, voire aucun acide aminé ; la délétion \textbf{décale tout le cadre de lecture} à partir du site de la mutation : tous les codons situés en aval sont modifiés (mutation « frameshift »). Ses conséquences sont donc beaucoup plus amples : protéine profondément différente, généralement non fonctionnelle.
\item Le codon n\textsuperscript{o} 3 (UAC, codant Tyr) devient UAA, un \cle{codon stop}. La traduction s'arrête prématurément : on obtient un polypeptide \textbf{tronqué} (Met -- Ala, 2 acides aminés), presque toujours non fonctionnel. C'est une mutation \emph{non sens}.
\item \textbf{Propriétés du code génétique dégagées :}
\begin{itemize}
\item il est \textbf{redondant} (dégénéré) : GCU et GCC codent le même acide aminé, d'où l'existence de mutations silencieuses ;
\item il est lu par \textbf{triplets non chevauchants} dans un cadre de lecture fixe : toute insertion ou délétion décale le cadre et bouleverse toute la séquence en aval.
\end{itemize}
\textbf{Relation gène -- protéine -- caractère :} la séquence nucléotidique du gène détermine, via l'ARNm, la séquence en acides aminés de la protéine ; celle-ci conditionne la structure tridimensionnelle et donc l'activité de la protéine (enzyme, transporteur, etc.), dont dépend le caractère observable. Une mutation modifiant la séquence du gène peut donc modifier la protéine et, en cascade, le phénotype : le gène contrôle le caractère \emph{par l'intermédiaire} de la protéine.
\end{enumerate}
\begin{attention}[Points de vigilance]
\begin{itemize}
\item Toujours repérer la polarité du brin donné avant de transcrire (ici le brin matrice est lu 3' $\rightarrow$ 5', l'ARNm s'écrit 5' $\rightarrow$ 3').
\item Pour la délétion, regrouper les codons \emph{depuis le début} de la séquence, pas depuis le point de mutation.
\item Distinguer clairement mutation synonyme (même acide aminé), faux sens (acide aminé différent) et non sens (codon stop).
\end{itemize}
\end{attention}
\end{corrige}

\begin{corrige}[Exercice 11 — Type Bac : quantité d'ADN au cours de la méiose et albinisme]
\textbf{Barème indicatif (sur 20) :} Partie A : q.1 : 4 pts ; q.2 : 4 pts ; q.3 : 2 pts ; q.4 : 4 pts. Partie B : q.5 : 4 pts ; q.6 : 2 pts.

\textbf{Partie A.}
\begin{enumerate}
\item Courbe de variation de la quantité d'ADN :
\begin{popfigure}
\begin{tikzpicture}
\begin{axis}[width=0.95\linewidth, height=6.2cm,
xlabel={Temps (h)}, ylabel={Quantité d'ADN},
xmin=0, xmax=25, ymin=0, ymax=2.3,
xtick={0,4,8,12,16,20,24},
ytick={0.5,1,1.5,2},
yticklabels={$Q/2$,$Q$,$1{,}5Q$,$2Q$},
grid=both, grid style={popline!40},
axis line style={popdark}]
\addplot[popcurve, mark=*, mark size=1.5pt] coordinates {(0,1) (4,1.5) (8,2) (12,2) (16,1) (20,1) (24,0.5)};
\node[font=\small, popdark] at (axis cs:4,2.15) {réplication};
\node[font=\small, popdark] at (axis cs:14,2.15) {méiose I};
\node[font=\small, popdark] at (axis cs:22,1.15) {méiose II};
\end{axis}
\end{tikzpicture}
\legende{Variation de la quantité d'ADN au cours de la méiose : montée progressive de $Q$ à $2Q$ (réplication), palier à $2Q$, chute à $Q$ (première division), palier à $Q$, puis chute à $Q/2$ (seconde division).}
\end{popfigure}
\item \textbf{Identification des phases :}
\begin{itemize}
\item \emph{Réplication (phase S de l'interphase)} : entre 0 et 8 h, la quantité d'ADN augmente progressivement de $Q$ à $2Q$ : c'est la duplication semi-conservative de l'ADN.
\item \emph{Première division (méiose I, division de réduction)} : entre 12 et 16 h, la quantité passe de $2Q$ à $Q$ : les deux cellules filles reçoivent chacune la moitié de l'ADN (chromosomes à deux chromatides, $n$ chromosomes).
\item \emph{Seconde division (méiose II, division équationnelle)} : entre 20 et 24 h, la quantité passe de $Q$ à $Q/2$, sans réplication préalable : séparation des chromatides sœurs.
\end{itemize}
\item Chacune des quatre cellules obtenues contient une quantité d'ADN égale à $Q/2$ et $n = 4$ chromosomes (à une chromatide).
\item \textbf{Comparaison mitose / méiose :}
\begin{center}
\begin{tabularx}{\linewidth}{|l|X|X|X|X|}
\hline
\rowcolor{popblueL}
 & \multicolumn{2}{c|}{\textbf{Mitose}} & \multicolumn{2}{c|}{\textbf{Méiose}} \\
\cline{2-5}
 & Avant division & Après division & Avant divisions & Après divisions \\
\hline
Quantité d'ADN & $2Q$ & $Q$ (par cellule) & $2Q$ & $Q/2$ (par cellule) \\
\hline
Chromosomes & $2n$ & $2n$ & $2n$ & $n$ \\
\hline
Cellules produites & \multicolumn{2}{c|}{2 cellules identiques} & \multicolumn{2}{c|}{4 cellules génétiquement différentes} \\
\hline
\end{tabularx}
\end{center}
\textbf{Particularité fondamentale de la méiose :} une seule réplication suivie de \emph{deux} divisions successives ; elle divise par deux le nombre de chromosomes ($2n \rightarrow n$) et par quatre la quantité d'ADN ($2Q \rightarrow Q/2$), assurant la production de gamètes haploïdes et le maintien de la diploïdie à la fécondation.
\end{enumerate}

\textbf{Partie B.}
\begin{enumerate}
\setcounter{enumi}{4}
\item \textbf{Chaîne de causalité :} la mutation modifie la séquence nucléotidique du gène de la tyrosinase $\rightarrow$ la séquence de l'ARNm transcrit est modifiée $\rightarrow$ la traduction produit une tyrosinase dont la séquence en acides aminés est altérée (ou tronquée) $\rightarrow$ le repliement tridimensionnel de l'enzyme est perturbé, son site actif n'est plus fonctionnel $\rightarrow$ la tyrosine n'est plus transformée en mélanine $\rightarrow$ absence de pigment dans la peau, les poils et l'iris : phénotype albinos. La mutation du gène agit donc sur le caractère \emph{par l'intermédiaire de la protéine} : gène $\rightarrow$ protéine $\rightarrow$ caractère.
\item Le code génétique étant lu par triplets, plusieurs types de mutations peuvent aboutir au même résultat : une mutation non sens créant un codon stop prématuré (enzyme tronquée), une délétion décalant le cadre de lecture, ou une mutation faux sens touchant un acide aminé essentiel du site actif. Dans tous les cas, la structure spatiale de la tyrosinase est compromise : des mutations différentes du même gène peuvent donc toutes produire une enzyme non fonctionnelle et le même phénotype albinos.
\end{enumerate}
\begin{attention}[Points de vigilance]
\begin{itemize}
\item La courbe doit montrer une \emph{montée progressive} (et non brutale) de $Q$ à $2Q$ : la réplication prend du temps (phase S).
\item Après la méiose I, les chromosomes sont encore à \emph{deux chromatides} : quantité $Q$ pour $n$ chromosomes — ne pas écrire $Q/2$ à cette étape.
\item Dans la partie B, expliciter chaque maillon de la chaîne de causalité : une réponse du type « la mutation empêche la fabrication de mélanine » sans passer par l'ARNm et la protéine est insuffisante au Bac.
\end{itemize}
\end{attention}
\end{corrige}

\newpage

\coursec{Fiche bilan}

\begin{popfigure}
\begin{center}\tcbox[colback=popblue,colframe=popblue,coltext=white,arc=9pt,boxsep=4pt,left=6pt,right=6pt]{\bfseries\small Du gène à la protéine}\end{center}
\vspace{-2pt}
\begin{tcbraster}[raster columns=3,raster equal height=rows,raster column skip=6pt,raster row skip=6pt,enhanced,blankest]
\begin{tcolorbox}[enhanced,colback=white,colframe=poporange,boxrule=0.9pt,arc=3pt,title={Acides nucléiques},fonttitle=\bfseries\scriptsize,coltitle=white,colbacktitle=poporange,left=4pt,right=4pt,top=2pt,bottom=2pt,boxsep=1pt,toptitle=1.5pt,bottomtitle=1.5pt,halign=left]
\begin{itemize}[nosep,leftmargin=*,itemsep=1pt,font=\scriptsize]
\item ADN : double hélice A--T, G--C
\item ARN : simple brin, U
\end{itemize}
\end{tcolorbox}
\begin{tcolorbox}[enhanced,colback=white,colframe=popgreen,boxrule=0.9pt,arc=3pt,title={Réplication},fonttitle=\bfseries\scriptsize,coltitle=white,colbacktitle=popgreen,left=4pt,right=4pt,top=2pt,bottom=2pt,boxsep=1pt,toptitle=1.5pt,bottomtitle=1.5pt,halign=left]
\begin{itemize}[nosep,leftmargin=*,itemsep=1pt,font=\scriptsize]
\item Semi-conservative
\item ADN polymérase, phase S
\end{itemize}
\end{tcolorbox}
\begin{tcolorbox}[enhanced,colback=white,colframe=poppurple,boxrule=0.9pt,arc=3pt,title={Transcription},fonttitle=\bfseries\scriptsize,coltitle=white,colbacktitle=poppurple,left=4pt,right=4pt,top=2pt,bottom=2pt,boxsep=1pt,toptitle=1.5pt,bottomtitle=1.5pt,halign=left]
\begin{itemize}[nosep,leftmargin=*,itemsep=1pt,font=\scriptsize]
\item Noyau
\item ARN polymérase $\to$ ARNm
\end{itemize}
\end{tcolorbox}
\begin{tcolorbox}[enhanced,colback=white,colframe=poppink,boxrule=0.9pt,arc=3pt,title={Traduction},fonttitle=\bfseries\scriptsize,coltitle=white,colbacktitle=poppink,left=4pt,right=4pt,top=2pt,bottom=2pt,boxsep=1pt,toptitle=1.5pt,bottomtitle=1.5pt,halign=left]
\begin{itemize}[nosep,leftmargin=*,itemsep=1pt,font=\scriptsize]
\item Cytoplasme, ribosome
\item codons $\to$ acides aminés
\end{itemize}
\end{tcolorbox}
\begin{tcolorbox}[enhanced,colback=white,colframe=popgold,boxrule=0.9pt,arc=3pt,title={Code génétique},fonttitle=\bfseries\scriptsize,coltitle=white,colbacktitle=popgold,left=4pt,right=4pt,top=2pt,bottom=2pt,boxsep=1pt,toptitle=1.5pt,bottomtitle=1.5pt,halign=left]
\begin{itemize}[nosep,leftmargin=*,itemsep=1pt,font=\scriptsize]
\item 64 codons, universel
\item gène $\to$ protéine $\to$ caractère
\end{itemize}
\end{tcolorbox}
\begin{tcolorbox}[enhanced,colback=white,colframe=popblue!70,boxrule=0.9pt,arc=3pt,title={ADN et méiose},fonttitle=\bfseries\scriptsize,coltitle=white,colbacktitle=popblue!70,left=4pt,right=4pt,top=2pt,bottom=2pt,boxsep=1pt,toptitle=1.5pt,bottomtitle=1.5pt,halign=left]
\begin{itemize}[nosep,leftmargin=*,itemsep=1pt,font=\scriptsize]
\item $Q \to 2Q \to Q \to Q/2$
\end{itemize}
\end{tcolorbox}
\end{tcbraster}

\legende{Carte mentale de la leçon : de la structure de l'ADN à l'expression du caractère.}
\end{popfigure}

\begin{bilan}
\textbf{Définitions clés}
\begin{itemize}
  \item \cle{Nucléotide} : unité de base d'un acide nucléique = acide phosphorique + sucre (désoxyribose ou ribose) + base azotée (A, T, G, C ou U).
  \item \cle{ADN} : double hélice antiparallèle ; liaisons hydrogène A=T (2 liaisons) et G$\equiv$C (3 liaisons) ; support de l'information génétique.
  \item \cle{Réplication} : duplication de l'ADN en phase S de l'interphase ; dite \cle{semi-conservative} car chaque molécule fille garde un brin parental.
  \item \cle{Transcription} : synthèse d'un ARNm complémentaire d'un brin d'ADN, dans le noyau, par l'ARN polymérase.
  \item \cle{Traduction} : synthèse d'une protéine au niveau des ribosomes (cytoplasme) à partir de l'ARNm, avec les ARNt et leurs anticodons.
  \item \cle{Code génétique} : correspondance entre triplets de nucléotides (codons) et acides aminés.
  \item \cle{Gène} : segment d'ADN codant une protéine (ou une chaîne polypeptidique).
\end{itemize}

\textbf{À connaître par cœur}
\begin{center}
\begin{tabularx}{\linewidth}{|l|X|X|}\hline
\rowcolor{popblueL} \textbf{Étape} & \textbf{Localisation} & \textbf{Acteurs / produit} \\ \hline
Réplication & Noyau (phase S) & ADN polymérase ; 2 ADN identiques \\ \hline
Transcription & Noyau & ARN polymérase ; ARNm \\ \hline
Traduction & Cytoplasme (ribosomes) & ARNm, ARNt, acides aminés ; protéine \\ \hline
\end{tabularx}
\end{center}

\begin{itemize}
  \item Complémentarité des bases : A--T (ADN) ; A--U (ARN) ; G--C. Sens de lecture de l'ARNm : $5' \to 3'$.
  \item Code génétique : $4^3 = 64$ codons ; codon d'initiation AUG (méthionine) ; codons stop UAA, UAG, UGA ; code dégénéré et universel.
  \item Relation fondamentale : gène (ADN) $\xrightarrow{\text{transcription}}$ ARNm $\xrightarrow{\text{traduction}}$ protéine $\to$ caractère.
  \item Quantité d'ADN au cours de la méiose : $Q \xrightarrow{\text{réplication}} 2Q \xrightarrow{\text{division I}} Q \xrightarrow{\text{division II}} Q/2$.
\end{itemize}

\textbf{Méthodes express}
\begin{itemize}
  \item Écrire le brin complémentaire : remplacer chaque base (A$\leftrightarrow$T ou A$\leftrightarrow$U, G$\leftrightarrow$C) et inverser le sens.
  \item Trouver la protéine : découper l'ARNm en codons à partir d'AUG, lire le tableau du code génétique, s'arrêter au codon stop.
  \item Courbe d'ADN : repérer la phase S (montée $Q \to 2Q$) puis les deux divisions (chutes successives).
\end{itemize}
\end{bilan}

\begin{aretenir}[Checklist avant le Bac]
\begin{itemize}
  \item[$\square$] Je sais nommer les 3 constituants d'un nucléotide et citer les 5 bases azotées.
  \item[$\square$] Je sais que l'ADN est une double hélice antiparallèle avec A=T et G$\equiv$C.
  \item[$\square$] Je sais différencier ADN et ARN (sucre, bases, nombre de brins).
  \item[$\square$] Je sais expliquer pourquoi la réplication est dite semi-conservative (expérience de Meselson et Stahl).
  \item[$\square$] Je sais localiser la transcription (noyau) et la traduction (cytoplasme).
  \item[$\square$] Je sais écrire l'ARNm complémentaire d'un brin d'ADN donné.
  \item[$\square$] Je sais utiliser le tableau du code génétique pour traduire un ARNm en protéine.
  \item[$\square$] Je connais le codon d'initiation (AUG) et les 3 codons stop.
  \item[$\square$] Je sais énoncer la relation gène $\to$ protéine $\to$ caractère avec un exemple (hémoglobine, albinisme).
  \item[$\square$] Je sais tracer et légender la courbe $Q \to 2Q \to Q \to Q/2$ de la quantité d'ADN au cours de la méiose.
  \item[$\square$] Je soigne mes schémas : double hélice, ribosome sur l'ARNm, légendes précises.
  \item[$\square$] Je vérifie toujours le sens $5' \to 3'$ avant de traduire une séquence.
\end{itemize}
\end{aretenir}

\findecours

\end{document}
